% Les problèmes d'émigration — Allemand Tle A — Cours signé OnBuch+
% Programme officiel MINESEC. Compiler avec : tectonic AL28-les-problemes-d-emigration.tex
\newcommand{\DOCMATIERE}{Allemand}
\newcommand{\DOCNIVEAU}{Tle A}
\newcommand{\DOCMODULE}{Chapitre 5 : Citoyenneté}
\newcommand{\DOCLECON}{AL28}
\newcommand{\DOCTITRE}{Les problèmes d'émigration}
% OnBuch+ — préambule « cours signés » ENRICHI (compilé avec tectonic / XeLaTeX)
% Système visuel « Pop » d'OnBuch+ : fond crème (jamais blanc), bordures encre
% épaisses, ombres « dures » décalées, titres Archivo Black en capitales.
% Version MATHÉMATIQUES : graphes (pgfplots/tikz), boîtes pédagogiques
% (définition, propriété, méthode, attention, le savais-tu, exercice résolu,
% exercice, TP, bilan), page de garde et signature OnBuch+ ancrée en bas.
%
% Macros attendues dans main.tex AVANT \input{preamble} :
%   \DOCMATIERE  \DOCNIVEAU  \DOCMODULE  \DOCLECON (numéro)  \DOCTITRE
\documentclass[11pt]{article}

\usepackage{fontspec}
\usepackage[ngerman,french]{babel} % français = langue principale ; l'allemand via \all{..} / textebox
\usepackage[a4paper,margin=2cm,top=3.4cm,bottom=2.8cm,headheight=1.4cm,footskip=1.3cm]{geometry}
\usepackage{amsmath,amssymb,amsfonts}
\usepackage{mathtools} % \xmapsto, \coloneqq... souvent utilisés par les modèles
\usepackage{cancel} % simplifications barrées (bilans de forces)
% \abs{z} (module, valeur absolue) et \norm{u} (norme), utilisés par les modèles
\providecommand{\abs}{}\let\abs\relax\DeclarePairedDelimiter{\abs}{\lvert}{\rvert}
\providecommand{\norm}{}\let\norm\relax\DeclarePairedDelimiter{\norm}{\lVert}{\rVert}
\usepackage[table]{xcolor} % [table] : fournit aussi \rowcolors (alternance de lignes)
\usepackage{pagecolor}
\usepackage{tikz}
\usetikzlibrary{arrows.meta,positioning,calc,shapes.geometric,shapes.misc,shapes.arrows,decorations.pathmorphing,decorations.pathreplacing,decorations.markings,patterns,shadows,mindmap,trees,fit,backgrounds,matrix,intersections,angles,quotes}
\usepackage{pgfplots}
\usepackage[version=4]{mhchem} % équations nucléaires \ce{^{235}_{92}U}
\usepackage[european,siunitx]{circuitikz} % schémas électriques
\pgfplotsset{compat=1.18}
\usepgfplotslibrary{fillbetween,groupplots} % aires entre courbes (calcul intégral)
\usepackage{enumitem}
\usepackage{fancyhdr}
% [most] charge toutes les bibliothèques tcolorbox (listings, minted, raster,
% documentation...) et gonfle inutilement la mémoire TeX sur un document long
% (~300 boîtes) : on ne charge que ce qui est réellement utilisé.
\usepackage[skins,breakable]{tcolorbox}
\usepackage{graphicx}
\usepackage{colortbl}
\usepackage{booktabs}
\usepackage{tabularx}
\usepackage{diagbox} % case d'en-tête diagonale des tableaux à double entrée
% Colonnes extensibles centrée / gauche / droite (C, L, R), souvent utilisées par les modèles
\newcolumntype{C}{>{\centering\arraybackslash}X}
\newcolumntype{L}{>{\raggedright\arraybackslash}X}
\newcolumntype{R}{>{\raggedleft\arraybackslash}X}
\usepackage{array}
\usepackage{multicol}
\usepackage{multirow}
\usepackage{siunitx}
% parse-numbers=false : accepte aussi \SI{2,5 \times 10^{-3}}{...}
\sisetup{locale=FR,output-decimal-marker={,},group-minimum-digits=4,parse-numbers=false}
\DeclareSIUnit\ohm{\ensuremath{\symup{\Omega}}} % U+2126 absent de la police
\DeclareSIUnit\division{div} % réglages d'oscilloscope (ms/div, V/div)
\DeclareSIUnit\tr{tr} % tours (vitesse de rotation, ex. \SI{1500}{\tr\per\minute})
\DeclareSIUnit\molar{M} % concentration molaire (ex. \SI{3}{\molar}, \SI{1}{\micro\molar})
\DeclareSIUnit\an{an} % année (le modèle confond parfois avec \year, un compteur TeX) — ex. \SI{5}{\kilo\meter\per\an}
\DeclareSIUnit\FCFA{FCFA} % franc CFA (ex. \SI{500}{\FCFA\per\kilo\gram})
\DeclareSIUnit\degreeBrix{\textdegree Brix} % teneur en sucre d'un sirop/jus (réfractométrie)
\DeclareSIUnit\month{mois} % le modèle utilise parfois \month, un compteur TeX, comme unité de durée
\usepackage{qrcode}
% \degree : commande fréquemment utilisée par les modèles pour un angle ou une
% température en dehors de \SI{}{} (non fournie par les paquets chargés).
\providecommand{\degree}{^{\circ}}
% \qed (fin de démonstration), utilisé par les modèles sans amsthm
\providecommand{\qed}{\hfill$\square$}
% \barre : double barre des valeurs interdites dans un tableau de variations (tkz-tab)
\providecommand{\barre}{\ensuremath{\|}}
% environnement proof (amsthm non chargé) utilisé par les modèles
\ifdefined\proof\else\newenvironment{proof}[1][Démonstration]{\par\noindent\textit{#1.}\ }{\qed\par}\fi
\usepackage{lastpage}
\usepackage{hyperref}
\hypersetup{hidelinks}

% ============================================================
% Palette « Pop » OnBuch+ — mêmes hex que la classe P (plus_ui.dart)
% ============================================================
\definecolor{poporange}{HTML}{F59321}
\definecolor{popdark}{HTML}{E07A0C}
\definecolor{popcream}{HTML}{FFF6E8}
\definecolor{popink}{HTML}{1C1714}
\definecolor{poppurple}{HTML}{7A5AE0}
\definecolor{popgreen}{HTML}{2E9E6A}
\definecolor{poppink}{HTML}{E0457B}
\definecolor{popblue}{HTML}{2F7DE1}
\definecolor{popgold}{HTML}{C98A12}
\definecolor{popmuted}{HTML}{8A7F76}
\definecolor{popline}{HTML}{EADFD2}
\definecolor{popgray}{HTML}{8A7F76}
\colorlet{popgrey}{popgray}
% Alias tolérants (le modèle nomme parfois les couleurs différemment)
\colorlet{popwhite}{white}
\colorlet{popblack}{popink}
\colorlet{popred}{poppink}
\colorlet{popyellow}{popgold}
% Teintes claires (fonds de tableaux/encadrés) — le modèle nomme parfois
% les couleurs avec un suffixe L (ex. popblueL) sans les avoir définies.
\colorlet{poporangeL}{poporange!20}
\colorlet{popdarkL}{popdark!20}
\colorlet{poppurpleL}{poppurple!20}
\colorlet{popgreenL}{popgreen!20}
\colorlet{poppinkL}{poppink!20}
\colorlet{popblueL}{popblue!20}
\colorlet{popgoldL}{popgold!20}
\colorlet{popredL}{popred!20}
\colorlet{popgrayL}{popgray!20}
% Teintes claires pour fonds de tableaux / zones
\colorlet{popblueL}{popblue!10!white}
\colorlet{poporangeL}{poporange!14!white}
\colorlet{popgreenL}{popgreen!12!white}
\colorlet{poppurpleL}{poppurple!12!white}
\colorlet{poppinkL}{poppink!12!white}
\colorlet{popgoldL}{popgold!14!white}

\pagecolor{popcream}

% ============================================================
% Typographie
% ============================================================
\defaultfontfeatures{Ligatures=TeX}
\setmainfont{PlusJakartaSans-Medium.ttf}[Path=../../../fonts/,BoldFont=PlusJakartaSans-Bold.ttf]
% Maths en sans-serif (Fira Math) avec chiffres et lettres droites de Plus
% Jakarta, pour que les formules se fondent dans le texte.
\usepackage[math-style=ISO,bold-style=ISO]{unicode-math}
\setmathfont{FiraMath-Regular.otf}[Path=../../../fonts/]
\setmathfont{PlusJakartaSans-Medium.ttf}[Path=../../../fonts/,range={up/num,up/{Latin,latin},\mathup}]
\usepackage{icomma} % virgule décimale sans espace en mode math
% Caractères mathématiques Unicode absents des polices de texte (Plus Jakarta,
% Archivo Black) : rendus par leur équivalent mathématique, en texte comme en
% formule — sinon ils s'affichent en carré vide (ex. « Arithmétique dans ℤ »).
\usepackage{newunicodechar}
\newunicodechar{ℝ}{\ensuremath{\mathbb{R}}}
\newunicodechar{ℤ}{\ensuremath{\mathbb{Z}}}
\newunicodechar{ℂ}{\ensuremath{\mathbb{C}}}
\newunicodechar{ℕ}{\ensuremath{\mathbb{N}}}
\newunicodechar{ℚ}{\ensuremath{\mathbb{Q}}}
\newunicodechar{𝔻}{\ensuremath{\mathbb{D}}}
\newunicodechar{α}{\ensuremath{\alpha}}
\newunicodechar{β}{\ensuremath{\beta}}
\newunicodechar{γ}{\ensuremath{\gamma}}
\newunicodechar{δ}{\ensuremath{\delta}}
\newunicodechar{ε}{\ensuremath{\varepsilon}}
\newunicodechar{θ}{\ensuremath{\theta}}
\newunicodechar{λ}{\ensuremath{\lambda}}
\newunicodechar{μ}{\ensuremath{\mu}}
\newunicodechar{σ}{\ensuremath{\sigma}}
\newunicodechar{τ}{\ensuremath{\tau}}
\newunicodechar{φ}{\ensuremath{\varphi}}
\newunicodechar{ψ}{\ensuremath{\psi}}
\newunicodechar{ω}{\ensuremath{\omega}}
\newunicodechar{Σ}{\ensuremath{\Sigma}}
% majuscules produites par \MakeUppercase dans les titres (α-aminés -> Α)
\newunicodechar{Α}{\ensuremath{\alpha}}
\newunicodechar{Β}{\ensuremath{\beta}}
\newunicodechar{Γ}{\ensuremath{\Gamma}}
\newunicodechar{Δ}{\ensuremath{\Delta}}
\newunicodechar{Ε}{\ensuremath{\varepsilon}}
\newunicodechar{Θ}{\ensuremath{\Theta}}
\newunicodechar{Λ}{\ensuremath{\Lambda}}
\newunicodechar{Μ}{\ensuremath{\mu}}
\newunicodechar{Τ}{\ensuremath{\tau}}
\newunicodechar{Φ}{\ensuremath{\Phi}}
\newunicodechar{Ψ}{\ensuremath{\Psi}}
\newunicodechar{Ω}{\ensuremath{\Omega}}
\newunicodechar{↦}{\ensuremath{\mapsto}}
\newunicodechar{∈}{\ensuremath{\in}}
\newunicodechar{∉}{\ensuremath{\notin}}
\newunicodechar{∧}{\ensuremath{\wedge}}
\newunicodechar{⇔}{\ensuremath{\Leftrightarrow}}
\newunicodechar{⟺}{\ensuremath{\Longleftrightarrow}}
\newunicodechar{⇒}{\ensuremath{\Rightarrow}}
\newunicodechar{⟹}{\ensuremath{\Longrightarrow}}
\newunicodechar{∘}{\ensuremath{\circ}}
\newunicodechar{∪}{\ensuremath{\cup}}
\newunicodechar{∩}{\ensuremath{\cap}}
\newunicodechar{≡}{\ensuremath{\equiv}}
\newunicodechar{⊂}{\ensuremath{\subset}}
\newunicodechar{⊥}{\ensuremath{\perp}}
\newunicodechar{∃}{\ensuremath{\exists}}
\newunicodechar{∀}{\ensuremath{\forall}}
\newunicodechar{∛}{\ensuremath{\sqrt[3]{\,}}}
\newunicodechar{∜}{\ensuremath{\sqrt[4]{\,}}}
\newunicodechar{⃗}{\ensuremath{{}^{\to}}}
\newunicodechar{ⁿ}{\ifmmode^{n}\else\textsuperscript{\ensuremath{n}}\fi}
\newunicodechar{ˣ}{\ifmmode^{x}\else\textsuperscript{\ensuremath{x}}\fi}
\newunicodechar{ᵇ}{\ifmmode^{b}\else\textsuperscript{\ensuremath{b}}\fi}
\newunicodechar{ʳ}{\ifmmode^{r}\else\textsuperscript{\ensuremath{r}}\fi}
\newunicodechar{ᵘ}{\ifmmode^{u}\else\textsuperscript{\ensuremath{u}}\fi}
\newunicodechar{ʸ}{\ifmmode^{y}\else\textsuperscript{\ensuremath{y}}\fi}
\newunicodechar{ᵃ}{\ifmmode^{a}\else\textsuperscript{\ensuremath{a}}\fi}
\newunicodechar{ᵉ}{\ifmmode^{e}\else\textsuperscript{\ensuremath{e}}\fi}
\newunicodechar{ᵏ}{\ifmmode^{k}\else\textsuperscript{\ensuremath{k}}\fi}
\newunicodechar{ᵖ}{\ifmmode^{p}\else\textsuperscript{\ensuremath{p}}\fi}
\newunicodechar{ᴺ}{\ifmmode^{N}\else\textsuperscript{\ensuremath{N}}\fi}
\newunicodechar{ᶜ}{\ifmmode^{c}\else\textsuperscript{\ensuremath{c}}\fi}
\newunicodechar{ʰ}{\ifmmode^{h}\else\textsuperscript{\ensuremath{h}}\fi}
\newunicodechar{ᵠ}{\ifmmode^{q}\else\textsuperscript{\ensuremath{q}}\fi}
\newunicodechar{ₙ}{\ifmmode_{n}\else\textsubscript{\ensuremath{n}}\fi}
\newunicodechar{ₐ}{\ifmmode_{a}\else\textsubscript{\ensuremath{a}}\fi}
\newunicodechar{ᵢ}{\ifmmode_{i}\else\textsubscript{\ensuremath{i}}\fi}
\newunicodechar{ᵣ}{\ifmmode_{r}\else\textsubscript{\ensuremath{r}}\fi}
\newunicodechar{ₖ}{\ifmmode_{k}\else\textsubscript{\ensuremath{k}}\fi}
\newunicodechar{ᵦ}{\ifmmode_{\beta}\else\textsubscript{\ensuremath{\beta}}\fi}
\newunicodechar{ₓ}{\ifmmode_{x}\else\textsubscript{\ensuremath{x}}\fi}
\newfontfamily\popdisplay{ArchivoBlack-Regular.ttf}[Path=../../../fonts/]
\newfontfamily\poplogo{SpaceGrotesk-Bold.ttf}[Path=../../../fonts/]
\newfontfamily\popsemi{PlusJakartaSans-SemiBold.ttf}[Path=../../../fonts/]
\newfontfamily\popxbold{PlusJakartaSans-ExtraBold.ttf}[Path=../../../fonts/]
\newfontfamily\popxboldls{PlusJakartaSans-ExtraBold.ttf}[Path=../../../fonts/,LetterSpace=3.0]

\setlist[enumerate]{leftmargin=1.7em,itemsep=5pt,topsep=4pt}
\setlist[itemize]{leftmargin=1.7em,itemsep=4pt,topsep=4pt,label={\color{poporange}\textbullet}}
\setlength{\parindent}{0pt}
\setlength{\parskip}{5pt}
\renewcommand{\arraystretch}{1.25}

% pgfplots : style Pop par défaut
\pgfplotsset{
  every axis/.append style={axis line style={popink,line width=1pt},tick style={popink},
    grid style={popline,line width=0.5pt},label style={font=\small},tick label style={font=\footnotesize}},
  popcurve/.style={poporange,line width=1.8pt,no markers,smooth},
}
% Même style disponible dans un \draw TikZ simple (hors pgfplots)
\tikzset{popcurve/.style={poporange,line width=1.8pt}}
\tikzset{popgrid/.style={popline,very thin}}
\colorlet{popcurve}{poporange} % le nom du style est parfois utilisé comme couleur

% ============================================================
% Pill / squiggle
% ============================================================
\newcommand{\poppill}[3]{%
  \tikz[baseline=-0.55ex]\node[draw=popink,line width=1.2pt,rounded corners=2.6pt,%
    fill=#2,inner xsep=6.5pt,inner ysep=3pt]{%
    \popxboldls\fontsize{7}{7}\selectfont\color{#3}\MakeUppercase{#1}};%
}
\newcommand{\popsquiggle}[1]{%
  \tikz[baseline=-0.4ex]{
    \draw[#1,line width=2.6pt,line cap=round]
      plot [smooth] coordinates {(0,0.09) (0.34,0.01) (0.68,0.21) (1.02,0.01) (1.36,0.10)};
  }%
}
% Mot-clé surligné (marqueur orange sous le texte)
\newcommand{\cle}[1]{{\popxbold\color{popdark}#1}}

% ============================================================
% En-tête / pied
% ============================================================
\pagestyle{fancy}
\fancyhf{}
\renewcommand{\headrulewidth}{0pt}
\renewcommand{\footrulewidth}{0pt}
\lhead{\raisebox{-0.15ex}{\poplogo\fontsize{13}{13}\selectfont\color{popink}OnBuch\color{poporange}+}}
\rhead{\poppill{\DOCMATIERE\ \textbullet\ \DOCNIVEAU\ \textbullet\ Leçon \DOCLECON}{white}{popink}}
\lfoot{\popxbold\fontsize{7.6}{7.6}\selectfont\color{popmuted}\MakeUppercase{Cours signé OnBuch+}\ \textbullet\ {\popsemi\DOCTITRE}}
\rfoot{\tikz[baseline=-0.6ex]\node[rounded corners=8.5pt,fill=popink,minimum height=17pt,minimum width=17pt,inner xsep=5pt,inner sep=0]{\popxbold\fontsize{8}{8}\selectfont\color{popcream}\thepage\,/\,\pageref*{LastPage}};}

% ============================================================
% Titres
% ============================================================
\newcommand{\chaptitre}[2]{%
  \begin{tcolorbox}[enhanced,colback=poporange,colframe=popink,boxrule=2.6pt,arc=11pt,
      drop shadow=popink,left=15pt,right=15pt,top=12pt,bottom=11pt,before skip=3pt,after skip=18pt]
    \poppill{#2}{popink}{popcream}\par\vspace{9pt}
    {\raggedright\hyphenpenalty=10000\exhyphenpenalty=10000\popdisplay\fontsize{22}{24}\selectfont\color{popcream}\MakeUppercase{#1}\par}
  \end{tcolorbox}
}
% Section numérotée : pastille orange avec le numéro + titre Archivo
\newcounter{popsec}
\newcommand{\coursec}[1]{%
  \stepcounter{popsec}\setcounter{popsub}{0}%
  \par\vspace{16pt}\noindent
  \begin{minipage}{\linewidth}
  \tikz[baseline=-0.7ex]\node[draw=popink,line width=1.6pt,fill=poporange,rounded corners=4pt,minimum size=22pt,inner sep=0]{\popdisplay\fontsize{12}{12}\selectfont\color{popcream}\thepopsec};%
  \hspace{8pt}{\popdisplay\fontsize{15}{17}\selectfont\color{popink}\MakeUppercase{#1}}\par
  \vspace{4pt}\noindent{\color{poporange}\rule{36pt}{3.6pt}}
  \end{minipage}\par\nopagebreak\vspace{8pt}
}
\newcounter{popsub}
\newcommand{\courssub}[1]{%
  \stepcounter{popsub}%
  \par\vspace{9pt}\noindent{\popxbold\fontsize{11.5}{13}\selectfont\color{popink}{\color{poporange}\thepopsec.\thepopsub}\hspace{6pt}#1}\par\nopagebreak\vspace{4pt}
}

% ============================================================
% Boîtes pédagogiques — même recette : fond blanc, bordure épaisse colorée,
% ombre dure encre, badge plein. Titre optionnel [#1].
% ============================================================
\tcbset{popbase/.style={breakable,enhanced,colback=white,boxrule=2.3pt,arc=9pt,
  shadow={3pt}{-3pt}{0pt}{popink},left=14pt,right=14pt,top=11pt,bottom=11pt,before skip=11pt,after skip=15pt,
  fontupper=\popsemi\fontsize{10.6}{14.5}\selectfont\color{popink},
  fontlower=\popsemi\fontsize{10.6}{14.5}\selectfont\color{popink}}}
% Coche dessinée (le glyphe ✓ n'existe pas dans les polices du document)
\newcommand{\popcheck}[1]{\tikz[baseline=-0.5ex]\draw[#1,line width=1.6pt,line cap=round,line join=round](-0.13,0.0)--(-0.04,-0.09)--(0.13,0.1);}
\providecommand{\coche}{\popcheck{popgreen}}
\providecommand{\resultat}[1]{\par\smallskip\noindent\fcolorbox{popgreen}{popgreenL}{\parbox{\dimexpr\linewidth-2\fboxsep-2\fboxrule\relax}{\textbf{Résultat.} #1}}\par\smallskip}
\newcommand{\popboxhead}[3]{\poppill{#1}{#2}{white}\ifx\relax#3\relax\else\hspace{6pt}{\popxbold\fontsize{10.6}{12}\selectfont\color{popink}#3}\fi\par\vspace{7pt}}

\newtcolorbox{aretenir}[1][]{popbase,colframe=popgreen,before upper={\popboxhead{À retenir}{popgreen}{#1}}}
% Texte en allemand (document de lecture, dialogue, extrait) : cadre
% d'étude. Un titre optionnel (source, auteur) s'affiche dans l'en-tête.
\newtcolorbox{textebox}[1][]{popbase,colframe=popblue,before upper={\popboxhead{Text}{popblue}{#1}\begin{otherlanguage}{ngerman}},after upper={\end{otherlanguage}}}
% Marques ✓ / ✗ (forme correcte / fautive) souvent utilisées par les modèles
\providecommand{\cross}{\textcolor{poppink}{\ensuremath{\times}}}
\providecommand{\xmark}{\cross}\providecommand{\crossmark}{\cross}\providecommand{\cmark}{\ensuremath{\checkmark}}
% Ligne de réponse / justification à compléter (inventée par les modèles)
\providecommand{\justification}[1]{\par\noindent\textit{Justification~:} \dotfill\par}
% \textipa (package tipa non chargé) : la police ne couvre pas l'API -> simple passage
\providecommand{\textipa}[1]{#1}
% Soulignement (ulem non chargé) : \uline, \ul
\providecommand{\uline}[1]{\underline{#1}}\providecommand{\ul}[1]{\underline{#1}}
% \glossaire{\item ...} inventée par les modèles : liste de glossaire
\providecommand{\glossaire}[1]{\begin{itemize}#1\end{itemize}}
% \alert (beamer) inventée par les modèles : texte mis en évidence
\providecommand{\alert}[1]{\textbf{\textcolor{poppink}{#1}}}
% variantes ASCII d'ordinaux écrites en macro
\providecommand{\iere}{\textsuperscript{re}}\providecommand{\ieme}{\textsuperscript{e}}\providecommand{\ere}{\textsuperscript{re}}\providecommand{\eme}{\textsuperscript{e}}\providecommand{\er}{\textsuperscript{er}}\providecommand{\ie}{\textsuperscript{e}}
% Ordinaux français écrits en macro par les modèles : 1\ière, 3\ième
\newcommand{\ière}{\textsuperscript{re}}\newcommand{\ième}{\textsuperscript{e}}
% \exemple (inventée par les modèles) : étiquette « Exemple : »
\providecommand{\exemple}{\textbf{Exemple~:}\ }
% \square utilisable aussi hors mode math (cases à cocher)
\AtBeginDocument{\let\squareMath\square\renewcommand{\square}{\ensuremath{\squareMath}}}
% Mot ou phrase en allemand à l'intérieur d'un paragraphe français
\newcommand{\all}[1]{\ifmmode\text{\textit{#1}}\else\begingroup\selectlanguage{ngerman}\itshape #1\endgroup\fi}
\let\esp\all % alias toléré
\newtcolorbox{exemplebox}[1][]{popbase,colframe=poporange,before upper={\popboxhead{Exemple}{poporange}{#1}}}
\newtcolorbox{definition}[1][]{popbase,colframe=popblue,before upper={\popboxhead{Définition}{popblue}{#1}}}
\newtcolorbox{propriete}[1][]{popbase,colframe=poppurple,before upper={\popboxhead{Propriété}{poppurple}{#1}}}
\newtcolorbox{methode}[1][]{popbase,colframe=popgold,before upper={\popboxhead{Méthode}{popgold}{#1}}}
\newtcolorbox{attention}[1][]{popbase,colframe=poppink,before upper={\popboxhead{Attention}{poppink}{#1}}}
\newtcolorbox{experience}[1][]{popbase,colframe=popblue,colback=popblueL,before upper={\popboxhead{Expérience}{popblue}{#1}}}
% « Le savais-tu ? » : carte violette pleine (contexte, culture, Cameroun)
\newtcolorbox{savaistu}[1][]{popbase,colback=poppurple,colframe=popink,
  fontupper=\popsemi\fontsize{10.4}{14.2}\selectfont\color{white},
  before upper={\poppill{Le savais-tu\,?}{white}{poppurple}\ifx\relax#1\relax\else\hspace{6pt}{\popxbold\color{white}#1}\fi\par\vspace{7pt}}}
% Exercice résolu : énoncé en haut, \tcblower puis la solution sur fond crème
\newcounter{popexo}\newcounter{popexr}
\newtcolorbox{exoresolu}[1][]{popbase,colframe=poporange,colbacklower=popcream,
  segmentation style={popink,line width=1.2pt,dashed},
  before upper={\stepcounter{popexr}\popboxhead{Exercice résolu \thepopexr}{poporange}{#1}},
  before lower={\poppill{Solution}{popink}{popcream}\par\vspace{6pt}}}
% Exercice (non résolu en place) — la correction est dans la partie Corrigés
\newtcolorbox{exercice}[1][]{popbase,colframe=popink,
  before upper={\stepcounter{popexo}\popboxhead{Exercice \thepopexo}{popink}{#1}}}
% Corrigé
\newtcolorbox{corrige}[1][]{popbase,colframe=popgreen,colback=popgreenL,
  before upper={\popboxhead{Corrigé}{popgreen}{#1}}}
% Objectifs / prérequis / bilan
\newtcolorbox{objectifs}{popbase,colframe=popink,colback=poporangeL,
  before upper={\popboxhead{Objectifs de la leçon}{popink}{}}}
\newtcolorbox{prerequis}{popbase,colframe=popink,colback=popblueL,
  before upper={\popboxhead{Prérequis}{popblue}{}}}
\newtcolorbox{bilan}{popbase,colframe=popink,colback=white,boxrule=2.8pt,
  before upper={\popboxhead{Fiche bilan}{poporange}{L'essentiel à connaître pour l'examen}}}
% Figure encadrée avec légende
\newtcolorbox{popfigure}[1][]{enhanced,breakable=false,colback=white,colframe=popline,boxrule=1.4pt,arc=8pt,
  left=10pt,right=10pt,top=10pt,bottom=8pt,before skip=10pt,after skip=12pt,halign=center,
  lower separated=false,halign lower=center,
  fontlower=\popsemi\fontsize{9}{11.5}\selectfont\color{popmuted},
  #1}
\newcommand{\legende}[1]{\par\vspace{6pt}{\popsemi\fontsize{9}{11.5}\selectfont\color{popmuted}#1}}

% ============================================================
% Signature OnBuch+
% ============================================================
\newcommand{\poplogoinline}[1]{{\poplogo\fontsize{#1}{#1}\selectfont\color{popink}OnBuch\color{poporange}+}}
\newcommand{\signatureonbuch}{%
  \begin{tcolorbox}[enhanced,colback=popink,colframe=popink,boxrule=2.2pt,arc=10pt,
      drop shadow=poporange,left=14pt,right=14pt,top=10pt,bottom=10pt,before skip=0pt,after skip=0pt]
    \tikz[baseline=-0.7ex]\node[circle,fill=popgreen,draw=popcream,line width=1.3pt,minimum size=22pt,inner sep=0]{\popcheck{white}};%
    \hspace{9pt}\begin{minipage}[c]{0.62\linewidth}
      {\popxbold\fontsize{12}{13}\selectfont\color{popcream}Signé\ \textbullet\ {\poplogo OnBuch\color{poporange}+}}\par\vspace{2pt}
      {\raggedright\popsemi\fontsize{8}{10.5}\selectfont\color{popline}\MakeUppercase{Équipe pédagogique OnBuch+}\\\MakeUppercase{Conforme au programme officiel MINESEC}\par}
    \end{minipage}\hfill
    {\poplogo\fontsize{22}{22}\selectfont\color{popcream}OnBuch\color{poporange}+}
  \end{tcolorbox}%
}

% ============================================================
% Page de garde
% #1 titre · #2 matière · #3 niveau · #4 module · #5 n° leçon · #6 durée ·
% #7 description / objectifs (texte ou itemize)
% ============================================================
\newcommand{\pagegarde}[7]{%
  \thispagestyle{empty}
  \newgeometry{a4paper,margin=1.7cm,top=1.6cm,bottom=1.4cm}%
  \begin{tikzpicture}[remember picture,overlay]
    \fill[poporange!18] ([xshift=-3.2cm,yshift=-3.6cm]current page.north east) circle (4.6cm);
    \fill[poppurple!14] ([xshift=2.4cm,yshift=7.6cm]current page.south west) circle (3.8cm);
  \end{tikzpicture}%
  {\poplogo\fontsize{30}{30}\selectfont\color{popink}OnBuch\color{poporange}+}\hfill
  \raisebox{0.6ex}{\poppill{Cours signé \textbullet\ Premium}{popink}{popcream}}\par
  \vspace{1pt}\popsquiggle{poppurple}\par
  \vspace{8pt}
  \begin{tcolorbox}[enhanced,colback=poporange,colframe=popink,boxrule=2.8pt,arc=13pt,
      drop shadow=popink,left=18pt,right=18pt,top=12pt,bottom=13pt,after skip=10pt]
    \poppill{#2 \textbullet\ #3}{popink}{popcream}\hspace{5pt}\poppill{#4}{white}{popink}\par\vspace{8pt}
    {\popdisplay\fontsize{14}{14}\selectfont\color{popink}LEÇON #5}\par\vspace{4pt}
    {\raggedright\hyphenpenalty=10000\exhyphenpenalty=10000\popdisplay\fontsize{25}{27}\selectfont\color{popcream}\MakeUppercase{#1}\par}
    \vspace{8pt}
    {\popxbold\fontsize{9.5}{9.5}\selectfont\color{popink}\MakeUppercase{Durée conseillée : #6}}
  \end{tcolorbox}
  \begin{tcolorbox}[enhanced,colback=white,colframe=popink,boxrule=2.2pt,arc=10pt,
      drop shadow=popink,left=16pt,right=16pt,top=10pt,bottom=10pt,after skip=8pt]
    \poppill{Dans cette leçon}{poppurple}{white}\par\vspace{5pt}
    \popsemi\fontsize{9.3}{12.6}\selectfont\color{popink}#7
  \end{tcolorbox}
  \noindent\begin{minipage}[t]{0.487\linewidth}
    \begin{tcolorbox}[enhanced,colback=popgreen,colframe=popink,boxrule=2.2pt,arc=10pt,
        drop shadow=popink,left=11pt,right=11pt,top=10pt,bottom=10pt,height=3.1cm,valign=center]
      \tcbox[colback=white,colframe=popink,boxrule=1.3pt,arc=5pt,boxsep=0pt,nobeforeafter,
        left=3pt,right=3pt,top=3pt,bottom=3pt,valign=center]{\qrcode[height=1.9cm]{https://play.google.com/store/apps/details?id=com.luvvix.onbuch&hl=fr}}%
      \hspace{8pt}\begin{minipage}[c]{0.5\linewidth}
        {\popdisplay\fontsize{11}{12.5}\selectfont\color{white}\MakeUppercase{Télécharge OnBuch}}\par\vspace{4pt}
        {\raggedright\popsemi\fontsize{8.4}{11}\selectfont\color{white}Gratuit \textbullet\ Google Play\\Tuteur IA, annales, concours, résultats\par}
      \end{minipage}
    \end{tcolorbox}
  \end{minipage}\hfill
  \begin{minipage}[t]{0.487\linewidth}
    \begin{tcolorbox}[enhanced,colback=poppurple,colframe=popink,boxrule=2.2pt,arc=10pt,
        drop shadow=popink,left=11pt,right=11pt,top=10pt,bottom=10pt,height=3.1cm,valign=center]
      \tikz[baseline=-0.7ex]\node[circle,fill=white,draw=popink,line width=1.3pt,minimum size=1.6cm,inner sep=0]{\popdisplay\fontsize{24}{24}\selectfont\color{poppurple}+};%
      \hspace{8pt}\begin{minipage}[c]{0.6\linewidth}
        {\popdisplay\fontsize{11}{12.5}\selectfont\color{white}\MakeUppercase{Passe à OnBuch+}}\par\vspace{4pt}
        {\raggedright\popsemi\fontsize{8.4}{11}\selectfont\color{white}Tous les cours signés, Tuteur IA illimité, fiches PDF — onglet OnBuch+ de l'app\par}
      \end{minipage}
    \end{tcolorbox}
  \end{minipage}
  \vspace{8pt}
  \popcontenu
  \vfill
  \signatureonbuch
  \restoregeometry
  \newpage
}

% Rangée « Ce cours contient » (page de garde)
\newcommand{\poptile}[3]{% #1 couleur #2 gros texte #3 légende
  \begin{tcolorbox}[enhanced,colback=#1,colframe=popink,boxrule=2pt,arc=9pt,shadow={2.5pt}{-2.5pt}{0pt}{popink},
      left=6pt,right=6pt,top=9pt,bottom=9pt,halign=center,valign=center,height=1.75cm,width=\linewidth,nobeforeafter]
    {\popdisplay\fontsize{12.5}{13}\selectfont\color{white}\MakeUppercase{#2}}\par\vspace{3pt}
    {\centering\popsemi\fontsize{7.8}{9.5}\selectfont\color{white}#3\par}
  \end{tcolorbox}}
\newcommand{\popcontenu}{%
  \par\noindent{\popxboldls\fontsize{7.5}{7.5}\selectfont\color{popmuted}CE COURS SIGNÉ CONTIENT}\par\vspace{7pt}
  \noindent\begin{minipage}[t]{0.235\linewidth}\poptile{popblue}{Cours}{complet, illustré, pas à pas}\end{minipage}\hfill
  \begin{minipage}[t]{0.235\linewidth}\poptile{poporange}{Méthodes}{et exercices résolus}\end{minipage}\hfill
  \begin{minipage}[t]{0.235\linewidth}\poptile{poppink}{Exercices}{type examen + corrigés}\end{minipage}\hfill
  \begin{minipage}[t]{0.235\linewidth}\poptile{popgreen}{Fiche bilan}{l'essentiel à retenir}\end{minipage}\par}

% Sommaire
\newtcolorbox{sommaire}{enhanced,colback=white,colframe=popink,boxrule=2.2pt,arc=10pt,
  drop shadow=popink,left=18pt,right=18pt,top=13pt,bottom=13pt,before skip=6pt,after skip=20pt,
  before upper={\poppill{Sommaire}{poppurple}{white}\par\vspace{9pt}\popsemi\fontsize{10.6}{14.5}\selectfont\color{popink}}}

% Signature de fin de cours — poussée tout en bas de la dernière page
\newcommand{\findecours}{%
  \par\vfill\nopagebreak
  \begin{center}\popsquiggle{poporange}\end{center}\vspace{4pt}
  \signatureonbuch
}
\AtBeginDocument{\def\llbracket{\mathopen{[\mkern-3.5mu[}}\def\rrbracket{\mathclose{]\mkern-3.5mu]}}} % intervalles d'entiers (glyphe ⟦ absent de la police)
% Glyphes absents de Fira Math : redéfinis à partir de glyphes disponibles
\AtBeginDocument{%
  \def\setminus{\mathbin{\backslash}}\let\smallsetminus\setminus
  \def\vdots{\mathord{\vbox{\baselineskip3.2pt\lineskiplimit0pt\kern4pt\hbox{.}\hbox{.}\hbox{.}}}}%
  \def\ddots{\mathinner{\mkern1mu\raise6.5pt\hbox{.}\mkern2mu\raise3.6pt\hbox{.}\mkern2mu\raise0.7pt\hbox{.}\mkern1mu}}%
  \def\triangle{\mathord{\scalebox{0.9}{$\symup{\Delta}$}}}%
  \def\nabla{\mathord{\raisebox{\depth}{\scalebox{1}[-1]{$\symup{\Delta}$}}}}%
  \def\checkmark{\popcheck{popdark}}%
  \def\star{\mathbin{\ast}}\def\bigstar{\mathord{\ast}}%
  \def\diamond{\mathbin{\scalebox{0.6}{$\Diamond$}}}%
  \def\bigcup{\mathop{\mathchoice{\raisebox{-0.3em}{\scalebox{1.7}{$\cup$}}}{\scalebox{1.25}{$\cup$}}{\cup}{\cup}}\displaylimits}%
  \def\bigcap{\mathop{\mathchoice{\raisebox{-0.3em}{\scalebox{1.7}{$\cap$}}}{\scalebox{1.25}{$\cap$}}{\cap}{\cap}}\displaylimits}%
  \def\lesseqqgtr{\mathrel{\lessgtr}}\def\bigoplus{\mathop{\oplus}\displaylimits}%
}
\providecommand{\ssi}{si et seulement si} % abréviation fréquente
\AtBeginDocument{\providecommand{\wideparen}{\overparen}} % arc de cercle

%% FIGFIX-BEGIN v5 — correctifs globaux des figures (docs/onbuch-generation/tools/figfix)
\usepackage{adjustbox}\usepackage{etoolbox}\tcbuselibrary{raster}
% toute figure TikZ/pgfplots plus large que la ligne est réduite à la largeur disponible
\BeforeBeginEnvironment{tikzpicture}{\begin{adjustbox}{max width=\linewidth}}
\AfterEndEnvironment{tikzpicture}{\end{adjustbox}}
% légendes pgfplots lisibles par-dessus les courbes
\pgfplotsset{every axis/.append style={legend style={fill=white,fill opacity=0.92,text opacity=1,draw=black!15,inner sep=3pt}}}
% étiquettes d'arêtes (syntaxe edge["..."]) avec fond blanc
\tikzset{every edge quotes/.append style={fill=white,inner sep=1.2pt}}
%% FIGFIX-END
\begin{document}

\pagegarde{Les problèmes d'émigration}{Allemand}{Tle A}{Chapitre 5 : Citoyenneté}{AL28}{2 h}{Cette leçon de type « texte » propose la lecture et le commentaire d'un texte allemand sur les problèmes d'émigration, enrichie d'un lexique thématique complet. Elle consolide la grammaire de la cause et du but, le parfait et le prétérit, ainsi que le subjonctif II, et aboutit à une production guidée.
\vspace{4pt}
\begin{multicols}{2}\small
\begin{itemize}
\item À la fin de cette leçon, je sais comprendre globalement puis en détail un texte allemand sur les problèmes d'émigration
\item Je sais employer correctement le vocabulaire thématique (die Auswanderung, die Flucht, das Visum, die Grenze, die Arbeitslosigkeit, das Ausland) avec articles et pluriels
\item Je sais exprimer la cause avec weil, denn et da en respectant la place du verbe
\item Je sais exprimer le but avec damit et um ... zu et choisir la bonne construction
\item Je sais distinguer et employer le parfait et le prétérit selon le contexte
\item Je sais former et utiliser le subjonctif II pour exprimer un souhait ou une hypothèse liés à l'émigration
\end{itemize}\end{multicols}}

\begin{sommaire}
\begin{multicols}{2}
\begin{enumerate}
\item Texte support : « Warum so viele junge Leute auswandern »
\item Compréhension du texte
\item Lexique thématique : l'émigration
\item Grammaire 1 : la cause et le but
\item Grammaire 2 : parfait, prétérit et subjonctif II
\item Commentaire et production guidée
\item Activité d'intégration
\item Méthodes et exercices résolus
\item Exercices
\item Corrigés des exercices
\item Fiche bilan
\end{enumerate}
\end{multicols}
\end{sommaire}

\begin{objectifs}
\begin{itemize}
\item À la fin de cette leçon, je sais comprendre globalement puis en détail un texte allemand sur les problèmes d'émigration
\item Je sais employer correctement le vocabulaire thématique (die Auswanderung, die Flucht, das Visum, die Grenze, die Arbeitslosigkeit, das Ausland) avec articles et pluriels
\item Je sais exprimer la cause avec weil, denn et da en respectant la place du verbe
\item Je sais exprimer le but avec damit et um ... zu et choisir la bonne construction
\item Je sais distinguer et employer le parfait et le prétérit selon le contexte
\item Je sais former et utiliser le subjonctif II pour exprimer un souhait ou une hypothèse liés à l'émigration
\item Je sais répondre à des questions de compréhension et résumer un texte en quelques phrases
\item Je sais rédiger un court paragraphe argumenté sur les raisons et les buts de l'émigration
\end{itemize}
\end{objectifs}

\begin{prerequis}
\begin{itemize}
\item La place du verbe dans la phrase allemande (V2, verbe en fin de subordonnée)
\item Le parfait avec haben et sein (Première)
\item Les pronoms personnels et les articles définis (der, die, das)
\item Le vocabulaire général du travail et des études (die Arbeit, die Schule, studieren)
\item Les connecteurs simples (und, aber, oder)
\end{itemize}
\end{prerequis}

\newpage

\coursec{Texte support : « Warum so viele junge Leute auswandern »}

Chaque année, des milliers de jeunes Camerounais rêvent de partir étudier ou travailler en Allemagne, en Autriche ou en Suisse. Mais entre le visa refusé, la langue à apprendre et le chômage qui pousse au départ, l'émigration est un parcours semé d'obstacles. Dans les pays germanophones aussi, on débat beaucoup de l'immigration et de l'intégration. Comment parler de ces réalités en allemand ? Comment expliquer \emph{pourquoi} on part (\all{weil}, \all{denn}, \all{da}) et \emph{dans quel but} (\all{damit}, \all{um ... zu}) ? C'est tout l'objet de cette leçon, qui commence par la lecture d'un article de presse.

\textbf{Problématique :} Pourquoi tant de jeunes quittent-ils leur pays, quels obstacles rencontrent-ils, et comment exprimer en allemand les causes et les buts de l'émigration ?

\courssub{Première approche : lecture globale du texte}

\begin{methode}[Lire un texte en trois étapes]
\begin{enumerate}
\item \textbf{Lecture globale sans dictionnaire :} lire le texte une première fois en silence, sans s'arrêter sur les mots inconnus. Objectif : identifier le \cle{thème}, le \cle{type de texte} et l'idée générale.
\item \textbf{Repérage des mots-clés et des connecteurs :} souligner les mots du champ lexical du thème (ici : l'émigration) et entourer les connecteurs logiques (cause, but, opposition).
\item \textbf{Lecture détaillée avec le glossaire :} relire phrase par phrase en s'aidant du glossaire, puis répondre aux questions de compréhension.
\end{enumerate}
\end{methode}

Avant de lire, observons le titre : \all{Warum wandern so viele junge Leute aus?} — « Pourquoi tant de jeunes émigrent-ils ? ». Le mot interrogatif \all{warum} annonce un texte qui donne des \cle{raisons}, des \cle{causes} : on peut donc s'attendre à trouver des connecteurs de cause. Le verbe \all{auswandern} (émigrer) annonce le thème. Le texte est un \cle{article de presse} : le locuteur est un journaliste qui décrit un phénomène de société de manière objective.

\begin{textebox}[Warum wandern so viele junge Leute aus?]
Immer mehr junge Menschen verlassen ihre Heimat, weil sie dort keine Arbeit finden. Die Arbeitslosigkeit ist hoch, und viele Familien sind arm. Da die wirtschaftliche Situation schwierig ist, beantragen viele junge Leute ein Visum, um im Ausland zu studieren oder zu arbeiten.

Viele Jugendliche träumen von einem Leben in Europa. Sie lernen Deutsch, damit sie sich in Deutschland, Österreich oder der Schweiz integrieren können. Denn ohne Sprachkenntnisse ist es fast unmöglich, einen guten Job zu finden. Manche besuchen Sprachkurse, andere lernen allein mit dem Handy.

Aber die Flucht ins Ausland ist nicht einfach: An der Grenze warten Kontrollen, und ohne Visum ist die Einreise unmöglich. Die Botschaften bekommen jedes Jahr Tausende von Anträgen, aber sie können nicht alle akzeptieren. Viele junge Menschen warten monatelang auf eine Antwort.

Viele Auswanderer haben ihre Familien verlassen und leben allein in einem fremden Land. Sie vermissen ihre Eltern, ihre Freunde und das Essen ihrer Heimat. Trotzdem hoffen sie auf eine bessere Zukunft. Sie arbeiten hart, damit ihre Kinder einmal ein besseres Leben haben. Einige schicken auch Geld an ihre Familien, denn sie wollen ihnen helfen.

Die Regierungen der Herkunftsländer müssen mehr Arbeitsplätze schaffen, damit die jungen Leute nicht weggehen. Denn ein Land ohne Jugend ist ein Land ohne Zukunft.
\end{textebox}

\textbf{Glossaire des mots difficiles :}

\begin{center}
\begin{tabularx}{\linewidth}{|X|X|}\hline
\rowcolor{popblueL} \textbf{Deutsch} & \textbf{Français} \\ \hline
die Heimat, -en & la patrie, le pays natal \\ \hline
die Arbeitslosigkeit, -en & le chômage \\ \hline
beantragen (trennbar) & demander officiellement, déposer une demande de \\ \hline
das Visum, die Visa & le visa \\ \hline
das Ausland (Sg.) & l'étranger \\ \hline
die Sprachkenntnisse (Pl.) & les connaissances linguistiques \\ \hline
die Grenze, -n & la frontière \\ \hline
die Einreise, -n & l'entrée (dans un pays) \\ \hline
die Botschaft, -en & l'ambassade \\ \hline
der Antrag, die Anträge & la demande officielle \\ \hline
vermissen & manquer (quelqu'un, quelque chose) \\ \hline
der Arbeitsplatz, die Arbeitsplätze & l'emploi, le poste de travail \\ \hline
\end{tabularx}
\end{center}

\courssub{Compréhension globale : de quoi parle le texte ?}

Après la première lecture silencieuse, l'enseignant lit le texte à voix haute (compréhension orale) : les élèves écoutent sans suivre, puis vérifient à la deuxième écoute. On peut alors dégager l'idée générale.

\begin{exemplebox}[Reformulation orale du contenu général (2-3 phrases)]
\all{Der Text erklärt, warum viele junge Menschen ihre Heimat verlassen: Sie finden keine Arbeit und wollen im Ausland studieren oder arbeiten. Die Auswanderung ist aber schwierig, weil man ein Visum braucht und die Sprache lernen muss. Trotzdem hoffen die Auswanderer auf eine bessere Zukunft.}

« Le texte explique pourquoi beaucoup de jeunes quittent leur pays : ils ne trouvent pas de travail et veulent étudier ou travailler à l'étranger. Mais l'émigration est difficile, car il faut un visa et apprendre la langue. Malgré tout, les émigrants espèrent un avenir meilleur. »
\end{exemplebox}

\begin{definition}[Auswanderung ou Einwanderung ? Une question de point de vue]
\all{Die Auswanderung} (-, -en) : le fait de \textbf{quitter} son pays pour s'installer à l'étranger = l'\emph{émigration} (point de vue du pays de départ).

\all{Die Einwanderung} (-, -en) : le fait d'\textbf{entrer} dans un pays pour s'y installer = l'\emph{immigration} (point de vue du pays d'accueil).

Les verbes correspondants : \all{auswandern} (émigrer) et \all{einwandern} (immigrer). Le migrant est donc à la fois \all{Auswanderer} (pour son pays) et \all{Einwanderer} (pour le pays d'accueil).
\end{definition}

\begin{attention}[Ne pas confondre auswandern et einwandern]
Le préfixe change le point de vue : \all{aus-} = « hors de » (on sort de son pays), \all{ein-} = « dans » (on entre dans un pays). Erreur fréquente des francophones : traduire « immigrer en Allemagne » par \all{auswandern}. Non !

\all{Viele Kameruner wandern nach Deutschland aus.} « Beaucoup de Camerounais émigrent vers l'Allemagne. »

\all{Viele Kameruner wandern in Deutschland ein.} « Beaucoup de Camerounais immigrent en Allemagne. »

Autre piège : \all{die Flucht} signifie « la fuite » (fuir la guerre, la misère), pas « le vol » ; et \all{die Botschaft} est « l'ambassade », pas « le message » (qui se dit \all{die Nachricht}).
\end{attention}

\courssub{Repérage : le champ lexical de l'émigration dans le texte}

Cherchons maintenant dans le texte tous les mots qui appartiennent au champ lexical de l'émigration. Cette étape prépare le travail de vocabulaire de la section 3.

\begin{center}
\begin{tabularx}{\linewidth}{|X|X|X|}\hline
\rowcolor{popblueL} \textbf{Mot du texte} & \textbf{Français} & \textbf{Famille de mots} \\ \hline
die Auswanderung & l'émigration & auswandern, der Auswanderer \\ \hline
verlassen (verlässt, verließ, hat verlassen) & quitter & der Abschied (les adieux) \\ \hline
die Heimat & la patrie & das Herkunftsland (pays d'origine) \\ \hline
die Arbeitslosigkeit & le chômage & arbeitslos, der Arbeitslose \\ \hline
das Visum & le visa & beantragen, der Antrag \\ \hline
das Ausland & l'étranger & der Ausländer, ausländisch \\ \hline
die Grenze & la frontière & die Kontrolle, die Einreise \\ \hline
die Flucht & la fuite & fliehen, der Flüchtling (réfugié) \\ \hline
sich integrieren & s'intégrer & die Integration \\ \hline
die Zukunft & l'avenir & zukünftig (futur, adj.) \\ \hline
\end{tabularx}
\end{center}

\begin{popfigure}
\begin{center}\tcbox[colback=popink,colframe=popink,coltext=white,arc=9pt,boxsep=4pt,left=6pt,right=6pt]{\bfseries\small die Auswanderung}\end{center}
\vspace{-2pt}
\begin{tcbraster}[raster columns=2,raster equal height=rows,raster column skip=6pt,raster row skip=6pt,enhanced,blankest]
\begin{tcolorbox}[enhanced,colback=white,colframe=poporange,boxrule=0.9pt,arc=3pt,title={Gründe},fonttitle=\bfseries\scriptsize,coltitle=white,colbacktitle=poporange,left=4pt,right=4pt,top=2pt,bottom=2pt,boxsep=1pt,toptitle=1.5pt,bottomtitle=1.5pt,halign=left]
\begin{itemize}[nosep,leftmargin=*,itemsep=1pt,font=\scriptsize]
\item Arbeitslosigkeit
\item Armut
\end{itemize}
\end{tcolorbox}
\begin{tcolorbox}[enhanced,colback=white,colframe=popgreen!80!black,boxrule=0.9pt,arc=3pt,title={Ziele},fonttitle=\bfseries\scriptsize,coltitle=white,colbacktitle=popgreen!80!black,left=4pt,right=4pt,top=2pt,bottom=2pt,boxsep=1pt,toptitle=1.5pt,bottomtitle=1.5pt,halign=left]
\begin{itemize}[nosep,leftmargin=*,itemsep=1pt,font=\scriptsize]
\item studieren
\item arbeiten
\end{itemize}
\end{tcolorbox}
\begin{tcolorbox}[enhanced,colback=white,colframe=poppurple,boxrule=0.9pt,arc=3pt,title={Probleme},fonttitle=\bfseries\scriptsize,coltitle=white,colbacktitle=poppurple,left=4pt,right=4pt,top=2pt,bottom=2pt,boxsep=1pt,toptitle=1.5pt,bottomtitle=1.5pt,halign=left]
\begin{itemize}[nosep,leftmargin=*,itemsep=1pt,font=\scriptsize]
\item das Visum
\item die Sprache
\end{itemize}
\end{tcolorbox}
\begin{tcolorbox}[enhanced,colback=white,colframe=popgold!90!black,boxrule=0.9pt,arc=3pt,title={Lösungen},fonttitle=\bfseries\scriptsize,coltitle=white,colbacktitle=popgold!90!black,left=4pt,right=4pt,top=2pt,bottom=2pt,boxsep=1pt,toptitle=1.5pt,bottomtitle=1.5pt,halign=left]
\begin{itemize}[nosep,leftmargin=*,itemsep=1pt,font=\scriptsize]
\item Arbeitsplätze
\item Sprachkurse
\end{itemize}
\end{tcolorbox}
\end{tcbraster}

\legende{Carte mentale : le champ lexical de l'émigration (die Auswanderung)}
\end{popfigure}

\courssub{Repérage : les connecteurs de cause et de but dans le texte}

Le titre pose la question \all{warum} (pourquoi) : le texte répond avec des connecteurs de \cle{cause} et de \cle{but}. Soulignons-les.

\begin{exemplebox}[Connecteurs relevés dans le texte]
\textbf{La cause :}
\begin{itemize}
\item \all{Immer mehr junge Menschen verlassen ihre Heimat, \textbf{weil} sie dort keine Arbeit finden.} « ... parce qu'ils n'y trouvent pas de travail. » (verbe à la fin)
\item \all{\textbf{Da} die wirtschaftliche Situation schwierig ist, beantragen viele ein Visum.} « Comme la situation économique est difficile, ... » (verbe à la fin)
\item \all{\textbf{Denn} ohne Sprachkenntnisse ist es fast unmöglich, einen guten Job zu finden.} « Car sans connaissances linguistiques, ... » (verbe en position 2)
\end{itemize}
\textbf{Le but :}
\begin{itemize}
\item \all{Sie beantragen ein Visum, \textbf{um} im Ausland \textbf{zu} studieren.} « ... pour étudier à l'étranger. »
\item \all{Sie lernen Deutsch, \textbf{damit} sie sich integrieren können.} « ... afin qu'ils puissent s'intégrer. »
\end{itemize}
\end{exemplebox}

\begin{aretenir}[Première observation grammaticale]
\all{weil} et \all{da} renvoient le verbe conjugué \textbf{à la fin} de la subordonnée ; \all{denn} ne change pas l'ordre des mots (le verbe reste en position 2). \all{um ... zu} + infinitif s'emploie quand le sujet des deux verbes est le même ; \all{damit} + subordonnée s'emploie quand les sujets sont différents. Ces règles seront étudiées en détail dans la section 4.
\end{aretenir}

\begin{exoresolu}[Questions de compréhension globale]
\textbf{Énoncé.} Répondez en allemand, par des phrases complètes :
\begin{enumerate}
\item Was für ein Text ist das? (ein Gedicht, ein Brief, ein Zeitungsartikel?)
\item Warum verlassen viele junge Menschen ihre Heimat?
\item Was machen die jungen Leute, um sich im Ausland zu integrieren?
\item Warum ist die Einreise ohne Visum unmöglich?
\end{enumerate}
\tcblower
\textbf{Solution détaillée.}
\begin{enumerate}
\item \all{Das ist ein Zeitungsartikel.} « C'est un article de journal. » (Le titre interrogatif, le ton objectif et le thème de société sont typiques de la presse.)
\item \all{Sie verlassen ihre Heimat, weil sie dort keine Arbeit finden und weil viele Familien arm sind.} « Ils quittent leur pays parce qu'ils n'y trouvent pas de travail et parce que beaucoup de familles sont pauvres. » (Reprise du connecteur \all{weil} du texte : verbe \all{finden} à la fin.)
\item \all{Sie lernen Deutsch, damit sie sich integrieren können.} « Ils apprennent l'allemand afin de pouvoir s'intégrer. »
\item \all{Ohne Visum ist die Einreise unmöglich, denn an der Grenze warten Kontrollen.} « Sans visa, l'entrée est impossible, car il y a des contrôles à la frontière. » (Ici \all{denn} : le verbe \all{warten} reste en position 2.)
\end{enumerate}
\end{exoresolu}

\begin{savaistu}[Willkommenskultur : la culture de l'accueil]
Depuis 2015, l'Allemagne est l'un des premiers pays d'accueil de migrants en Europe. Le mot \all{Willkommenskultur} (« culture de l'accueil ») est alors devenu célèbre dans le monde entier : il désigne l'attitude d'ouverture envers les réfugiés et les migrants. Le mot a même été choisi « mot de l'année » en Allemagne. Saviez-vous aussi que le Cameroun fut une colonie allemande de 1884 à 1916 ? C'est l'une des raisons pour lesquelles l'allemand est aujourd'hui enseigné dans les lycées camerounais.
\end{savaistu}

\begin{exercice}[À vous de jouer]
\begin{enumerate}
\item Retrouvez dans le texte trois mots de la famille de \all{arbeiten} et traduisez-les.
\item Recopiez deux phrases du texte contenant un connecteur de cause et soulignez le verbe conjugué : où se trouve-t-il ?
\item Reformulez oralement, en deux phrases allemandes, l'idée générale du texte avec vos propres mots.
\end{enumerate}
\end{exercice}

\begin{corrige}[Exercice : À vous de jouer]
\begin{enumerate}
\item \all{die Arbeit} (le travail), \all{die Arbeitslosigkeit} (le chômage), \all{der Arbeitsplatz} (l'emploi, le poste).
\item Exemples : \all{... weil sie dort keine Arbeit \textbf{finden}.} et \all{\textbf{Da} die Situation schwierig \textbf{ist}, ...} — dans les deux cas, le verbe conjugué est à la \textbf{fin} de la subordonnée introduite par \all{weil} ou \all{da}.
\item Modèle : \all{Viele junge Leute wandern aus, weil sie in der Heimat keine Arbeit finden. Sie hoffen auf eine bessere Zukunft im Ausland.}
\end{enumerate}
\end{corrige}

\coursec{Compréhension du texte}

Dans la section précédente, nous avons lu et découvert le texte support \all{„Warum so viele junge Leute auswandern“}, qui décrit les raisons pour lesquelles de nombreux jeunes quittent leur pays, les démarches qu'ils entreprennent (demande de visa, apprentissage de la langue) et les difficultés qu'ils rencontrent (frontières, intégration, nostalgie). Nous allons maintenant exploiter ce texte de manière méthodique : d'abord une compréhension \cle{globale}, puis une compréhension \cle{détaillée}, ensuite des exercices de vrai/faux, un travail sur les connecteurs de cause et de but, un résumé guidé et enfin une discussion orale. Cette démarche est exactement celle exigée à l'épreuve d'allemand du Baccalauréat : lire, comprendre, justifier, résumer, s'exprimer.

\courssub{La démarche de lecture : du global au détaillé}

\begin{popfigure}
\begin{tikzpicture}[
  node distance=0.55cm and 0.9cm,
  etape/.style={rectangle, rounded corners=3pt, draw=popblue, fill=popblueL, thick, align=center, minimum width=3.4cm, minimum height=1.1cm, font=\small},
  fleche/.style={-{Stealth[length=3mm]}, thick, popdark}
]
\node[etape, align=center] (l1){1. Lecture globale\\ \emph{Worum geht es?}};
\node[etape, right=of l1, align=center] (l2){2. Lecture détaillée\\ \emph{Warum? Was? Wie?}};
\node[etape, right=of l2, align=center] (l3){3. Résumé\\ \emph{Zusammenfassung}};
\node[etape, right=of l3, align=center] (l4){4. Discussion\\ \emph{Diskussion}};
\draw[fleche] (l1) -- (l2);
\draw[fleche] (l2) -- (l3);
\draw[fleche] (l3) -- (l4);
\end{tikzpicture}
\legende{Les quatre étapes de la compréhension d'un texte : lecture globale, lecture détaillée, résumé, discussion.}
\end{popfigure}

\begin{methode}[Répondre à une question de compréhension]
\begin{enumerate}
\item \textbf{Repérer les mots-clés de la question} : l'interrogatif (\all{warum}, \all{was}, \all{wer}, \all{wo}, \all{wohin}, \all{wie}) indique le type d'information attendu.
\item \textbf{Localiser le passage du texte} qui contient la réponse et le souligner.
\item \textbf{Reprendre les mots de la question} dans la réponse : à \all{Warum beantragen viele junge Leute ein Visum?} on répond \all{Viele junge Leute beantragen ein Visum, weil \ldots} et non par une phrase sans rapport avec la question.
\item \textbf{Citer le passage du texte} qui justifie la réponse, entre guillemets allemands : \all{Im Text steht: „\ldots“}.
\item \textbf{Formuler une phrase complète}, avec le verbe conjugué à la deuxième position (ou en finale après \all{weil}, \all{da}, \all{damit}).
\end{enumerate}
\end{methode}

\begin{aretenir}[Les mots interrogatifs à connaître]
\begin{center}
\begin{tabularx}{\linewidth}{|l|l|X|}
\hline
\rowcolor{popblueL} Allemand & Français & Exemple \\
\hline
\all{wer} & qui & \all{Wer will auswandern?} « Qui veut émigrer ? » \\
\hline
\all{was} & quoi, que & \all{Was lernen sie?} « Qu'apprennent-ils ? » \\
\hline
\all{warum} & pourquoi & \all{Warum verlassen sie die Heimat?} « Pourquoi quittent-ils leur pays ? » \\
\hline
\all{wo} & où (lieu) & \all{Wo arbeiten sie?} « Où travaillent-ils ? » \\
\hline
\all{wohin} & où (direction) & \all{Wohin wandern sie aus?} « Où émigrent-ils ? » \\
\hline
\all{wie} & comment & \all{Wie bekommt man ein Visum?} « Comment obtient-on un visa ? » \\
\hline
\end{tabularx}
\end{center}
Attention à la distinction \all{wo} (lieu où l'on est) / \all{wohin} (lieu vers lequel on va) : le français « où » traduit les deux.
\end{aretenir}

\courssub{Compréhension globale : de quoi parle le texte ?}

La lecture globale répond à trois questions simples : \emph{de quoi} parle le texte (\all{Worum geht es?}), \emph{qui} est concerné (\all{Wer?}) et \emph{où} se situe l'action (\all{Wo? Wohin?}).

\begin{exemplebox}[Questions de compréhension globale et réponses modèles]
\begin{enumerate}
\item \all{Worum geht es in diesem Text?} — \all{Der Text handelt von der Auswanderung junger Leute und von ihren Problemen.} « De quoi parle ce texte ? — Le texte traite de l'émigration des jeunes et de leurs problèmes. »
\item \all{Wer ist betroffen?} — \all{Betroffen sind vor allem junge Leute, die ihre Heimat verlassen wollen.} « Qui est concerné ? — Ce sont surtout des jeunes qui veulent quitter leur patrie. »
\item \all{Wohin wandern die jungen Leute aus?} — \all{Sie wandern ins Ausland aus, zum Beispiel nach Deutschland, nach Österreich oder in die Schweiz.} « Où émigrent-ils ? — Ils émigrent à l'étranger, par exemple en Allemagne, en Autriche ou en Suisse. »
\end{enumerate}
\end{exemplebox}

\begin{attention}[« ins Ausland » ou « im Ausland » ?]
\all{ins Ausland} (in + accusatif) exprime la \textbf{direction} : \all{Sie gehen ins Ausland.} « Ils partent à l'étranger. » \all{im Ausland} (in + datif) exprime le \textbf{lieu} : \all{Sie arbeiten im Ausland.} « Ils travaillent à l'étranger. » Même logique que \all{in die Schule gehen} / \all{in der Schule sein}.
\end{attention}

\begin{aretenir}[Les mots de la thématique à reconnaître dès la première lecture]
\all{die Auswanderung, -en} (l'émigration), \all{auswandern} (émigrer, verbe à particule séparable : \all{er wandert aus} ; au parfait : \all{er ist ausgewandert}), \all{der Auswanderer, -} (l'émigrant), \all{die Heimat} (la patrie, le pays natal), \all{das Ausland} (l'étranger), \all{das Visum, die Visa} (le visa), \all{die Grenze, -n} (la frontière), \all{die Arbeitslosigkeit} (le chômage), \all{die Flucht} (la fuite), \all{der Flüchtling, -e} (le réfugié), \all{die Botschaft, -en} (l'ambassade).
\end{aretenir}

\courssub{Compréhension détaillée : pourquoi partent-ils ? Quels obstacles ?}

La lecture détaillée exige de repérer les informations précises : les \cle{causes} du départ et les \cle{obstacles} rencontrés. Observons d'abord un échange type.

\begin{exemplebox}[Exemple traduit]
\all{Warum verlassen die jungen Leute ihre Heimat? — Weil sie dort keine Arbeit finden.}\\
« Pourquoi les jeunes quittent-ils leur patrie ? — Parce qu'ils n'y trouvent pas de travail. »\\
On observe : la question commence par \all{warum} (verbe conjugué en 2\up{e} position) ; la réponse commence par \all{weil} et le verbe conjugué \all{finden} se place \textbf{en fin de phrase}.
\end{exemplebox}

\begin{textebox}[Questions types sur le texte]
1. Warum beantragen viele junge Leute ein Visum?\\
2. Was lernen sie, damit sie sich integrieren können?\\
3. Welche Probleme gibt es an der Grenze?
\end{textebox}

\textbf{Réponses modèles commentées.}\par
1. \all{Viele junge Leute beantragen ein Visum, weil sie im Ausland arbeiten oder studieren wollen.} « Beaucoup de jeunes demandent un visa parce qu'ils veulent travailler ou étudier à l'étranger. » — La cause est introduite par \all{weil}, verbe conjugué \all{wollen} en finale.\par
2. \all{Sie lernen die deutsche Sprache, damit sie sich integrieren können.} « Ils apprennent la langue allemande afin de pouvoir s'intégrer. » — Le but est introduit par \all{damit}, verbe conjugué \all{können} en finale.\par
3. \all{An der Grenze gibt es lange Kontrollen, und manche Flüchtlinge müssen viele Tage warten.} « À la frontière, il y a de longs contrôles, et certains réfugiés doivent attendre plusieurs jours. » — \all{an der Grenze} : préposition \all{an} + datif (lieu).

\begin{attention}[Répondre à « warum » : le piège de « denn »]
En allemand, la réponse à une question en \all{warum} commence presque toujours par \all{weil}, avec le \textbf{verbe conjugué en fin de phrase} : \all{Weil sie dort keine Arbeit finden.} N'utilisez jamais \all{denn} seul en tête de réponse orale : \all{denn} est une conjonction de coordination qui ne renvoie pas le verbe à la fin et qui sonne mal en début de réponse spontanée. On réserve \all{denn} à la rédaction, à l'intérieur d'une phrase : \all{Sie bleiben nicht, denn sie finden keine Arbeit.} (verbe en 2\up{e} position après \all{denn}).
\end{attention}

\courssub{Vrai ou faux : justifier par une citation du texte}

\begin{exoresolu}[Vrai/Faux avec justification]
Énoncé. Répondez par \all{richtig} ou \all{falsch} et citez le passage du texte qui justifie votre réponse.\par
a) \all{Die jungen Leute verlassen ihre Heimat, weil sie dort keine Arbeit finden.}\par
b) \all{Alle Auswanderer bekommen sofort ein Visum.}\par
c) \all{Viele lernen Deutsch, damit sie sich im Ausland integrieren können.}
\tcblower
Solution détaillée.\par
a) \all{Richtig.} Justification : le texte dit que le chômage pousse les jeunes au départ ; on cite : \all{„Weil sie dort keine Arbeit finden“} (ou la phrase équivalente du texte). « Vrai : parce qu'ils n'y trouvent pas de travail. »\par
b) \all{Falsch.} Le texte montre au contraire que la demande de visa est longue et difficile ; on cite le passage qui parle des \all{„langen Wartezeiten bei der Botschaft“} (longs délais d'attente à l'ambassade). « Faux : tous les émigrants n'obtiennent pas immédiatement un visa. »\par
c) \all{Richtig.} On cite la phrase avec \all{damit} : \all{„Sie lernen die deutsche Sprache, damit sie sich integrieren können.“} « Vrai : ils apprennent l'allemand afin de pouvoir s'intégrer. »\par
Méthode : une réponse vrai/faux sans citation n'est jamais complète ; la citation entre guillemets allemands est la preuve exigée par le correcteur.
\end{exoresolu}

\courssub{Recherche dans le texte : expressions de cause et de but}

Le texte support contient des connecteurs de \cle{cause} (\all{weil}, \all{denn}, \all{da}) et de \cle{but} (\all{damit}, \all{um \ldots\ zu}). Le travail consiste à les repérer, à les traduire et à comparer leur construction avec le français.

\begin{center}
\begin{tabularx}{\linewidth}{|l|X|X|X|}
\hline
\rowcolor{popblueL} Connecteur & Sens en français & Place du verbe & Exemple du texte \\
\hline
\all{weil} & parce que & verbe en finale & \all{Weil sie keine Arbeit finden, \ldots} \\
\hline
\all{da} & puisque, comme & verbe en finale & \all{Da die Arbeitslosigkeit hoch ist, \ldots} \\
\hline
\all{denn} & car & verbe en 2\up{e} position & \all{\ldots, denn sie finden keine Arbeit.} \\
\hline
\all{damit} & afin que, pour que & verbe en finale & \all{\ldots, damit sie sich integrieren können.} \\
\hline
\all{um \ldots\ zu} & pour, afin de & infinitif en finale & \all{\ldots, um ein Visum zu beantragen.} \\
\hline
\end{tabularx}
\end{center}

\begin{exemplebox}[Traduction des expressions repérées]
\begin{itemize}
\item \all{Weil die Arbeitslosigkeit in der Heimat sehr hoch ist, wollen viele junge Leute auswandern.} « Parce que le chômage est très élevé dans leur pays, beaucoup de jeunes veulent émigrer. » — Quand la subordonnée est en tête, le verbe de la principale vient juste après la virgule (inversion : \all{wollen viele junge Leute}).
\item \all{Sie sparen Geld, denn die Reise ins Ausland ist teuer.} « Ils économisent de l'argent, car le voyage à l'étranger est coûteux. »
\item \all{Sie lernen Deutsch, um in Deutschland studieren zu können.} « Ils apprennent l'allemand pour pouvoir étudier en Allemagne. » (même sujet pour les deux verbes : \all{um \ldots\ zu}).
\item \all{Die Eltern helfen ihnen, damit sie das Visum bekommen.} « Les parents les aident afin qu'ils obtiennent le visa. » (sujets différents : \all{damit}).
\end{itemize}
\end{exemplebox}

\begin{attention}[« um … zu » ou « damit » ?]
Si le sujet des deux actions est \textbf{le même}, on emploie \all{um \ldots\ zu} + infinitif : \all{Er lernt Deutsch, um in Deutschland zu studieren.} Si les sujets sont \textbf{différents}, on emploie \all{damit} + verbe conjugué en finale : \all{Er lernt Deutsch, damit seine Familie stolz auf ihn ist.} Erreur fréquente des francophones : \all{*um sie sich integrieren zu können} est impossible, car le sujet change. Autre piège : avec un verbe à particule, \all{zu} s'insère entre la particule et le verbe : \all{um auszuwandern}, \all{um sich einzuleben}.
\end{attention}

\courssub{Résumé guidé du texte}

\begin{methode}[Rédiger un résumé en 3 à 4 phrases]
\begin{enumerate}
\item Phrase 1 : le thème général avec \all{Der Text handelt von \ldots} ou \all{Es geht um \ldots}.
\item Phrase 2 : la cause principale avec \all{weil} ou \all{da}.
\item Phrase 3 : les démarches ou les moyens avec \all{um \ldots\ zu} ou \all{damit}.
\item Phrase 4 : les difficultés ou la conclusion avec \all{aber} ou \all{deshalb}.
\end{enumerate}
Connecteurs imposés : \all{weil}, \all{um \ldots\ zu}, \all{aber}, \all{deshalb}.
\end{methode}

\begin{exemplebox}[Résumé modèle]
\all{Der Text handelt von der Auswanderung junger Leute. Weil sie in ihrer Heimat keine Arbeit finden und die Arbeitslosigkeit sehr hoch ist, wollen sie ins Ausland gehen. Sie beantragen ein Visum und lernen Deutsch, um sich dort integrieren zu können. Aber an der Grenze und bei der Botschaft gibt es viele Probleme, deshalb ist der Weg ins Ausland oft lang und schwierig.}\\
« Le texte traite de l'émigration des jeunes. Parce qu'ils ne trouvent pas de travail dans leur pays et que le chômage est très élevé, ils veulent partir à l'étranger. Ils demandent un visa et apprennent l'allemand pour pouvoir s'y intégrer. Mais à la frontière et à l'ambassade, il y a beaucoup de problèmes ; c'est pourquoi le chemin vers l'étranger est souvent long et difficile. »
\end{exemplebox}

\begin{exercice}[À vous : résumé personnel]
Rédigez votre propre résumé du texte en 3 à 4 phrases en utilisant obligatoirement les connecteurs \all{da}, \all{damit}, \all{aber} et \all{deshalb}. Soulignez les verbes conjugués et vérifiez leur place.
\end{exercice}

\begin{corrige}[Exercice : résumé personnel]
Modèle possible : \all{Da viele junge Leute in der Heimat arbeitslos sind, träumen sie vom Ausland. Sie lernen eine Fremdsprache, damit sie dort studieren oder arbeiten können. Aber das Visum ist schwer zu bekommen, deshalb warten manche viele Monate an der Grenze oder bei der Botschaft.} « Comme beaucoup de jeunes sont au chômage dans leur pays, ils rêvent de l'étranger. Ils apprennent une langue étrangère afin de pouvoir y étudier ou travailler. Mais le visa est difficile à obtenir ; c'est pourquoi certains attendent de nombreux mois à la frontière ou à l'ambassade. » Vérifications : après \all{da} et \all{damit}, verbe conjugué en finale ; après \all{deshalb}, inversion du sujet (\all{warten manche}) ; \all{vom Ausland} = \all{von dem Ausland} (datif après \all{von}).
\end{corrige}

\courssub{Discussion orale : et au Cameroun ?}

\begin{experience}[Débat guidé : comparer le texte avec la situation camerounaise]
Par groupes de quatre, discutez les questions suivantes en allemand, puis présentez vos conclusions à la classe.
\begin{enumerate}
\item \all{Gibt es in Kamerun auch viele junge Leute, die auswandern wollen? Warum?} « Y a-t-il aussi au Cameroun beaucoup de jeunes qui veulent émigrer ? Pourquoi ? »
\item \all{Wohin wollen die jungen Kameruner gehen — nach Deutschland, nach Frankreich oder in ein anderes Land?} « Où veulent aller les jeunes Camerounais — en Allemagne, en France ou dans un autre pays ? »
\item \all{Welche Probleme haben sie: das Visum, das Geld, die Sprache?} « Quels problèmes ont-ils : le visa, l'argent, la langue ? »
\item \all{Ist die Auswanderung eine gute Lösung? Was kann man in der Heimat tun?} « L'émigration est-elle une bonne solution ? Que peut-on faire dans son pays ? »
\end{enumerate}
Expressions utiles : \all{Ich meine, dass \ldots} (je pense que), \all{Meiner Meinung nach \ldots} (à mon avis), \all{Ich bin der Meinung, dass \ldots} (je suis d'avis que), \all{Einerseits \ldots, andererseits \ldots} (d'une part \ldots, d'autre part \ldots). Rappel : après \all{dass}, le verbe conjugué se place en finale.
\end{experience}

\begin{savaistu}[Le Cameroun et l'Allemagne : des liens anciens]
Le Cameroun a été un protectorat allemand de 1884 à 1916 : on disait \all{„Kamerun“}. C'est pourquoi l'allemand est aujourd'hui enseigné dans de nombreux lycées camerounais et pourquoi l'Allemagne est une destination prisée des étudiants camerounais. Beaucoup préparent les examens du Goethe-Institut à Yaoundé ou à Douala, \all{um in Deutschland studieren zu können} — « pour pouvoir étudier en Allemagne » : un bel exemple de but exprimé avec \all{um \ldots\ zu} !
\end{savaistu}

\begin{aretenir}[Ce qu'il faut retenir de cette section]
\begin{itemize}
\item On lit un texte en quatre étapes : lecture globale, lecture détaillée, résumé, discussion.
\item Une réponse de compréhension reprend les mots de la question et s'appuie sur une citation du texte entre guillemets allemands.
\item À \all{warum}, on répond par \all{weil} + verbe en finale ; \all{denn} ne se place pas en tête de réponse.
\item Cause : \all{weil}, \all{da} (verbe en finale), \all{denn} (verbe en 2\up{e} position) ; but : \all{damit} (sujets différents), \all{um \ldots\ zu} (même sujet).
\item \all{ins Ausland} (direction, accusatif) s'oppose à \all{im Ausland} (lieu, datif).
\item Un résumé se construit avec des connecteurs imposés et des verbes correctement placés.
\end{itemize}
\end{aretenir}

\coursec{Lexique thématique : l'émigration}

Après avoir lu et compris le texte « Warum so viele junge Leute auswandern », nous disposons d'un contexte riche : des jeunes qui quittent leur pays, des frontières, des visas, le chômage, l'espoir d'une vie meilleure. Il est temps maintenant de \cle{fixer le vocabulaire} de ce thème de manière systématique. Un texte du Bac sur l'émigration mobilise toujours les mêmes familles de mots ; les connaître avec leurs \cle{articles} et leurs \cle{pluriels}, c'est déjà gagner la moitié de la compréhension.

\courssub{Observer d'abord : un mini-texte de réactivation}

\begin{textebox}[Vokabeln im Kontext — « Aufbruch nach Europa »]
Viele junge Menschen verlassen ihre \textbf{Heimat}, weil sie in ihrem Land keine \textbf{Arbeit} finden. Die \textbf{Arbeitslosigkeit} und die \textbf{Armut} sind die wichtigsten Gründe für die \textbf{Auswanderung}. Wer nach Europa will, braucht ein \textbf{Visum}. Man muss einen \textbf{Antrag} stellen und oft lange warten. An der \textbf{Grenze} gibt es strenge \textbf{Kontrollen}, und ohne \textbf{Pass} ist die \textbf{Einreise} unmöglich. Manche Menschen wählen die \textbf{Flucht}: Sie \textbf{fliehen} vor Krieg oder Hunger und überqueren die Grenze illegal. Im \textbf{Ausland} suchen alle dasselbe: einen guten \textbf{Lohn}, Sicherheit und eine bessere \textbf{Zukunft}. Aber das Leben dort ist nicht einfach: Die \textbf{Integration} braucht Zeit, und viele \textbf{Ausländer} fühlen sich am Anfang fremd.
\end{textebox}

\textbf{Glossaire rapide :} der Aufbruch (le départ) ; streng (strict, sévère) ; illegal (illégal, adjectif) ; fremd (étranger, inconnu) ; brauchen (avoir besoin de).

Observez : presque tous les mots en gras appartiennent à quatre champs lexicaux. C'est exactement l'organisation que nous allons suivre.

\courssub{Champ 1 — Partir et fuir}

\begin{definition}[Les mots du départ]
\begin{itemize}
\item \all{die Auswanderung, die Auswanderungen} : l'émigration (le fait de quitter son pays pour s'installer ailleurs)
\item \all{auswandern} (wandert aus, ist ausgewandert) : émigrer — verbe faible (régulier) à particule séparable \all{aus-}, auxiliaire \all{sein} (déplacement)
\item \all{die Flucht, die Fluchten} : la fuite
\item \all{fliehen} (floh, ist geflohen) : fuir — verbe fort irrégulier, auxiliaire \all{sein}
\item \all{verlassen} (verließ, hat verlassen) : quitter, abandonner — verbe fort, préfixe \all{ver-} inséparable, transitif (auxiliaire \all{haben})
\item \all{die Heimat, die Heimaten} : la patrie, le pays natal (souvent employé au singulier seulement)
\end{itemize}
\end{definition}

\begin{exemplebox}[Exemples traduits]
\all{Viele Jugendliche sind letztes Jahr ausgewandert.} « Beaucoup de jeunes ont émigré l'année dernière. »

\all{Die Familie ist vor dem Krieg geflohen.} « La famille a fui devant la guerre. »

\all{Er hat seine Heimat schweren Herzens verlassen.} « Il a quitté sa patrie le cœur lourd. »
\end{exemplebox}

\begin{attention}[Constructions verbales : pièges à éviter]
\begin{itemize}
\item \all{verlassen} est \textbf{transitif direct} (accusatif) : \all{Er verlässt das Land.} « Il quitte le pays. » Interdit : \all{*Er verlässt aus dem Land}.
\item \all{fliehen} se construit avec \textbf{\all{vor} + datif} : \all{Sie fliehen vor der Armut.} « Ils fuient la pauvreté. »
\item \all{auswandern} est intransitif, auxiliaire \textbf{\all{sein}} au Perfekt : \all{Er ist ausgewandert.} (pas \all{*hat ausgewandert}).
\end{itemize}
\end{attention}

\courssub{Champ 2 — Frontières et papiers}

\begin{definition}[Les mots de l'administration]
\begin{itemize}
\item \all{das Visum, die Visa} : le visa (pluriel latin irrégulier, à retenir !)
\item \all{die Grenze, die Grenzen} : la frontière
\item \all{der Pass, die Pässe} : le passeport (Umlaut au pluriel)
\item \all{die Einreise, die Einreisen} : l'entrée (dans un pays)
\item \all{die Kontrolle, die Kontrollen} : le contrôle
\item \all{der Antrag, die Anträge} : la demande (officielle), le formulaire de demande (Umlaut au pluriel)
\item \all{beantragen} (beantragte, hat beantragt) : demander officiellement, solliciter — verbe faible
\end{itemize}
\end{definition}

\begin{exemplebox}[Exemples traduits]
\all{Er hat ein Visum beantragt, um in Deutschland zu studieren.} « Il a demandé un visa pour étudier en Allemagne. »

\all{Ohne Pass ist die Einreise nicht möglich.} « Sans passeport, l'entrée n'est pas possible. »

\all{An der Grenze kontrolliert die Polizei alle Papiere.} « À la frontière, la police contrôle tous les papiers. »
\end{exemplebox}

\courssub{Champ 3 — Travail et pauvreté}

\begin{definition}[Les mots économiques]
\begin{itemize}
\item \all{die Arbeitslosigkeit} (kein Plural) : le chômage (nom abstrait, toujours au singulier)
\item \all{der Arbeitslose, die Arbeitslosen} : le/la chômeur(se) — adjectif substantivé, déclinaison N (N-Deklination : \all{der Arbeitslose, des Arbeitslosen, dem Arbeitslosen, den Arbeitslosen} ; Pl. \all{die Arbeitslosen})
\item \all{die Arbeit, die Arbeiten} : le travail, l'ouvrage
\item \all{die Armut} (kein Plural) : la pauvreté (nom abstrait, toujours au singulier)
\item \all{der Lohn, die Löhne} : le salaire (Umlaut au pluriel)
\end{itemize}
\end{definition}

\begin{exemplebox}[Exemples traduits]
\all{Wegen der hohen Arbeitslosigkeit sind viele Jugendliche ausgewandert.} « À cause du chômage élevé, beaucoup de jeunes ont émigré. »

\all{Die Arbeitslosen finden in ihrer Heimat keine Arbeit.} « Les chômeurs ne trouvent pas de travail dans leur pays. »

\all{Der Lohn ist hier sehr niedrig.} « Le salaire est très bas ici. »
\end{exemplebox}

\begin{attention}[Singulier seul et faux ami]
\begin{itemize}
\item \all{die Armut} et \all{die Arbeitslosigkeit} sont des abstraits \textbf{invariables au singulier} (comme en français). On n'écrit jamais \all{*die Armuten} ni \all{*die Arbeitslosigkeiten}.
\item \textbf{Faux ami} : \all{die Arbeit} = « le travail » (l'activité, l'emploi). « L'armée » se dit \all{die Armee}.
\item \textbf{Género} : \all{das Visum} (neutre !), \all{der Pass} (masculin), \all{die Grenze} (féminin).
\end{itemize}
\end{attention}

\courssub{Champ 4 — Le pays d'accueil}

\begin{definition}[Les mots de l'arrivée]
\begin{itemize}
\item \all{das Ausland} (kein Plural) : l'étranger (ensemble des pays étrangers), toujours singulier
\item \all{im Ausland} (= \all{in dem Ausland}, datif) : à l'étranger (lieu)
\item \all{ins Ausland} (= \all{in das Ausland}, accusatif) : à l'étranger (direction)
\item \all{die Integration} (kein Plural) : l'intégration
\item \all{sich integrieren} (integrierte sich, hat sich integriert) : s'intégrer — verbe réfléchi, auxiliaire \all{haben}
\item \all{der Ausländer, die Ausländer} : l'étranger (personne) ; féminin \all{die Ausländerin, die Ausländerinnen}
\item \all{die Zukunft} (kein Plural) : l'avenir, le futur
\end{itemize}
\end{definition}

\begin{attention}[Das Ausland : piège prépositionnel]
\all{Das Ausland} est \textbf{toujours singulier}. On dit :
\begin{itemize}
\item \all{im Ausland leben / sein} (lieu, datif) — « vivre / être à l'étranger »
\item \all{ins Ausland gehen / fahren / reisen} (direction, accusatif) — « partir / aller à l'étranger »
\end{itemize}
Jamais \all{*in den Ausland}, \all{*aus den Ausland}. Pour « venir de l'étranger » : \all{aus dem Ausland kommen}.
\end{attention}

\begin{exemplebox}[Exemples traduits]
\all{Mein Onkel lebt seit zehn Jahren im Ausland.} « Mon oncle vit à l'étranger depuis dix ans. »

\all{Viele Ausländer integrieren sich schnell in Deutschland.} « Beaucoup d'étrangers s'intègrent vite en Allemagne. »

\all{Sie suchen im Ausland eine bessere Zukunft.} « Ils cherchent un avenir meilleur à l'étranger. »
\end{exemplebox}

\courssub{Tableau récapitulatif des quatre champs}

\begin{center}
\begin{tabularx}{\linewidth}{|l|l|l|X|}
\hline
\rowcolor{popblueL} Champ & Deutsch (Artikel + Plural) & Français & Beispiel \\
\hline
Partir & die Auswanderung, -en & l'émigration & Die Auswanderung nimmt zu. \\
\hline
Partir & fliehen (floh, ist geflohen) & fuir & Sie sind vor dem Krieg geflohen. \\
\hline
Partir & die Heimat, -en & la patrie & Er vermisst seine Heimat. \\
\hline
Papiers & das Visum, die Visa & le visa & Sie hat ein Visum bekommen. \\
\hline
Papiers & die Grenze, -n & la frontière & Die Grenze ist geschlossen. \\
\hline
Papiers & der Antrag, die Anträge & la demande & Er stellt einen Antrag. \\
\hline
Travail & die Arbeitslosigkeit (Sg.) & le chômage & Die Arbeitslosigkeit steigt. \\
\hline
Travail & der Lohn, die Löhne & le salaire & Der Lohn ist niedrig. \\
\hline
Accueil & das Ausland (Sg.) & l'étranger & Sie geht ins Ausland. \\
\hline
Accueil & die Integration (Sg.) & l'intégration & Integration braucht Zeit. \\
\hline
\end{tabularx}
\end{center}

\courssub{Les familles de mots : un mot en cache trois}

\begin{definition}[Qu'est-ce qu'une famille de mots ?]
Une \cle{famille de mots} regroupe des mots formés sur la \cle{même racine} : ils partagent un sens de base et se construisent par préfixes, suffixes ou changements de voyelle (Ablaut). Exemple : \all{fliehen} (fuir) $\rightarrow$ \all{die Flucht} (la fuite) $\rightarrow$ \all{der Flüchtling, -e} (le réfugié).
\end{definition}

\begin{propriete}[Famille 1 : wandern (marcher, randonner)]
\all{wandern} $\rightarrow$ \all{auswandern} (émigrer, particule séparable \all{aus-}) $\rightarrow$ \all{die Auswanderung} (l'émigration, suffixe \all{-ung} $\rightarrow$ féminin) $\rightarrow$ \all{der Auswanderer, -} (l'émigrant). \\
Inverse : \all{einwandern} (immigrer) $\rightarrow$ \all{die Einwanderung} $\rightarrow$ \all{der Einwanderer}.
\end{propriete}

\begin{propriete}[Famille 2 : arbeiten (travailler)]
\all{arbeiten} $\rightarrow$ \all{die Arbeit} (le travail) $\rightarrow$ \all{arbeitslos} (au chômage, adjectif, suffixe \all{-los} « sans ») $\rightarrow$ \all{die Arbeitslosigkeit} (le chômage, suffixe \all{-keit} $\rightarrow$ féminin) $\rightarrow$ \all{der Arbeitslose} (le chômeur, N-Deklination).
\end{propriete}

\begin{popfigure}
\begin{center}
\begin{tikzpicture}[
  root/.style={rectangle, rounded corners, draw=popblue, fill=popblue!20, font=\bfseries, align=center, minimum width=2.5cm},
  child node/.style={rectangle, rounded corners, draw=poporange, fill=poporange!20, align=center, font=\small, minimum width=2.2cm},
  edge from parent/.style={draw=popmuted, thick, ->}
]
\node[root, align=center]{arbeiten\\ (travailler)}
  child { node[child node, align=center]{die Arbeit\\ (le travail)} }
  child { node[child node, align=center]{arbeitslos\\ (sans travail)}
    child { node[child node, align=center]{die Arbeitslosigkeit\\ (le chômage)} }
    child { node[child node, align=center]{der Arbeitslose\\ (le chômeur)} }
  };
\end{tikzpicture}
\end{center}
\legende{Carte mentale de la famille de mots « arbeiten » : une racine, cinq mots clés.}
\end{popfigure}

\begin{savaistu}[Der Gastarbeiter : un mot chargé d'histoire]
\all{Der Gastarbeiter, die Gastarbeiter} (« travailleur invité ») désignait les migrants venus travailler en RFA dans les années 1960-70 (Italie, Turquie, Yougoslavie, Espagne, Grèce, Maroc, Tunisie, Portugal, Corée du Sud). Le terme est aujourd'hui historique — on préfère \all{der Migrant}, \all{der Einwanderer} ou \all{die Fachkraft} — mais il revient très souvent dans les textes du Bac sur l'histoire de l'immigration en Allemagne. Le connaître évite une méprise de compréhension !
\end{savaistu}

\courssub{Expressions utiles à réemployer (Redemittel)}

\begin{center}
\begin{tabularx}{\linewidth}{|X|X|}
\hline
\rowcolor{popgreenL} Deutsch & Français \\
\hline
ein Visum beantragen & demander un visa \\
\hline
ins Ausland gehen / reisen & partir à l'étranger \\
\hline
eine bessere Zukunft suchen & chercher un avenir meilleur \\
\hline
die Grenze überqueren & traverser la frontière \\
\hline
vor der Armut / dem Krieg fliehen & fuir la pauvreté / la guerre \\
\hline
seine Heimat verlassen & quitter sa patrie \\
\hline
sich (gut) integrieren & bien s'intégrer \\
\hline
einen Antrag stellen & déposer une demande \\
\hline
\end{tabularx}
\end{center}

\begin{methode}[Comment apprendre ce vocabulaire efficacement pour le Bac]
\begin{enumerate}
\item Apprends chaque nom \textbf{avec son article et son pluriel} sous forme explicite : \all{die Grenze, die Grenzen} ; \all{der Pass, die Pässe} ; \all{das Visum, die Visa}.
\item Réemploie \textbf{immédiatement} chaque mot dans une phrase personnelle à l'oral et à l'écrit : \all{Ich möchte nächstes Jahr ins Ausland gehen.}
\item Regroupe les mots par \textbf{familles morphologiques} : apprendre \all{auswandern} débloque \all{die Auswanderung} et \all{der Auswanderer}.
\item Vérifie la conjugaison des \textbf{verbes forts/irréguliers} : \all{fliehen — floh — ist geflohen} ; \all{verlassen — verließ — hat verlassen}.
\item Attention aux \textbf{prépositions régies} : \all{fliehen vor + Dat.}, \all{sich integrieren in + Akk.}, \all{wegen + Gen.}.
\end{enumerate}
\end{methode}

\begin{exoresolu}[Complétez avec le mot qui convient (forme correcte)]
Énoncé. Complétez les phrases avec : \all{Visum}, \all{Arbeitslosigkeit}, \all{Grenze}, \all{Heimat}, \all{Ausland}.

1. Wegen der hohen \dots\ finden viele Jugendliche keine Arbeit.
2. Er hat ein \dots\ für Deutschland beantragt.
3. An der \dots\ kontrolliert die Polizei die Pässe.
4. Viele Menschen verlassen ihre \dots, um Arbeit zu suchen.
5. Meine Schwester lebt seit fünf Jahren im \dots.
\tcblower
Solution détaillée.

1. \all{Arbeitslosigkeit} — Contexte « keine Arbeit finden » $\rightarrow$ cause = chômage. \all{wegen} + \textbf{génitif} $\rightarrow$ \all{wegen der hohen Arbeitslosigkeit}.

2. \all{Visum} — Verbe \all{beantragen} + accusatif $\rightarrow$ \all{ein Visum}. (Pluriel : \all{die Visa}).

3. \all{Grenze} — Lieu « à la frontière » $\rightarrow$ \all{an} + \textbf{datif} $\rightarrow$ \all{an der Grenze}.

4. \all{Heimat} — \all{verlassen} + accusatif $\rightarrow$ \all{ihre Heimat verlassen}.

5. \all{Ausland} — \all{im Ausland leben} (lieu, datif) ; toujours singulier, toujours \all{im}.
\end{exoresolu}

\begin{attention}[Les trois pièges classiques des francophones — Bilan]
\begin{enumerate}
\item \textbf{Genre} : \all{\textbf{das} Visum} (neutre, pas *der/*die) ; \all{\textbf{der} Pass} (masc.) ; \all{\textbf{die} Grenze} (fém., comme « frontière »).
\item \textbf{Pluriel irrégulier} : \all{die Visa} (latin), \all{die Pässe} (Umlaut), \all{die Anträge} (Umlaut), \all{die Löhne} (Umlaut). Pas de \all{*Visen}, \all{*Pässe} (sans Umlaut), etc.
\item \textbf{Faux ami \all{kontrollieren}} : signifie « vérifier, contrôler » (contrôle d'identité, de qualité), \textbf{pas} « surveiller avec suspicion » (c'est \all{überwachen}). \all{Die Kontrolle} = « le contrôle / la vérification ».
\end{enumerate}
\end{attention}

Ce vocabulaire est désormais votre boîte à outils. Dans la prochaine section, nous verrons comment relier ces mots entre eux grâce aux connecteurs de \cle{cause} (\all{weil}, \all{denn}, \all{da}) et de \cle{but} (\all{damit}, \all{um \dots\ zu}) — indispensables pour expliquer \emph{pourquoi} les gens émigrent et \emph{dans quel but}.

\coursec{Grammaire 1 : la cause et le but}

Dans la section précédente, nous avons constitué le lexique thématique de l'émigration : \all{die Auswanderung}, \all{die Flucht}, \all{das Visum}, \all{die Grenze}, \all{die Arbeitslosigkeit}, \all{das Ausland}. Ce vocabulaire ne prend tout son sens que lorsqu'on peut l'articuler dans un raisonnement : \emph{pourquoi} les jeunes quittent-ils leur pays ? \emph{dans quel but} apprennent-ils une langue étrangère ? Pour répondre à ces questions, l'allemand dispose de deux familles de connecteurs : les marqueurs de la \cle{cause} (\all{weil}, \all{da}, \all{denn}) et les marqueurs du \cle{but} (\all{um ... zu}, \all{damit}). C'est précisément là que le francophone se trompe le plus souvent, car la place du verbe change selon le connecteur choisi. Observons d'abord un court texte.

\begin{textebox}[Warum gehen sie? — Mini-texte d'observation]
Viele junge Menschen verlassen ihr Heimatland, \textbf{weil} sie dort keine Arbeit finden. \textbf{Da} die Arbeitslosigkeit in vielen Regionen hoch ist, sehen sie keine Zukunft mehr. Manche sparen jahrelang Geld, \textbf{denn} sie wollen die Reise und das Visum bezahlen. Sie lernen Deutsch oder Englisch, \textbf{um} im Ausland einen Beruf \textbf{zu} finden. Andere schicken Geld nach Hause, \textbf{damit} ihre Familien besser leben können. Doch die Auswanderung ist nicht einfach, \textbf{denn} die Grenzen sind streng kontrolliert, und viele warten lange auf ein Visum, \textbf{weil} die Botschaften sehr viele Anträge bekommen.
\par\smallskip
\emph{Glossaire :} das Heimatland (-„er) : le pays natal ; die Zukunft : l'avenir ; der Antrag (-„e) : la demande (administrative) ; die Botschaft (-en) : l'ambassade ; streng : sévèrement.
\end{textebox}

\textbf{Questions d'observation.} Relevez dans le texte les cinq connecteurs soulignés. Pour chacun : (1) où se trouve le verbe conjugué ? (2) la proposition peut-elle commencer la phrase ? (3) les deux parties de la phrase ont-elles le même sujet ? Gardez vos réponses : elles annoncent les règles qui suivent.

\courssub{La cause avec \all{weil}}

\begin{definition}[La subordonnée de cause avec \all{weil}]
\all{Weil} signifie « parce que ». C'est une \cle{conjonction de subordination} : elle introduit une subordonnée dans laquelle le \cle{verbe conjugué va à la toute fin}. C'est la réponse typique à la question \all{Warum?} (Pourquoi ?).
\end{definition}

\begin{propriete}[Après \all{weil}, le verbe conjugué va à la fin]
\all{Er lernt Deutsch, weil er in Deutschland studieren \textbf{will}.} (Il apprend l'allemand parce qu'il veut étudier en Allemagne.)
\begin{itemize}
\item Structure : Hauptsatz (verbe en position 2) + \all{weil} + sujet + ... + verbe conjugué \emph{en dernière position}.
\item S'il y a un infinitif ou un participe, le verbe conjugué se place \emph{après} lui : \all{..., weil er ein Visum beantragen \textbf{muss}.} (..., parce qu'il doit demander un visa.)
\item La subordonnée peut aussi ouvrir la phrase ; la principale commence alors par son verbe : \all{Weil er keine Arbeit findet, wandert er aus.} (Parce qu'il ne trouve pas de travail, il émigre.)
\end{itemize}
\end{propriete}

\begin{exemplebox}[Exemples avec \all{weil}]
\all{Viele junge Leute wandern aus, weil die Arbeitslosigkeit hoch ist.} — Beaucoup de jeunes émigrent parce que le chômage est élevé.\\
\all{Sie flieht, weil sie in ihrer Heimat nicht sicher ist.} — Elle fuit parce qu'elle n'est pas en sécurité dans son pays.\\
\all{Wir lernen Deutsch, weil wir später in Österreich arbeiten wollen.} — Nous apprenons l'allemand parce que nous voulons travailler plus tard en Autriche.\\
\all{Er hat das Visum nicht bekommen, weil seine Papiere unvollständig waren.} — Il n'a pas obtenu le visa parce que ses documents étaient incomplets.
\end{exemplebox}

\courssub{La cause avec \all{da}}

\begin{propriete}[\all{Da} : la cause déjà connue, souvent en tête de phrase]
\all{Da} signifie « comme, puisque, étant donné que ». Il fonctionne exactement comme \all{weil} (subordonnée, verbe conjugué à la fin), mais avec deux nuances d'emploi :
\begin{enumerate}
\item Il présente une cause \emph{déjà connue} de l'interlocuteur ou évidente dans le contexte ;
\item Il se place \emph{très souvent en tête de phrase} : la subordonnée ouvre l'énoncé, et la principale suit avec le verbe en position 2.
\end{enumerate}
\all{Da die Arbeitslosigkeit hoch ist, wandern viele aus.} (Comme le chômage est élevé, beaucoup émigrent.)
\end{propriete}

\begin{attention}[\all{Da} en tête : attention à la place du verbe de la principale]
Quand la subordonnée en \all{da} (ou \all{weil}) ouvre la phrase, elle occupe la \emph{première position} ; le verbe de la principale vient \emph{immédiatement après la virgule}, puis le sujet : \all{Da er kein Visum hat, \textbf{bleibt} \emph{er} im Land.} et non \all{*Da er kein Visum hat, er bleibt...}. La subordonnée entière compte comme un seul élément de phrase.
\end{attention}

\begin{exemplebox}[Exemples avec \all{da}]
\all{Da sie die Grenze nicht überqueren konnte, blieb sie in Douala.} — Comme elle ne pouvait pas franchir la frontière, elle est restée à Douala.\\
\all{Da das Visum teuer ist, sparen viele Familien lange.} — Comme le visa coûte cher, beaucoup de familles économisent longtemps.
\end{exemplebox}

\courssub{La cause avec \all{denn}}

\begin{definition}[\all{Denn} : une coordination, pas une subordination]
\all{Denn} signifie « car ». Contrairement à \all{weil} et \all{da}, c'est une \cle{conjonction de coordination} : la proposition qui suit est une phrase normale, avec le \cle{verbe conjugué en position 2}. \all{Denn} ne peut \emph{jamais} ouvrir la phrase.
\end{definition}

\begin{propriete}[Après \all{denn}, le verbe reste en position 2]
\all{Er lernt Deutsch, denn er \textbf{will} in Deutschland studieren.} (Il apprend l'allemand car il veut étudier en Allemagne.)
\begin{itemize}
\item Structure : Hauptsatz + \all{denn} + sujet + verbe conjugué (position 2) + reste.
\item Registre : \all{denn} est un peu plus soutenu et écrit que \all{weil} ; à l'oral, on préfère souvent \all{weil}.
\item Place : toujours \emph{après} la proposition principale, jamais en tête.
\end{itemize}
\end{propriete}

\begin{exemplebox}[Exemples avec \all{denn}]
\all{Sie bleibt im Ausland, denn dort verdient sie mehr Geld.} — Elle reste à l'étranger, car là-bas elle gagne plus d'argent.\\
\all{Die Reise ist gefährlich, denn die Grenzen sind streng kontrolliert.} — Le voyage est dangereux, car les frontières sont strictement contrôlées.\\
\all{Er schreibt seiner Familie oft, denn er vermisst sie sehr.} — Il écrit souvent à sa famille, car elle lui manque beaucoup.
\end{exemplebox}

\begin{attention}[L'erreur n° 1 du francophone : la place du verbe]
Ne dites jamais \all{*weil er will studieren} ni \all{*denn er studieren will}. Retenez le réflexe :
\begin{center}
\all{weil} / \all{da} = verbe conjugué \textbf{à la fin} ; \quad \all{denn} = verbe conjugué \textbf{en 2\up{e} position}.
\end{center}
Autre piège : \all{denn} ne signifie jamais « alors » au sens temporel ; et ne confondez pas \all{denn} (car) avec \all{dann} (ensuite) : \all{Zuerst beantragt er das Visum, dann reist er ab.} (D'abord il demande le visa, ensuite il part.)
\end{attention}

\begin{popfigure}
\begin{tikzpicture}[font=\small, node distance=0.45cm]
\node[draw=popblue, fill=popblueL, rounded corners=2pt, align=center, minimum width=5.2cm] (hs) {Hauptsatz : \all{Er lernt Deutsch,}\\ verbe en position 2};
\node[draw=popgreen, fill=popgreenL, rounded corners=2pt, align=center, minimum width=6.8cm, below=of hs] (weil) {\all{weil} + sujet + ... + \textbf{verbe à la fin}\\ \all{weil er studieren \textbf{will}}};
\node[draw=poporange, fill=poporangeL, rounded corners=2pt, align=center, minimum width=6.8cm, below=of weil] (denn) {\all{denn} + sujet + \textbf{verbe en position 2}\\ \all{denn er \textbf{will} studieren}};
\draw[-{Stealth}, thick, popgreen] (hs.south) -- (weil.north);
\draw[-{Stealth}, thick, poporange] (hs.south) -- (denn.north);
\node[draw=poppurple, fill=poppurpleL, rounded corners=2pt, align=center, right=1.1cm of weil] (memo) {Subordination :\\ le verbe « saute »\\ à la fin};
\draw[-{Stealth}, dashed, poppurple] (memo.west) -- (weil.east);
\end{tikzpicture}
\legende{Position du verbe conjugué : subordonnée avec \all{weil} (verbe final) contre coordination avec \all{denn} (verbe en position 2).}
\end{popfigure}

\courssub{Le but avec \all{um ... zu} + infinitif}

\begin{propriete}[\all{Um ... zu} : « pour + infinitif », même sujet obligatoire]
Pour exprimer le but (« pour, afin de »), l'allemand utilise \all{um} en tête et \all{zu} + infinitif à la fin. Condition absolue : \cle{le sujet des deux actions doit être le même}.
\par\smallskip
\all{Sie lernt Deutsch, um im Ausland \textbf{zu arbeiten}.} (Elle apprend l'allemand pour travailler à l'étranger.)
\begin{itemize}
\item Structure : Hauptsatz + \all{um} + compléments + \all{zu} + infinitif (en dernier).
\item Avec un verbe à particule séparable, \all{zu} s'insère dans le verbe : \all{um ein Visum \textbf{beanzutragen}} (pour demander un visa).
\item La construction \all{um ... zu} peut aussi ouvrir la phrase : \all{Um Geld zu verdienen, zieht er nach Deutschland.} (Pour gagner de l'argent, il part en Allemagne.)
\end{itemize}
\end{propriete}

\begin{exemplebox}[Exemples avec \all{um ... zu}]
\all{Viele beantragen ein Visum, um die Grenze legal zu überqueren.} — Beaucoup demandent un visa pour franchir la frontière légalement.\\
\all{Er spart Geld, um die Reise zu bezahlen.} — Il économise de l'argent pour payer le voyage.\\
\all{Wir wiederholen die Grammatik, um die Prüfung zu bestehen.} — Nous révisons la grammaire pour réussir l'examen.
\end{exemplebox}

\courssub{Le but avec \all{damit}}

\begin{propriete}[\all{Damit} : « afin que », sujets différents ou style soutenu]
\all{Damit} introduit une \cle{subordonnée} (verbe conjugué à la fin, comme avec \all{weil}). On l'emploie :
\begin{enumerate}
\item \emph{obligatoirement} quand les deux parties de la phrase ont des \emph{sujets différents} ;
\item \emph{au choix} quand le sujet est identique, dans un style plus soutenu.
\end{enumerate}
\all{Sie lernt Deutsch, damit ihre Kinder später studieren \textbf{können}.} (Elle apprend l'allemand afin que ses enfants puissent étudier plus tard.)
\end{propriete}

\begin{exemplebox}[Exemples avec \all{damit}]
\all{Viele beantragen ein Visum, damit sie die Grenze überqueren können.} — Beaucoup demandent un visa afin de pouvoir franchir la frontière.\\
\all{Er schickt Geld nach Hause, damit seine Eltern ein Haus bauen können.} — Il envoie de l'argent au pays pour que ses parents puissent construire une maison.\\
\all{Die Regierung bildet junge Leute aus, damit sie im Land Arbeit finden.} — Le gouvernement forme les jeunes afin qu'ils trouvent du travail dans le pays.
\end{exemplebox}

\begin{attention}[Le choix \all{um ... zu} ou \all{damit} : regardez le sujet]
Même sujet $\rightarrow$ \all{um ... zu} : \all{Ich lerne Deutsch, um in Deutschland zu studieren.} (J'apprends l'allemand pour étudier en Allemagne.)\\
Sujets différents $\rightarrow$ \all{damit} : \all{Ich lerne Deutsch, damit mein Bruder in Deutschland studieren kann.} (J'apprends l'allemand pour que mon frère puisse étudier en Allemagne.)\\
Interdit : \all{*um mein Bruder zu studieren} — \all{um ... zu} ne contient jamais de sujet propre ni de verbe conjugué.
\end{attention}

\courssub{Tableau comparatif et transformations}

\begin{center}
\begin{tabularx}{\linewidth}{|l|X|X|X|}
\hline
\rowcolor{popblueL} \textbf{Connecteur} & \textbf{Français} & \textbf{Structure / place du verbe} & \textbf{Emploi} \\
\hline
\all{weil} & parce que & subordonnée, verbe à la fin & réponse à \all{Warum?}, oral et écrit \\
\hline
\all{da} & comme, puisque & subordonnée, verbe à la fin & cause connue, souvent en tête de phrase \\
\hline
\all{denn} & car & coordination, verbe en position 2 & jamais en tête, style plutôt écrit \\
\hline
\all{um ... zu} & pour + infinitif & \all{zu} + infinitif à la fin & même sujet dans les deux parties \\
\hline
\all{damit} & afin que, pour que & subordonnée, verbe à la fin & sujets différents ou style soutenu \\
\hline
\end{tabularx}
\end{center}

\begin{methode}[Transformer une phrase : \all{weil} $\leftrightarrow$ \all{denn}, \all{damit} $\leftrightarrow$ \all{um ... zu}]
\begin{enumerate}
\item \textbf{Identifier} le type de proposition : le verbe est-il à la fin (subordonnée) ou en position 2 (coordination) ?
\item \textbf{Changer le connecteur} : \all{weil} $\rightarrow$ \all{denn} ou l'inverse.
\item \textbf{Déplacer le verbe conjugué} en conséquence : fin $\leftrightarrow$ position 2. C'est l'étape que les francophones oublient !
\item Pour \all{damit} $\rightarrow$ \all{um ... zu} : vérifier que les sujets sont identiques, supprimer le sujet de la subordonnée, mettre le verbe à l'infinitif avec \all{zu}. Si les sujets diffèrent, la transformation est \emph{impossible}.
\end{enumerate}
\end{methode}

\begin{exoresolu}[Transformations guidées]
\textbf{Énoncé.} a) Transformez avec \all{denn} : \all{Er bleibt im Ausland, weil er dort Arbeit gefunden hat.} — b) Transformez avec \all{weil} : \all{Sie lernt Deutsch, denn sie will in der Schweiz studieren.} — c) Transformez avec \all{um ... zu} : \all{Er spart Geld, damit er ein Visum beantragen kann.}
\tcblower
\textbf{Solution détaillée.}
\begin{itemize}
\item a) \all{weil} est une subordination (verbe \all{hat} à la fin) ; avec \all{denn}, le verbe passe en position 2 : \all{Er bleibt im Ausland, denn er \textbf{hat} dort Arbeit gefunden.} (Il reste à l'étranger, car il y a trouvé du travail.)
\item b) \all{denn} est une coordination (verbe \all{will} en position 2) ; avec \all{weil}, le verbe va à la fin : \all{Sie lernt Deutsch, weil sie in der Schweiz studieren \textbf{will}.} (Elle apprend l'allemand parce qu'elle veut étudier en Suisse.)
\item c) Les sujets sont identiques (\all{er} = \all{er}) : transformation possible. On supprime le sujet et le modal, et on met l'infinitif avec \all{zu} (inséré dans le verbe séparable \all{beantragen}) : \all{Er spart Geld, um ein Visum \textbf{beanzutragen}.} (Il économise de l'argent pour demander un visa.)
\end{itemize}
\end{exoresolu}

\begin{exercice}[À vous de jouer]
\begin{enumerate}
\item Complétez avec \all{weil}, \all{da} ou \all{denn} : \all{... die Arbeitslosigkeit hoch ist, wandern viele junge Leute aus.} / \all{Er wartet lange, ... die Botschaft viele Anträge bekommt.} / \all{Sie flieht, ... sie in ihrer Heimat nicht sicher ist.}
\item Transformez avec \all{um ... zu} ou justifiez l'impossibilité : a) \all{Ich lerne Deutsch, damit ich im Ausland arbeiten kann.} b) \all{Er arbeitet viel, damit seine Kinder zur Schule gehen können.}
\item Traduisez : « Comme beaucoup de jeunes ne trouvent pas de travail, ils apprennent l'allemand pour émigrer. »
\end{enumerate}
\end{exercice}

\begin{corrige}[Exercice « À vous de jouer »]
\begin{enumerate}
\item \all{\textbf{Da} die Arbeitslosigkeit hoch ist, wandern viele junge Leute aus.} (cause connue, en tête) ; \all{Er wartet lange, \textbf{weil} die Botschaft viele Anträge bekommt.} (verbe \all{bekommt} à la fin) ; \all{Sie flieht, \textbf{denn} sie ist in ihrer Heimat nicht sicher.} (verbe \all{ist} en position 2).
\item a) Même sujet : \all{Ich lerne Deutsch, um im Ausland zu arbeiten.} b) Sujets différents (\all{er} / \all{seine Kinder}) : transformation impossible, on garde \all{damit}.
\item \all{Da viele junge Leute keine Arbeit finden, lernen sie Deutsch, um auszuwandern.} (Comme beaucoup de jeunes ne trouvent pas de travail, ils apprennent l'allemand pour émigrer.)
\end{enumerate}
\end{corrige}

\begin{aretenir}[L'essentiel de la section]
\begin{itemize}
\item \all{weil} et \all{da} : subordination $\rightarrow$ verbe conjugué \textbf{à la fin} ; \all{da} aime la tête de phrase.
\item \all{denn} : coordination $\rightarrow$ verbe \textbf{en position 2}, jamais en tête de phrase.
\item \all{um ... zu} : même sujet, infinitif avec \all{zu} à la fin ; \all{damit} : sujets différents, verbe conjugué à la fin.
\item Réflexe de survie : après avoir choisi le connecteur, vérifiez \emph{toujours} la place du verbe.
\end{itemize}
\end{aretenir}

\coursec{Grammaire 2 : parfait, prétérit et subjonctif II}

Dans la section précédente, nous avons appris à exprimer la \cle{cause} (\all{weil}, \all{denn}, \all{da}) et le \cle{but} (\all{damit}, \all{um ... zu}) : \all{Viele junge Leute wandern aus, weil sie keine Arbeit finden, um ihren Familien zu helfen.} Mais pour raconter l'émigration — ce qui s'est passé hier, l'année dernière — et pour exprimer les rêves et les regrets des candidats au départ, il nous faut maîtriser trois outils temporels et modaux : le \cle{parfait} (Perfekt), le \cle{prétérit} (Präteritum) et le \cle{subjonctif II} (Konjunktiv II). C'est l'objet de cette section.

\courssub{Observation : les temps du passé dans un texte de presse}

Lisons d'abord ce court article original, rédigé dans le style de la presse germanophone, et observons les formes verbales soulignées.

\begin{textebox}[Aus der Presse : „Die neue Auswanderungswelle“]
Immer mehr junge Fachkräfte verließen im vergangenen Jahr ihr Heimatland, um im Ausland Arbeit zu suchen. Ein junger Ingenieur aus Yaoundé erklärte unserer Zeitung: „Ich habe drei Jahre lang nach einer Stelle gesucht, aber ich habe nichts gefunden. Deshalb bin ich nach Deutschland ausgewandert.“ Er hatte Glück: Er bekam ein Visum und fand schnell eine Arbeit in München. Andere waren weniger erfolgreich. Viele mussten monatelang auf eine Antwort der Botschaft warten, und einige gaben ihre Pläne schließlich auf. „Wenn ich die Situation früher gekannt hätte, wäre ich vielleicht nicht gegangen“, sagte eine junge Krankenschwester aus Douala. „Aber jetzt bin ich hier, und ich möchte bleiben.“ Experten warnen: Wenn die Arbeitslosigkeit weiter steigt, werden noch mehr junge Menschen fliehen. Die Regierung versprach gestern neue Programme für die Jugend. Ob sie helfen werden, ist noch unklar. Viele junge Leute sagen: „Wir würden gern bleiben, wenn wir hier eine Zukunft hätten.“

\textbf{Glossaire :} die Fachkraft, die Fachkräfte (personnel qualifié) ; die Stelle, die Stellen (le poste de travail) ; das Visum, die Visa (le visa) ; die Botschaft, die Botschaften (l'ambassade) ; aufgeben, gab auf, hat aufgegeben (renoncer) ; die Krankenschwester, die Krankenschwestern (l'infirmière) ; warnen (avertir) ; versprechen, versprach, hat versprochen (promettre) ; die Zukunft (l'avenir).
\end{textebox}

\textbf{Questions d'observation :}
\begin{enumerate}
\item Relevez dans le texte trois verbes au \cle{prétérit} et trois verbes au \cle{parfait}.
\item Quels verbes au parfait utilisent l'auxiliaire \all{sein} ? Pourquoi ?
\item Relevez deux phrases au subjonctif II. Qu'expriment-elles ?
\end{enumerate}

On constate que le texte mélange deux temps du passé : \all{verließen}, \all{hatte}, \all{bekam}, \all{mussten}, \all{waren} (prétérit) et \all{habe gesucht}, \all{habe gefunden}, \all{bin ausgewandert} (parfait). Et deux phrases expriment l'irréel : \all{wäre ich nicht gegangen}, \all{würden gern bleiben, wenn wir ... hätten}. Expliquons tout cela.

\courssub{Le parfait (Perfekt) : rappel}

\begin{propriete}[Formation du parfait]
Le parfait se forme avec l'auxiliaire \all{haben} ou \all{sein} au présent + le \cle{participe II} placé en fin de phrase :
\begin{itemize}
\item Verbes réguliers (faibles) : \all{ge- + radical + -t} ; \all{beantragen → beantragt} (pas de \all{ge-} car préfixe inséparable \all{be-}) ;
\item Verbes forts (irréguliers) : \all{ge- + radical modifié + -en} ; \all{fliehen → geflohen} ;
\item Verbes forts à préfixe inséparable (\all{ver-}, \all{be-}, \all{ent-}, \all{er-}, \all{ge-}, \all{miss-}, \all{zer-}) : pas de \all{ge-} ; \all{verlassen → verlassen} ;
\item Verbes à particule séparable : \all{ge-} s'insère entre la particule et le radical ; \all{auswandern → ausgewandert}.
\end{itemize}
\end{propriete}

\begin{propriete}[Le choix de l'auxiliaire : haben ou sein]
On utilise \all{sein} quand le verbe exprime un \cle{déplacement} (changement de lieu) ou un \cle{changement d'état}, ainsi que pour \all{sein} et \all{bleiben}. Tous les autres verbes prennent \all{haben}. Les verbes clés de notre leçon :
\begin{itemize}
\item Avec \all{sein} : \all{auswandern} (émigrer), \all{fliehen} (fuir), \all{gehen}, \all{kommen}, \all{reisen} ;
\item Avec \all{haben} : \all{beantragen} (demander officiellement), \all{suchen}, \all{finden}, \all{verlassen} (quitter qqch/qqn : transitif).
\end{itemize}
\end{propriete}

\begin{exemplebox}[Le parfait dans le discours direct]
\all{Ich habe drei Jahre lang nach einer Stelle gesucht.} « J'ai cherché un poste pendant trois ans. »

\all{Deshalb bin ich nach Deutschland ausgewandert.} « C'est pourquoi j'ai émigré en Allemagne. »

\all{Viele sind geflohen, weil sie keine Zukunft gesehen haben.} « Beaucoup ont fui parce qu'ils n'ont pas vu d'avenir. »
\end{exemplebox}

\begin{attention}[verlassen : transitif, donc haben]
Les francophones hésitent souvent : \all{verlassen} (quitter) est \textbf{transitif} — on quitte quelque chose — donc il prend \all{haben} : \all{Er hat sein Land verlassen.} « Il a quitté son pays. » En revanche \all{weggehen} (partir, s'en aller), intransitif et déplacement, prend \all{sein} : \all{Er ist weggegangen.}
\end{attention}

\courssub{Le prétérit (Präteritum) : le temps du récit écrit}

\begin{propriete}[Emploi du prétérit]
Le prétérit est le temps du \cle{récit écrit} : presse, romans, rapports, biographies. À l'oral, on préfère le parfait. \textbf{Exception majeure :} \all{sein} (\all{war}), \all{haben} (\all{hatte}) et les verbes modaux (\all{konnte}, \all{musste}, \all{wollte}) s'emploient au prétérit \textbf{même à l'oral} — le parfait (\all{ich bin gewesen}) sonne lourd et artificiel.
\end{propriete}

\begin{exemplebox}[Prétérit à l'écrit et à l'oral]
\all{Im vergangenen Jahr verließen viele Fachkräfte ihr Heimatland.} « L'année dernière, de nombreux cadres qualifiés quittèrent leur pays. » (presse)

\all{Er hatte Glück und bekam ein Visum.} « Il a eu de la chance et a obtenu un visa. »

\all{Ich war gestern bei der Botschaft.} « J'étais hier à l'ambassade. » (oral : prétérit de \all{sein}, tout à fait normal)

\all{Ich hatte keine Zeit.} « Je n'avais pas le temps. » (oral : prétérit de \all{haben}, normal)
\end{exemplebox}

\begin{attention}[Parfait ou prétérit ?]
Ne dites pas à l'oral : \all{Ich bin gestern bei der Botschaft gewesen.} Dites : \all{Ich war gestern bei der Botschaft.} Règle pratique : à l'oral, parfait pour tous les verbes \textbf{sauf} \all{sein}, \all{haben} et les modaux ; à l'écrit (récit, article), prétérit partout.
\end{attention}

\begin{center}
\begin{tabularx}{\linewidth}{|l|X|X|X|}
\hline
\rowcolor{popblueL} \textbf{Infinitif} & \textbf{Prétérit (3\textsuperscript{e} pers. sg.)} & \textbf{Parfait} & \textbf{Français} \\
\hline
auswandern & wanderte aus & ist ausgewandert & émigrer \\
\hline
fliehen & floh & ist geflohen & fuir \\
\hline
verlassen & verließ & hat verlassen & quitter \\
\hline
beantragen & beantragte & hat beantragt & demander (officiellement) \\
\hline
sein & war & ist gewesen & être \\
\hline
haben & hatte & hat gehabt & avoir \\
\hline
\end{tabularx}
\end{center}

\begin{popfigure}
\begin{tikzpicture}[scale=1]
\draw[-{Stealth[length=3mm]}, thick, popdark] (0,0) -- (12.5,0);
\foreach \x/\t in {1.5/Vergangenheit, 6.2/Gegenwart, 10.8/Irrealis} {
  \fill[poporange] (\x,0) circle (0.09);
  \node[above, align=center, font=\small\bfseries, popdark] at (\x,0.15) {\t};
}
\node[below, align=center, font=\small, text=popmuted] at (1.5,-0.25) {Präteritum (écrit)\\Perfekt (oral)};
\node[below, align=center, font=\small, text=popmuted] at (6.2,-0.25) {Präsens\\ich wandere aus};
\node[below, align=center, font=\small, text=popmuted] at (10.8,-0.25) {Konjunktiv II\\souhait, hypothèse};
\draw[-{Stealth[length=2.5mm]}, thick, poppurple] (1.5,0.9) to[bend left=25] node[above, font=\small, text=poppurple, align=center]{Base du\\Präteritum + Umlaut} (10.2,0.9);
\node[font=\small, text=poppurple, align=center] at (6.0,-1.6) {war $\rightarrow$ wäre \quad hatte $\rightarrow$ hätte \quad floh $\rightarrow$ flöhe};
\end{tikzpicture}
\legende{Frise des temps : du passé (prétérit/parfait) au présent, puis au souhait irréel (subjonctif II), formé à partir du radical du prétérit.}
\end{popfigure}

\courssub{Le subjonctif II (Konjunktiv II) : le monde de l'irréel}

Revenons à la phrase de la jeune infirmière : \all{Wenn ich die Situation früher gekannt hätte, wäre ich vielleicht nicht gegangen.} « Si j'avais connu la situation plus tôt, je ne serais peut-être pas partie. » Elle parle d'une réalité qui n'a \textbf{pas} eu lieu : c'est le domaine du subjonctif II, l'équivalent de notre conditionnel français.

\begin{propriete}[Formation du subjonctif II]
Le subjonctif II se forme sur le \cle{radical du prétérit}, avec un \cle{tréma} (Umlaut) sur les voyelles \textbf{a, o, u} pour les verbes forts et les auxiliaires :
\begin{itemize}
\item \all{sein} : war → \all{wäre} (ich wäre, du wärst/wärest, er wäre, wir wären, ihr wärt/wäret, sie wären) ;
\item \all{haben} : hatte → \all{hätte} (ich hätte, du hättest, er hätte, wir hätten, ihr hättet, sie hätten) ;
\item Modaux : konnte → \all{könnte}, musste → \all{müsste}, durfte → \all{dürfte}, wollte → \all{wollte} (pas de tréma !), sollte → \all{sollte} ;
\item Verbes forts : kam → \all{käme}, gab → \all{gäbe}, wüsste (de \all{wissen}).
\end{itemize}
\end{propriete}

\begin{popfigure}
\begin{center}
\begin{tabularx}{\linewidth}{|X|X|X|X|}
\hline
\rowcolor{popblueL} \textbf{Infinitiv} & \textbf{Präteritum} & \textbf{Konjunktiv II} & \textbf{Français} \\
\hline
sein & war & wäre & je serais \\
\hline
haben & hatte & hätte & j'aurais \\
\hline
können & konnte & könnte & je pourrais \\
\hline
müssen & musste & müsste & je devrais (obligation) \\
\hline
werden & wurde & würde & (conditionnel général) \\
\hline
\end{tabularx}
\end{center}
\legende{Les cinq formes de subjonctif II indispensables : du prétérit au subjonctif II par ajout du tréma (sauf wollen/sollen).}
\end{popfigure}

\begin{propriete}[würde + infinitif : le cas général]
Pour la plupart des verbes, on évite les formes simples du subjonctif II (souvent identiques au prétérit ou trop lourdes) et on utilise \all{würde} + \cle{infinitif} en fin de phrase :
\all{ich würde bleiben} « je resterais », \all{wir würden helfen} « nous aiderions ».
Les formes simples qu'il faut absolument connaître : \all{wäre}, \all{hätte}, \all{könnte}, \all{müsste}, \all{sollte}, \all{gäbe} (\all{es gäbe} « il y aurait »).
\end{propriete}

\begin{exemplebox}[Les trois grands emplois du subjonctif II]
\textbf{1. Le souhait irréel} (avec \all{doch}, \all{nur}, \all{bloß}) :

\all{Wenn ich doch ein Visum hätte!} « Si seulement j'avais un visa ! »

\all{Hätte ich nur mehr Geld!} « Si seulement j'avais plus d'argent ! » (inversion du verbe)

\textbf{2. L'hypothèse irréelle} (si + subjonctif II) :

\all{Wenn ich reich wäre, würde ich nicht auswandern.} « Si j'étais riche, je n'émigrerais pas. »

\all{Wenn ich ein Visum hätte, würde ich in Deutschland studieren.} « Si j'avais un visa, j'étudierais en Allemagne. »

\textbf{3. Le conseil poli} :

\all{Du solltest zuerst Deutsch lernen.} « Tu devrais d'abord apprendre l'allemand. »

\all{An deiner Stelle würde ich bei der Botschaft nachfragen.} « À ta place, je me renseignerais à l'ambassade. »
\end{exemplebox}

\begin{exemplebox}[Exemple tiré de la leçon]
\all{Viele sind ausgewandert, weil sie keine Arbeit hatten.} « Beaucoup ont émigré parce qu'ils n'avaient pas de travail. » — On remarque ici la combinaison parfait (\all{sind ausgewandert}, verbe de mouvement) + prétérit de \all{haben} (\all{hatten}) dans la subordonnée causale avec \all{weil} : exactement la grammaire de notre chapitre.
\end{exemplebox}

\begin{center}
\begin{tabularx}{\linewidth}{|X|X|X|}
\hline
\rowcolor{popblueL} \textbf{Français} & \textbf{Deutsch} & \textbf{Remarque} \\
\hline
Si j'avais un visa... & Wenn ich ein Visum hätte... & conditionnel français = Konj. II \\
\hline
j'étudierais en Allemagne & würde ich in Deutschland studieren & würde + infinitif en fin \\
\hline
je voudrais un visa & ich hätte gern ein Visum / ich möchte ein Visum & politesse, pas « ich will » \\
\hline
tu devrais partir & du solltest auswandern & conseil \\
\hline
si seulement j'étais riche ! & Wenn ich doch reich wäre! & souhait irréel \\
\hline
\end{tabularx}
\end{center}

\begin{attention}[Le piège du souhait irréel]
Ne dites pas \all{Ich will ein Visum haben} pour exprimer un souhait : cela signifie « je veux avoir un visa », une exigence directe, presque impolie. Pour un souhait irréel ou un désir poli, dites : \all{Ich hätte gern ein Visum.} « J'aimerais avoir un visa. » ou \all{Wenn ich doch ein Visum hätte!} « Si seulement j'avais un visa ! »
\end{attention}

\begin{attention}[Ordre des mots avec würde]
\all{würde} est le verbe conjugué : il occupe la \textbf{2\textsuperscript{e} position} dans la phrase principale, et l'infinitif va \textbf{en fin de phrase} : \all{Ich würde in Deutschland studieren.} Jamais : \all{Ich würde studieren in Deutschland.} Dans la subordonnée avec \all{wenn}, le verbe conjugué va en fin : \all{Wenn ich ein Visum hätte, ...}
\end{attention}

\begin{exoresolu}[Transformation au conditionnel]
\textbf{Énoncé.} Mettez les phrases suivantes au subjonctif II (hypothèse irréelle) :
\begin{enumerate}
\item \all{Ich habe kein Visum. Ich kann nicht nach Deutschland reisen.}
\item \all{Er hat keine Arbeit. Er wandert aus.}
\item \all{Wir haben kein Geld. Wir bleiben nicht hier.}
\end{enumerate}
\tcblower
\textbf{Solution détaillée.}
\begin{enumerate}
\item \all{Wenn ich ein Visum hätte, könnte ich nach Deutschland reisen.} « Si j'avais un visa, je pourrais voyager en Allemagne. » — \all{haben} → \all{hätte} (prétérit \all{hatte} + tréma) ; \all{können} → \all{könnte}.
\item \all{Wenn er Arbeit hätte, würde er nicht auswandern.} « S'il avait du travail, il n'émigrerait pas. » — verbe plein \all{auswandern} : on utilise \all{würde} + infinitif \all{auswandern} en fin de phrase (la particule \all{aus-} reste soudée à l'infinitif !).
\item \all{Wenn wir Geld hätten, würden wir hier bleiben.} « Si nous avions de l'argent, nous resterions ici. » — \all{wir hätten}, \all{wir würden} : terminaison \all{-en} à la 1\textsuperscript{re} personne du pluriel.
\end{enumerate}
\textbf{Méthode :} (1) repérer le verbe de la condition → le mettre au subjonctif II dans la subordonnée \all{wenn} (verbe en fin) ; (2) le verbe principal : \all{wäre}/\all{hätte}/\all{könnte} si possible, sinon \all{würde} + infinitif final.
\end{exoresolu}

\begin{savaistu}[Le subjonctif II au Bac]
Au baccalauréat, le subjonctif II est très souvent testé par la \textbf{transformation} d'une phrase au conditionnel ou par la traduction de phrases du type « Si j'avais..., je... ». Bonne nouvelle : maîtriser seulement quatre formes — \all{wäre}, \all{hätte}, \all{könnte} et \all{würde} + infinitif — suffit dans environ 90 \% des cas. Apprenez-les par cœur avec leurs traductions : « je serais », « j'aurais », « je pourrais », « je ferais ».
\end{savaistu}

\begin{aretenir}[L'essentiel de la section]
\begin{itemize}
\item \textbf{Parfait} (oral) : \all{haben}/\all{sein} + participe II ; \all{sein} pour les verbes de mouvement : \all{ich bin ausgewandert}, \all{sie ist geflohen}.
\item \textbf{Prétérit} (récit écrit, presse) : \all{verließ}, \all{floh}, \all{beantragte} ; mais \all{war}, \all{hatte}, \all{konnte} s'emploient aussi à l'oral.
\item \textbf{Subjonctif II} = prétérit + tréma : \all{war → wäre}, \all{hatte → hätte}, \all{konnte → könnte} ; cas général : \all{würde} + infinitif en fin de phrase.
\item Emplois : souhait irréel (\all{Wenn ich doch ein Visum hätte!}), hypothèse (\all{Wenn ich reich wäre, würde ich bleiben.}), conseil (\all{Du solltest...}).
\item Le conditionnel français correspond au subjonctif II allemand : « j'aurais » = \all{ich hätte}, « je resterais » = \all{ich würde bleiben}.
\end{itemize}
\end{aretenir}

\begin{exercice}[À vous de jouer]
\begin{enumerate}
\item Conjuguez au parfait : \all{auswandern} (er), \all{fliehen} (sie plur.), \all{beantragen} (ich), \all{verlassen} (wir).
\item Mettez au prétérit : \all{Er hat kein Visum. Er kann nicht fliehen.}
\item Traduisez en allemand : « Si j'avais de l'argent, je n'émigrerais pas. » / « Si seulement je pouvais rester ! » / « Tu devrais demander un visa à l'ambassade. »
\end{enumerate}
\end{exercice}

\begin{corrige}[Exercice « À vous de jouer »]
\begin{enumerate}
\item \all{er ist ausgewandert} (mouvement → sein) ; \all{sie sind geflohen} (mouvement → sein, participe fort \all{geflohen}) ; \all{ich habe beantragt} (préfixe inséparable \all{be-} : pas de \all{ge-}) ; \all{wir haben verlassen} (transitif → haben).
\item \all{Er hatte kein Visum. Er konnte nicht fliehen.}
\item \all{Wenn ich Geld hätte, würde ich nicht auswandern.} / \all{Wenn ich doch bleiben könnte!} / \all{Du solltest bei der Botschaft ein Visum beantragen.}
\end{enumerate}
\end{corrige}

\coursec{Commentaire et production guidée}

Dans la section précédente, nous avons étudié le parfait, le prétérit et le subjonctif II. Ces temps ne sont pas de simples exercices de grammaire : ce sont les outils qui vont vous permettre de \emph{parler du texte} (résumer au prétérit ce que l'auteur dit) et de \emph{donner votre avis} (formuler des hypothèses et des conseils au subjonctif II). Cette dernière section met donc tout en œuvre : lire, comprendre, commenter, débattre et rédiger. C'est exactement ce qui vous sera demandé au baccalauréat.

\courssub{La méthode du commentaire de texte}

Commenter un texte, en allemand comme en français, ce n'est ni le réciter ni le résumer mot à mot : c'est le \cle{présenter}, en \cle{dégager les idées principales} et \cle{prendre position}. En allemand, on suit toujours la même architecture en trois parties : \all{die Einleitung} (l'introduction), \all{der Hauptteil} (le développement), \all{der Schluss} (la conclusion).

\begin{methode}[La structure du commentaire en trois étapes]
\begin{enumerate}
\item \textbf{Einleitung (introduction)} : présenter le texte en une ou deux phrases — le thème, éventuellement l'auteur ou la source, et la problématique. Formule de base : \all{Der Text handelt von ...} « Le texte traite de ... ». On peut ajouter : \all{Der Autor beschreibt die Situation der jungen Auswanderer.} « L'auteur décrit la situation des jeunes émigrants. »
\item \textbf{Hauptteil (développement)} : dégager les idées principales et les appuyer par des éléments du texte (courtes citations ou paraphrases). Formule de base : \all{Der Autor zeigt, dass ...} « L'auteur montre que ... ». Utiliser les connecteurs de cause et de but étudiés dans cette leçon : \all{weil}, \all{da}, \all{denn}, \all{um ... zu}, \all{damit}.
\item \textbf{Schluss (conclusion)} : donner son avis personnel, éventuellement avec une hypothèse au subjonctif II. Formule de base : \all{Meiner Meinung nach ...} « À mon avis ... ».
\end{enumerate}
\end{methode}

\begin{popfigure}
\begin{center}
\begin{tikzpicture}[
  block/.style={rectangle, rounded corners=3pt, draw=popblue, fill=popblueL, text width=3.4cm, align=center, minimum height=1.6cm, font=\small},
  fleche/.style={-{Stealth[length=3mm]}, thick, poporange}
]
\node[block, align=center] (e) at (0,0){\textbf{Einleitung}\\ Der Text handelt von ...};
\node[block, draw=poppurple, fill=poppurpleL, align=center] (h) at (5.2,0){\textbf{Hauptteil}\\ Argumente + Beispiele\\ Der Autor zeigt, dass ...};
\node[block, draw=popgreen, fill=popgreenL, align=center] (s) at (10.4,0){\textbf{Schluss}\\ Meine Meinung\\ Meiner Meinung nach ...};
\draw[fleche] (e) -- (h);
\draw[fleche] (h) -- (s);
\end{tikzpicture}
\end{center}
\legende{Les trois blocs du commentaire de texte : introduction, développement, conclusion.}
\end{popfigure}

\begin{aretenir}[Les temps du commentaire]
Pour parler du contenu du texte, on emploie le \cle{présent} (\all{Der Autor zeigt, dass viele Jugendliche keine Arbeit finden.} « L'auteur montre que beaucoup de jeunes ne trouvent pas de travail. ») ou le \cle{prétérit} pour rapporter des faits passés (\all{Viele junge Leute verließen ihre Heimat.} « Beaucoup de jeunes quittèrent leur pays. »). Pour l'avis personnel et les hypothèses, on emploie le \cle{subjonctif II} : \all{Wenn die Heimat mehr Arbeit anbieten würde, würden weniger Menschen fliehen.} « Si le pays offrait plus de travail, moins de gens fuiraient. »
\end{aretenir}

\courssub{Un commentaire modèle}

Lisez le commentaire suivant, rédigé sur notre texte support « Warum so viele junge Leute auswandern ». Repérez les trois parties et les connecteurs.

\begin{textebox}[Commentaire modèle (8 à 10 lignes)]
Der Text handelt von den Problemen der Auswanderung. Viele junge Leute verlassen ihre Heimat, weil sie keine Arbeit finden. Sie lernen Deutsch, um sich im Ausland zu integrieren. Aber die Einreise ist schwierig, denn man braucht ein Visum. Meiner Meinung nach ist die Auswanderung keine einfache Lösung: Wenn die Heimat mehr Arbeit anbieten würde, würden weniger Menschen fliehen.
\end{textebox}

\textbf{Analyse du modèle.} Phrase 1 : introduction avec \all{handelt von}. Phrases 2 à 4 : développement avec deux causes (\all{weil}, \all{denn}) et un but (\all{um ... zu}). Dernière phrase : conclusion avec \all{Meiner Meinung nach} et une hypothèse au subjonctif II (\all{würde anbieten}, \all{würden fliehen}). C'est exactement le schéma attendu à l'examen.

\courssub{Les expressions pour donner son avis}

\begin{propriete}[Exprimer son opinion]
\begin{itemize}
\item \all{Meiner Meinung nach ...} « À mon avis ... » — tournure impersonnelle, très appréciée à l'écrit.
\item \all{Ich finde, dass ...} / \all{Ich denke, dass ...} / \all{Ich glaube, dass ...} « Je trouve / pense / crois que ... » — suivies d'une subordonnée avec le verbe conjugué à la fin.
\item \all{Einerseits ... andererseits ...} « D'un côté ... de l'autre ... » — pour présenter le pour et le contre.
\item \all{Ich bin der Meinung, dass ...} « Je suis d'avis que ... » — registre plus soutenu.
\end{itemize}
\end{propriete}

\begin{exemplebox}[Exemples traduits]
\all{An deiner Stelle würde ich zuerst Deutsch lernen.} « À ta place, j'apprendrais d'abord l'allemand. »

\all{Einerseits sucht man im Ausland Arbeit, andererseits vermisst man die Familie.} « D'un côté on cherche du travail à l'étranger, de l'autre la famille nous manque. »

\all{Ich finde, dass die Regierung mehr Jobs schaffen muss.} « Je trouve que le gouvernement doit créer plus d'emplois. »

\all{Meiner Meinung nach ist die Flucht gefährlich.} « À mon avis, la fuite est dangereuse. »
\end{exemplebox}

\begin{attention}[La place du verbe après « meiner Meinung nach »]
\all{Meiner Meinung nach} occupe la première position de la phrase : le \cle{verbe conjugué vient donc en deuxième position}, immédiatement après. On dit : \all{Meiner Meinung nach \textbf{ist} die Auswanderung schwierig.} et non \emph{Meiner Meinung nach die Auswanderung ist ...}. Même piège avec \all{einerseits} et \all{andererseits} : \all{Einerseits \textbf{sucht} man Arbeit} (verbe en position 2). C'est l'erreur la plus fréquente des francophones, car en français l'ordre ne change pas (« À mon avis, l'émigration est difficile »). En revanche, après \all{ich finde, dass ...}, le verbe de la subordonnée va \cle{à la fin} : \all{Ich finde, dass die Auswanderung schwierig \textbf{ist}.}
\end{attention}

\begin{center}
\begin{tabularx}{\linewidth}{|X|X|X|}
\hline
\rowcolor{popblueL} Français & Deutsch & Remarque \\
\hline
À mon avis, ... & Meiner Meinung nach ... & verbe en position 2 \\
\hline
Je trouve que ... & Ich finde, dass ... & verbe à la fin dans la subordonnée \\
\hline
D'un côté ... de l'autre ... & Einerseits ... andererseits ... & verbe en position 2 après chaque connecteur \\
\hline
À ta place, je ... & An deiner Stelle würde ich ... & subjonctif II de conseil \\
\hline
Je suis d'avis que ... & Ich bin der Meinung, dass ... & registre soutenu \\
\hline
\end{tabularx}
\end{center}

\courssub{Production écrite guidée}

\begin{methode}[Rédiger le paragraphe « Warum wandern junge Kameruner aus ? »]
\begin{enumerate}
\item Commencez par une phrase d'introduction : \all{Viele junge Kameruner wandern aus.} « Beaucoup de jeunes Camerounais émigrent. »
\item Donnez \textbf{au moins deux causes} avec \all{weil}, \all{da} ou \all{denn} : le chômage (\all{die Arbeitslosigkeit}), la pauvreté (\all{die Armut}), le manque de perspectives (\all{keine Perspektive}).
\item Donnez \textbf{au moins deux buts} avec \all{um ... zu} ou \all{damit} : trouver du travail, étudier, soutenir la famille.
\item Terminez par une phrase d'opinion avec \all{Meiner Meinung nach} et, si possible, une hypothèse au subjonctif II.
\item Relisez : verbe en position 2 dans la phrase principale, verbe à la fin après \all{weil}/\all{da}/\all{dass}, infinitif avec \all{zu} en fin de phrase après \all{um}.
\end{enumerate}
\end{methode}

\begin{exoresolu}[Production écrite : un paragraphe complet]
Énoncé : rédigez un paragraphe de 6 à 8 lignes intitulé \all{Warum wandern junge Kameruner aus ?} avec au moins deux causes (\all{weil}/\all{da}/\all{denn}) et deux buts (\all{um ... zu}/\all{damit}).
\tcblower
Solution détaillée :

\all{Viele junge Kameruner wandern heute aus. Sie verlassen ihre Heimat, weil die Arbeitslosigkeit sehr hoch ist. Da sie zu Hause keine Perspektive haben, suchen sie ihr Glück im Ausland. Viele lernen Deutsch oder Englisch, um im Ausland studieren zu können. Andere arbeiten hart, damit sie Geld an ihre Familien schicken können. Meiner Meinung nach ist die Auswanderung verständlich, aber traurig: Wenn Kamerun mehr Arbeit anbieten würde, würden die jungen Leute gern bleiben.}

Traduction : « Beaucoup de jeunes Camerounais émigrent aujourd'hui. Ils quittent leur pays parce que le chômage est très élevé. Comme ils n'ont pas de perspective chez eux, ils cherchent leur bonheur à l'étranger. Beaucoup apprennent l'allemand ou l'anglais pour pouvoir étudier à l'étranger. D'autres travaillent dur afin de pouvoir envoyer de l'argent à leurs familles. À mon avis, l'émigration est compréhensible mais triste : si le Cameroun offrait plus de travail, les jeunes resteraient volontiers. »

Vérification : deux causes (\all{weil}, \all{da}), deux buts (\all{um ... zu können}, \all{damit}), avis personnel, une phrase au subjonctif II. Toutes les règles de place du verbe sont respectées : verbe à la fin après \all{weil} et \all{da}, verbe en position 2 après \all{Meiner Meinung nach}.
\end{exoresolu}

\courssub{Production orale : le débat « Faut-il partir ou rester ? »}

\begin{experience}[Débat en binômes : « Auswandern oder bleiben ? »]
Par groupes de deux, l'un défend le départ (\all{die Auswanderung}), l'autre défend le fait de rester au pays. Chaque élève doit :
\begin{itemize}
\item donner au moins deux arguments avec \all{weil}, \all{da}, \all{denn} ou \all{um ... zu} ;
\item donner un conseil à son partenaire avec le subjonctif II : \all{An deiner Stelle würde ich ...} « À ta place, je ... » ;
\item réagir à l'argument de l'autre avec \all{einerseits ... andererseits ...}.
\end{itemize}
Exemples de répliques : \all{An deiner Stelle würde ich zuerst mein Studium hier beenden.} « À ta place, je terminerais d'abord mes études ici. » — \all{Ich würde nicht illegal fliehen, denn das ist zu gefährlich.} « Je ne fuirais pas illégalement, car c'est trop dangereux. » — \all{Wenn ich ein Visum hätte, würde ich in Deutschland arbeiten.} « Si j'avais un visa, je travaillerais en Allemagne. »
\end{experience}

\courssub{Correction collective : la grille d'évaluation}

Après la production écrite ou orale, on corrige ensemble avec une grille simple en trois critères. C'est la grille utilisée au baccalauréat.

\begin{center}
\begin{tabularx}{\linewidth}{|l|X|X|}
\hline
\rowcolor{popblueL} Critère & Was wird geprüft ? (ce qu'on vérifie) & Beispiele (exemples) \\
\hline
Lexique & vocabulaire du thème, articles et pluriels corrects & \all{die Auswanderung}, \all{das Visum}, \all{die Grenze}, \all{die Arbeitslosigkeit} \\
\hline
Grammaire & place du verbe, connecteurs de cause et de but, subjonctif II & \all{weil sie keine Arbeit finden} ; \all{um zu studieren} ; \all{ich würde bleiben} \\
\hline
Cohérence & structure en trois parties, idées liées, avis personnel clair & Einleitung — Hauptteil — Schluss \\
\hline
\end{tabularx}
\end{center}

\begin{exercice}[À vous de jouer]
\begin{enumerate}
\item Remettez les mots dans l'ordre : \all{Meiner Meinung nach / ist / gefährlich / die Flucht / sehr}.
\item Complétez avec \all{weil}, \all{denn} ou \all{um ... zu} : \all{Er lernt Deutsch, ... in Deutschland zu arbeiten. Er bleibt nicht hier, ... er keine Arbeit findet.}
\item Mettez le conseil au subjonctif II : \all{Ich lerne zuerst Deutsch.} → \all{An deiner Stelle ...}
\item Rédigez un court commentaire (6 à 8 lignes) sur le texte support « Warum so viele junge Leute auswandern » en suivant le schéma Einleitung — Hauptteil — Schluss.
\end{enumerate}
\end{exercice}

\begin{corrige}[Exercice « À vous de jouer »]
\begin{enumerate}
\item \all{Meiner Meinung nach ist die Flucht sehr gefährlich.} (verbe \all{ist} en position 2, juste après l'expression d'opinion.)
\item \all{Er lernt Deutsch, \textbf{um} in Deutschland zu arbeiten. Er bleibt nicht hier, \textbf{weil} er keine Arbeit findet.} (avec \all{weil}, le verbe \all{findet} va à la fin ; avec \all{denn}, il faudrait dire : \all{Er bleibt nicht hier, denn er findet keine Arbeit.} — verbe en position 2.)
\item \all{An deiner Stelle \textbf{würde} ich zuerst Deutsch \textbf{lernen}.} (subjonctif II : \all{würde} + infinitif à la fin.)
\item Modèle de commentaire : \all{Der Text handelt von der Auswanderung junger Menschen. Der Autor zeigt, dass viele Jugendliche ihre Heimat verlassen, weil sie keine Arbeit finden. Sie lernen Fremdsprachen, um sich im Ausland zu integrieren. Aber die Einreise ist schwierig, denn man braucht ein Visum. Meiner Meinung nach sollte die Heimat mehr Chancen bieten: Wenn es mehr Jobs gäbe, würden weniger junge Leute auswandern.} Traduction : « Le texte traite de l'émigration des jeunes. L'auteur montre que beaucoup de jeunes quittent leur pays parce qu'ils ne trouvent pas de travail. Ils apprennent des langues étrangères pour s'intégrer à l'étranger. Mais l'entrée est difficile, car il faut un visa. À mon avis, le pays devrait offrir plus de chances : s'il y avait plus d'emplois, moins de jeunes émigreraient. »
\end{enumerate}
\end{corrige}

\begin{savaistu}[Le jour du Bac]
Au baccalauréat, la production écrite sur ce thème demande souvent 8 à 10 lignes. Pour valoriser votre copie : utilisez \cle{au moins trois connecteurs différents} (par exemple \all{weil}, \all{denn}, \all{um ... zu}), glissez \cle{une phrase au subjonctif II} (\all{Wenn ... würde ...}) et terminez par un avis personnel avec \all{Meiner Meinung nach}. Les correcteurs récompensent la variété des structures bien plus que la longueur du texte. Et n'oubliez jamais : majuscule à tous les noms, verbe en position 2 dans la phrase principale !
\end{savaistu}

\newpage

\coursec{Activité d'intégration}

\courssub{Objectifs de l'activité}

Cette activité d'intégration clôt la leçon sur les problèmes d'émigration. Elle vise à faire réutiliser, dans une situation de communication orale et écrite authentique, l'ensemble des acquis de la séquence :

\begin{itemize}
\item le \cle{lexique thématique} de l'émigration : \all{die Auswanderung}, \all{die Flucht}, \all{das Visum}, \all{die Grenze}, \all{die Arbeitslosigkeit}, \all{das Ausland}, \all{die Botschaft}, \all{der Antrag} ;
\item les \cle{connecteurs de cause} \all{weil}, \all{denn}, \all{da} et les \cle{connecteurs de but} \all{damit}, \all{um ... zu} ;
\item le \cle{parfait} et le \cle{prétérit} pour raconter un parcours migratoire ;
\item le \cle{subjonctif II} pour exprimer l'hypothèse irréelle et le regret ;
\item la capacité à poser des questions, à répondre à l'oral, puis à rédiger un court article au passé.
\end{itemize}

À l'issue de l'activité, l'élève doit être capable de conduire un entretien simple en allemand et de le transformer en texte rédigé, en respectant l'ordre des mots, la place du verbe et les temps du passé. Ces compétences correspondent exactement aux attendus du Baccalauréat : compréhension de l'oral et de l'écrit, expression orale en interaction et production écrite guidée.

\courssub{Situation}

La radio scolaire germanophone du lycée de Yaoundé, \all{„Radio Welle Kamerun“}, prépare une émission spéciale consacrée à l'émigration des jeunes. L'animateur a invité un ancien élève du lycée, Serge Mbarga, 24 ans, qui a quitté Douala il y a deux ans pour faire une formation professionnelle (\all{eine Ausbildung}) en Allemagne. L'émission sera ensuite résumée dans le journal de l'école.

\begin{textebox}[Annonce de l'émission — „Radio Welle Kamerun“]
„Liebe Hörerinnen und Hörer! Heute sprechen wir über ein wichtiges Thema: die Auswanderung junger Kameruner. Unser Gast ist Serge Mbarga aus Douala. Er ist vor zwei Jahren nach Deutschland ausgewandert, weil er dort eine Ausbildung machen wollte. Wie hat er das Visum bekommen? Was hat er erlebt? Was würde er anders machen, wenn er noch einmal anfangen könnte? Bleiben Sie dran!“
\end{textebox}

\begin{definition}[Glossaire de l'annonce]
\all{der Hörer, -} : l'auditeur ; \all{die Auswanderung, -en} : l'émigration ; \all{auswandern (ist ausgewandert)} : émigrer ; \all{die Ausbildung, -en} : la formation professionnelle ; \all{erleben (hat erlebt)} : vivre, éprouver ; \all{noch einmal anfangen (fing an, hat angefangen)} : recommencer ; \all{Bleiben Sie dran!} : restez à l'écoute !
\end{definition}

\begin{attention}[Observation grammaticale dans l'annonce]
Repérez dans le texte trois structures clés de la leçon : la subordonnée de cause avec verbe à la fin (\all{weil er dort eine Ausbildung machen wollte}), le parfait avec l'auxiliaire \all{sein} pour un verbe de déplacement (\all{ist ... ausgewandert}) et le subjonctif II pour l'hypothèse (\all{würde er anders machen, wenn er ... anfangen könnte}).
\end{attention}

\courssub{Consignes}

Travail par binômes : l'élève A joue le journaliste de la radio allemande, l'élève B joue Serge, le jeune émigré. Suivez les étapes dans l'ordre.

\begin{enumerate}
\item \textbf{Préparation du journaliste.} Rédigez cinq questions avec les interrogatifs corrects (\all{Warum}, \all{Wie}, \all{Was}, \all{Wann}, \all{Wohin}). Exemples : \all{Warum sind Sie ausgewandert?} \all{Was haben Sie gemacht, um ein Visum zu bekommen?} \all{Wie war Ihre erste Woche in Deutschland?}
\item \textbf{Préparation de l'émigré.} Notez des idées de réponses. Vos réponses doivent contenir obligatoirement : une phrase avec \all{weil}, une avec \all{denn}, une avec \all{um ... zu}, une avec \all{damit}, et une phrase au subjonctif II du type \all{Wenn ich kein Visum bekommen hätte, wäre ich in Kamerun geblieben.}
\item \textbf{Jeu de rôle oral.} Jouez l'interview devant votre voisin de table (5 minutes). Le journaliste note les informations essentielles.
\item \textbf{Rédaction de l'article.} Ensemble, rédigez un court article de 8 à 10 lignes pour le journal de l'école, au prétérit et au parfait, avec un titre, une introduction, le résumé des réponses et une conclusion.
\item \textbf{Restitution.} Deux binômes présentent leur interview devant la classe ; correction collective au tableau (lexique, connecteurs, temps du passé, ordre des mots).
\end{enumerate}

\begin{attention}[Pièges à éviter]
Après \all{weil} et \all{damit}, le verbe conjugué va à la fin : \all{..., weil er eine Ausbildung machen wollte.} Après \all{denn}, le verbe reste en deuxième position : \all{Er ging weg, denn er hatte keine Arbeit.} Ne confondez pas \all{damit} (sujets différents dans les deux propositions) et \all{um ... zu} (même sujet). Au passé, les verbes de déplacement comme \all{auswandern} se conjuguent avec \all{sein} : \all{er ist ausgewandert}. Enfin, n'oubliez pas la majuscule à tous les noms : \all{das Visum}, \all{die Grenze}, \all{die Arbeit}.
\end{attention}

\courssub{Ressources à mobiliser}

\begin{propriete}[Rappel : cause et but]
\begin{itemize}
\item \all{weil} / \all{da} + verbe conjugué à la fin (subordonnée) : \all{Er wanderte aus, weil er keine Arbeit fand.} (\all{da} se place plutôt en tête de phrase : \all{Da er keine Arbeit fand, wanderte er aus.})
\item \all{denn} + verbe en position 2 (coordination, sans inversion) : \all{Er wanderte aus, denn er fand keine Arbeit.}
\item \all{um ... zu} + infinitif : même sujet dans les deux parties. \all{Er sparte Geld, um das Visum zu bezahlen.}
\item \all{damit} + verbe conjugué à la fin : sujets différents. \all{Er lernte Deutsch, damit seine Familie stolz auf ihn sein konnte.}
\end{itemize}
\end{propriete}

\begin{propriete}[Rappel : passé et subjonctif II]
Récit oral et conversation : \all{ich bin ausgewandert}, \all{ich habe Deutsch gelernt} ; article écrit et presse : \all{er wanderte aus}, \all{er hatte}, \all{er war}, \all{er bekam}. Subjonctif II du passé : \all{hätte} / \all{wäre} + participe II : \all{Wenn ich mehr Geld gehabt hätte, wäre ich früher geflogen.} On choisit \all{hätte} avec les verbes qui prennent \all{haben} au parfait, et \all{wäre} avec ceux qui prennent \all{sein}.
\end{propriete}

Lexique utile : \all{die Arbeitslosigkeit} (le chômage), \all{die Botschaft, -en} (l'ambassade), \all{der Antrag, die Anträge} (la demande officielle), \all{die Grenze, -n} (la frontière), \all{das Ausland} (l'étranger), \all{die Heimat} (la patrie), \all{sich integrieren} (s'intégrer), \all{verdienen} (gagner de l'argent), \all{vermissen} (regretter, éprouver le manque de), \all{einen Antrag stellen} (déposer une demande), \all{das Sprachdiplom, -e} (le diplôme de langue).

\courssub{Proposition de production et corrigé commenté}

\begin{exoresolu}[Modèle d'article pour le journal de l'école]
Rédigez un article de 8 à 10 lignes résumant l'interview de Serge Mbarga.
\tcblower
\textbf{„Von Douala nach München : Serge erzählt“}

\all{Gestern hat unsere Radiosendung Serge Mbarga empfangen, einen ehemaligen Schüler unseres Lyzeums. Er ist vor zwei Jahren nach Deutschland ausgewandert, weil er in Kamerun keine Arbeit gefunden hat. Er sagte: „Ich habe ein Jahr Deutsch gelernt, um das Sprachdiplom zu bekommen.“ Dann hat er bei der deutschen Botschaft einen Antrag gestellt, damit die Beamten ihm ein Visum gaben. Die erste Woche in München war schwer, denn er vermisste seine Familie und das Essen aus Douala. Heute macht er eine Ausbildung und verdient sein eigenes Geld. „Wenn ich kein Visum bekommen hätte, wäre ich in Kamerun geblieben“, erklärte er zum Schluss. Sein Rat an die Jugend: „Lernt Sprachen und glaubt an eure Zukunft!“}

\textbf{Commentaire des choix de langue :}
\begin{itemize}
\item Temps : parfait pour le récit (\all{hat empfangen}, \all{ist ausgewandert}, \all{hat gelernt}) ; prétérit pour \all{war}, \all{gaben} et \all{vermisste}, comme dans la presse écrite.
\item Cause : \all{weil} + verbe final (\all{weil er ... gefunden hat}) ; \all{denn} + verbe en position 2 (\all{denn er vermisste}).
\item But : \all{um ... zu} avec le même sujet (Serge apprend, Serge obtient le diplôme) ; \all{damit} avec deux sujets différents (Serge dépose la demande, les fonctionnaires donnent le visa) — c'est la règle de choix essentielle.
\item Subjonctif II du passé : \all{hätte ... bekommen}, \all{wäre ... geblieben} pour l'hypothèse irréelle ; \all{wäre} s'impose devant \all{geblieben} car \all{bleiben} se conjugue avec \all{sein}.
\item Discours rapporté : guillemets allemands „...“ et verbes introducteurs \all{sagte}, \all{erklärte}.
\end{itemize}
\end{exoresolu}

\begin{methode}[De l'interview à l'article : la démarche]
\begin{enumerate}
\item Relever dans les notes les faits essentiels : dates, lieux, raisons du départ, étapes du parcours.
\item Classer les informations : introduction (qui ? quand ?), développement (pourquoi ? comment ?), conclusion (bilan, conseil).
\item Choisir les temps : parfait pour raconter les faits, prétérit pour \all{sein}, \all{haben} et les verbes modaux, subjonctif II pour les regrets.
\item Vérifier chaque connecteur : même sujet = \all{um ... zu} ; sujets différents = \all{damit} ; cause en coordination = \all{denn}.
\item Relire en contrôlant la place du verbe, les majuscules des noms et les Umlaute.
\end{enumerate}
\end{methode}

\courssub{Bilan de l'activité}

L'activité a permis de mobiliser simultanément les quatre savoir-faire de la leçon : comprendre et poser des questions, employer le lexique de l'émigration dans un contexte réaliste, articuler cause et but avec les bons connecteurs, et raconter au passé en alternant parfait, prétérit et subjonctif II. La restitution orale a montré l'importance de la place du verbe après \all{weil} et \all{damit}, point faible fréquent des francophones, ainsi que le choix entre \all{um ... zu} et \all{damit} selon l'identité ou la différence des sujets. Pour aller plus loin, les élèves pourront enregistrer leur interview et rédiger une lettre à un ami resté au Cameroun, en réutilisant les mêmes outils grammaticaux — un excellent entraînement à l'épreuve de production écrite du Baccalauréat.

\newpage

\coursec{Méthodes et exercices résolus}

\courssub{Savoir-faire 1 : comprendre un texte sur l'émigration (global puis détaillé)}

\begin{methode}[Comprendre un texte de presse sur l'émigration]
\begin{enumerate}
\item \textbf{Lire le titre et la source.} Le titre \all{„Warum so viele junge Leute auswandern“} annonce un texte \emph{explicatif} (question avec \all{warum}) : on attend des \cle{causes} et des \cle{conséquences}.
\item \textbf{Première lecture globale} : repérer le thème (ici : \all{die Auswanderung} / \all{die Abwanderung}), le point de vue de l'auteur, et les chiffres ou mots-clés (\all{Arbeitslosigkeit, Visum, Grenze, Ausland}).
\item \textbf{Deuxième lecture détaillée} : souligner les connecteurs logiques. \all{weil, denn, da} introduisent des causes ; \all{damit, um … zu} des buts ; \all{deshalb, darum} des conséquences. Ce sont eux qui structurent l'argumentation.
\item \textbf{Repérer les temps} : \all{Perfekt} ou \all{Präteritum} pour les faits passés, \all{Präsens} pour la situation actuelle, \all{Konjunktiv II} pour les souhaits et hypothèses (\all{Viele Jugendliche würden gern im Ausland arbeiten}).
\item \textbf{Répondre aux questions} en reformulant avec ses propres mots, jamais en recopiant le texte mot à mot ; citer une ligne du texte seulement pour justifier (\all{Im Text steht: „…“}).
\item \textbf{Vérifier} : chaque réponse doit être une phrase complète, avec le verbe conjugué en deuxième position (ou en fin dans une subordonnée).
\end{enumerate}
\end{methode}

\begin{exoresolu}[Compréhension de texte — type Bac]
\textbf{Texte} : \all{Viele junge Kameruner verlassen ihr Land, weil sie keine Arbeit finden. Sie hoffen, dass sie im Ausland ein besseres Leben haben. Aber die Auswanderung ist oft schwierig: Man braucht ein Visum, und die Reise ist teuer und gefährlich. Trotzdem träumen viele Jugendliche davon, in Deutschland oder in der Schweiz zu studieren.}

« Beaucoup de jeunes Camerounais quittent leur pays parce qu'ils ne trouvent pas de travail. Ils espèrent avoir une vie meilleure à l'étranger. Mais l'émigration est souvent difficile : il faut un visa, et le voyage est cher et dangereux. Pourtant, beaucoup de jeunes rêvent d'étudier en Allemagne ou en Suisse. »

\textbf{Questions} : 1. \all{Warum verlassen viele junge Kameruner ihr Land?} 2. \all{Was braucht man für die Auswanderung?} 3. \all{Wovon träumen viele Jugendliche?}
\tcblower
\textbf{Raisonnement} : la question 1 porte sur la \emph{cause} : on cherche la subordonnée introduite par \all{weil}. La question 2 porte sur une \emph{condition matérielle} (\all{man braucht …}). La question 3 porte sur un \emph{souhait} (\all{träumen von + Dativ} ; à l'interrogatif : \all{wovon}).

\textbf{Réponses rédigées} :

1. \all{Viele junge Kameruner verlassen ihr Land, weil sie keine Arbeit finden.} — On reprend la subordonnée causale du texte ; attention : verbe \all{finden} en \textbf{fin de phrase} après \all{weil}.

2. \all{Für die Auswanderung braucht man ein Visum.} — \all{für + Akkusativ} ; \all{man} est le sujet, le verbe \all{braucht} reste en deuxième position car \all{für die Auswanderung} occupe la première.

3. \all{Viele Jugendliche träumen davon, in Deutschland oder in der Schweiz zu studieren.} — Construction figée : \all{träumen davon, … zu + Infinitiv} ; l'infinitif \all{zu studieren} se place en fin.

\textbf{Conclusion} : chaque réponse est une phrase complète, introduite par le thème de la question, sans recopier inutilement tout le paragraphe.
\end{exoresolu}

\courssub{Savoir-faire 2 : employer le vocabulaire de l'émigration}

\begin{methode}[Construire et utiliser le lexique thématique]
\begin{enumerate}
\item \textbf{Apprendre chaque nom avec son article et son pluriel} : \all{die Auswanderung (-en)}, \all{die Flucht (-en)}, \all{das Visum (die Visa)}, \all{die Grenze (-n)}, \all{die Arbeitslosigkeit (Sg.)}, \all{das Ausland (Sg.)}, \all{der Flüchtling (-e)}, \all{der Auswanderer (-) / die Auswanderin (-nen)}.
\item \textbf{Mémoriser les familles de mots} : \all{auswandern} (verbe) → \all{der Auswanderer} (personne) → \all{die Auswanderung} (action) ; \all{fliehen / flüchten} → \all{die Flucht} → \all{der Flüchtling}.
\item \textbf{Apprendre les collocations} (verbe + nom) : \all{ein Visum beantragen} (demander un visa), \all{die Grenze überqueren} (franchir la frontière), \all{ins Ausland gehen} (partir à l'étranger), \all{vor der Arbeitslosigkeit fliehen} (fuir le chômage), \all{in der Heimat bleiben} (rester au pays).
\item \textbf{Distinguer les paires} : \all{die Auswanderung} (quitter son pays, vu du départ) / \all{die Einwanderung} (entrer dans un pays, vu de l'arrivée) ; \all{das Ausland} ≠ \all{die Heimat}.
\item \textbf{Réutiliser} au moins cinq mots du champ lexical dans toute production écrite sur le sujet : c'est un critère de notation explicite au Bac.
\end{enumerate}
\end{methode}

\begin{center}
\begin{tabularx}{\linewidth}{|l|l|X|}
\hline
\rowcolor{popblueL} \textbf{Deutsch} & \textbf{Pluriel} & \textbf{Français / remarque} \\
\hline
\all{die Auswanderung} & \all{-en} & l'émigration (départ) \\
\hline
\all{die Einwanderung} & \all{-en} & l'immigration (arrivée) \\
\hline
\all{der Auswanderer / die Auswanderin} & \all{- / -nen} & l'émigrant / l'émigrante \\
\hline
\all{die Flucht} & \all{-en} & la fuite \\
\hline
\all{der Flüchtling} & \all{-e} & le réfugié \\
\hline
\all{das Visum} & \all{die Visa} & le visa ; \all{ein Visum beantragen} \\
\hline
\all{die Grenze} & \all{-n} & la frontière ; \all{die Grenze überqueren} \\
\hline
\all{die Arbeitslosigkeit} & Sg. & le chômage \\
\hline
\all{das Ausland} & Sg. & l'étranger ; \all{ins Ausland} (mouvement), \all{im Ausland} (lieu) \\
\hline
\all{die Heimat} & \all{-en} & la patrie, le pays natal \\
\hline
\all{der Arbeitsplatz} & \all{-„e} & l'emploi, le poste de travail \\
\hline
\all{die Reise} & \all{-n} & le voyage \\
\hline
\end{tabularx}
\end{center}

\begin{exoresolu}[Lexique — compléter un texte]
\textbf{Énoncé} : Complétez avec le mot qui convient : \all{Visum, Grenze, Arbeitslosigkeit, Ausland, Flüchtlinge}.

\all{1. Wegen der hohen ……… finden viele Jugendliche keinen Job. 2. Er hat ein ……… für Deutschland beantragt. 3. Viele ……… kommen über das Mittelmeer nach Europa. 4. Die ……… zwischen den beiden Ländern ist geschlossen. 5. Mein Bruder arbeitet seit zwei Jahren im ……….}
\tcblower
\textbf{Raisonnement} : on identifie d'abord le genre grammatical exigé par l'article ou la préposition, puis le sens.

1. \all{Arbeitslosigkeit} — \all{wegen + Genitiv} (\all{der hohen Arbeitslosigkeit}) ; sens : « à cause du chômage élevé ».

2. \all{Visum} — \all{ein Visum beantragen} : collocation obligatoire ; \all{das Visum}, pluriel \all{die Visa}.

3. \all{Flüchtlinge} — sujet pluriel (verbe \all{kommen}) ; \all{der Flüchtling, -e}.

4. \all{Grenze} — \all{die Grenze zwischen … ist geschlossen} : « la frontière est fermée ».

5. \all{Ausland} — \all{im Ausland} = \all{in dem Ausland} (datif, lieu) ; « à l'étranger ». À ne pas confondre avec \all{ins Ausland} = \all{in das Ausland} (accusatif, mouvement) : \all{Er geht ins Ausland.}

\textbf{Conclusion} : toujours vérifier l'accord article/nom et la collocation : on \emph{beantragt} un visa, on ne \emph{macht} pas un visa (calque du français « faire un visa »).
\end{exoresolu}

\courssub{Savoir-faire 3 : exprimer la cause (weil, denn, da)}

\begin{methode}[Choisir et construire la proposition causale]
\begin{enumerate}
\item \textbf{\all{weil} (parce que)} : subordonnée, verbe conjugué \textbf{en fin}. Répond à \all{Warum …?}. Peut se placer après ou avant la principale : \all{Weil er kein Visum hatte, blieb er in Yaoundé.} (la subordonnée en tête occupe la première position : le verbe de la principale, \all{blieb}, vient juste après la virgule).
\item \textbf{\all{da} (comme, puisque)} : subordonnée aussi, verbe en fin ; exprime une cause \emph{déjà connue} du lecteur ; se place volontiers en tête : \all{Da die Arbeitslosigkeit steigt, wandern viele aus.}
\item \textbf{\all{denn} (car)} : \textbf{coordonnée}, verbe en \textbf{deuxième position} (comme une phrase normale), jamais en tête de texte : \all{Er blieb, denn er hatte kein Visum.}
\item \textbf{Transformation} : pour passer de \all{weil} à \all{denn}, on déplace le verbe de la fin à la deuxième position : \all{Sie flieht, weil sie Angst hat.} → \all{Sie flieht, denn sie hat Angst.}
\item \textbf{Cause nominale} : avec \all{wegen + Genitiv} (ou datif à l'oral) : \all{wegen der Arbeitslosigkeit} = « à cause du chômage ».
\end{enumerate}
\end{methode}

\begin{center}
\begin{tabularx}{\linewidth}{|l|X|X|}
\hline
\rowcolor{popblueL} \textbf{Conjonction} & \textbf{Place du verbe} & \textbf{Exemple} \\
\hline
\all{weil} (parce que) & fin (subordonnée) & \all{Er bleibt, weil er kein Visum hat.} \\
\hline
\all{da} (comme, puisque) & fin (subordonnée) & \all{Da er kein Visum hat, bleibt er.} \\
\hline
\all{denn} (car) & 2\up{e} position (coordonnée) & \all{Er bleibt, denn er hat kein Visum.} \\
\hline
\all{wegen + Gen.} (à cause de) & — (groupe nominal) & \all{Er bleibt wegen des Visums.} \\
\hline
\end{tabularx}
\end{center}

\begin{exoresolu}[Grammaire — transformation weil / denn]
\textbf{Énoncé} : Transformez avec \all{denn} : \all{1. Viele Jugendliche wandern aus, weil sie keine Arbeit finden. 2. Er konnte nicht reisen, weil er kein Visum bekommen hatte.}
\tcblower
\textbf{Raisonnement} : \all{denn} est une conjonction de coordination : la phrase qui suit garde l'ordre sujet–verbe (verbe en deuxième position). On supprime \all{weil}, on remet le verbe conjugué juste après le sujet.

1. \all{Viele Jugendliche wandern aus, denn sie finden keine Arbeit.} — \all{finden} passe de la fin à la deuxième position ; le reste est inchangé.

2. \all{Er konnte nicht reisen, denn er hatte kein Visum bekommen.} — Au plus-que-parfait, c'est l'auxiliaire \all{hatte} qui prend la deuxième position ; le participe \all{bekommen} reste en fin.

\textbf{Conclusion} : la transformation ne change ni le sens ni les temps ; seule la \emph{place du verbe} change. C'est exactement ce que le correcteur vérifie.
\end{exoresolu}

\courssub{Savoir-faire 4 : exprimer le but (damit, um … zu)}

\begin{methode}[Construire la proposition finale]
\begin{enumerate}
\item \textbf{\all{um … zu + Infinitiv}} : quand le sujet des deux verbes est \textbf{le même}. L'infinitif avec \all{zu} va en fin : \all{Er lernt Deutsch, um in Deutschland zu studieren.}
\item \textbf{\all{damit} + subordonnée} (verbe en fin) : quand les sujets sont \textbf{différents} : \all{Die Eltern sparen Geld, damit ihre Kinder im Ausland studieren können.}
\item \textbf{Test rapide} : même sujet → \all{um … zu} ; sujets différents → \all{damit}. Ne jamais écrire \all{um dass} ni \all{damit … zu}.
\item \textbf{Verbes à particule} : le \all{zu} s'insère entre la particule et le radical : \all{auszuwandern}, \all{fortzugehen} : \all{Sie spart Geld, um auszuwandern.}
\item \textbf{Équivalents français} : « pour », « afin de » → \all{um … zu} ; « pour que, afin que » → \all{damit}.
\end{enumerate}
\end{methode}

\begin{exoresolu}[Thème — traduction avec but et cause]
\textbf{Énoncé} : Traduisez en allemand : 1. « Beaucoup de jeunes apprennent l'allemand pour travailler en Allemagne. » 2. « Le gouvernement crée des emplois pour que les jeunes restent au pays. » 3. « Il est parti à l'étranger parce qu'il ne trouvait pas de travail. »
\tcblower
\textbf{Raisonnement} : phrase 1 : même sujet (\all{viele junge Leute}) → \all{um … zu}. Phrase 2 : sujets différents (\all{die Regierung} / \all{die Jugendlichen}) → \all{damit}. Phrase 3 : cause → \all{weil}, verbe en fin, prétérit de \all{finden} : \all{fand}.

1. \all{Viele junge Leute lernen Deutsch, um in Deutschland zu arbeiten.} — infinitif \all{zu arbeiten} en fin ; majuscule à \all{Deutsch} et \all{Deutschland}.

2. \all{Die Regierung schafft Arbeitsplätze, damit die Jugendlichen im Land bleiben.} — « créer des emplois » = \all{Arbeitsplätze schaffen} ; dans la subordonnée avec \all{damit}, le verbe \all{bleiben} va en fin.

3. \all{Er ist ins Ausland gegangen, weil er keine Arbeit fand.} (ou \all{… weil er keine Arbeit gefunden hat.}) — « partir » exprimé au passé : \all{Perfekt} (\all{ist … gegangen}, verbe de déplacement → auxiliaire \all{sein}) ; la cause peut être au \all{Präteritum} (\all{fand}) ou au \all{Perfekt}.

\textbf{Conclusion} : avant de traduire « pour », toujours se poser la question du sujet ; avant de traduire « parce que », vérifier la place du verbe en fin de subordonnée.
\end{exoresolu}

\courssub{Savoir-faire 5 : manier parfait, prétérit et subjonctif II}

\begin{methode}[Choisir le temps et former le Konjunktiv II]
\begin{enumerate}
\item \textbf{Perfekt} : langue parlée et récits personnels ; auxiliaire \all{haben} (majorité) ou \all{sein} (déplacement, changement d'état : \all{gehen, kommen, auswandern, sterben, aufstehen}) + participe en fin : \all{Er ist nach Deutschland ausgewandert.}
\item \textbf{Präteritum} : langue écrite, presse, récits ; obligatoire pour \all{sein} (\all{war}), \all{haben} (\all{hatte}) et les modaux (\all{musste, konnte, wollte}) même à l'oral : \all{Viele Flüchtlinge kamen über das Meer.}
\item \textbf{Konjunktiv II} : souhait, hypothèse, conseil. Forme de base : \all{würde + Infinitiv} ; formes irrégulières à connaître : \all{wäre} (sein), \all{hätte} (haben), \all{könnte, müsste, dürfte, sollte, wüsste, gäbe}.
\item \textbf{Emplois types} : souhait irréel : \all{Wenn ich ein Visum hätte, würde ich nach Deutschland reisen.} ; rêve : \all{Ich würde gern im Ausland studieren.} ; conseil : \all{Du solltest mehr Deutsch lernen.}
\item \textbf{Phrase hypothétique} : \all{wenn} + KII (verbe en fin) → principale avec \all{würde} en deuxième position : \all{Wenn es mehr Arbeit gäbe, würden die Jugendlichen bleiben.}
\end{enumerate}
\end{methode}

\begin{center}
\begin{tabular}{|l|l|l|}
\hline
\rowcolor{popblueL} \textbf{Infinitiv} & \textbf{Präteritum} & \textbf{Konjunktiv II} \\
\hline
\all{sein} & \all{war} & \all{wäre} \\
\hline
\all{haben} & \all{hatte} & \all{hätte} \\
\hline
\all{werden} & \all{wurde} & \all{würde} \\
\hline
\all{können} & \all{konnte} & \all{könnte} \\
\hline
\all{müssen} & \all{musste} & \all{müsste} \\
\hline
\all{dürfen} & \all{durfte} & \all{dürfte} \\
\hline
\all{sollen} & \all{sollte} & \all{sollte} \\
\hline
\all{wissen} & \all{wusste} & \all{wüsste} \\
\hline
\all{geben (es gibt)} & \all{es gab} & \all{es gäbe} \\
\hline
\end{tabular}
\end{center}

\begin{aretenir}[Formation du Konjunktiv II]
Le KII se forme sur le radical du \textbf{prétérit} ; pour les verbes forts et les modaux qui ont un Umlaut à l'infinitif (\all{a, o, u}), on ajoute l'Umlaut : \all{war → wäre}, \all{hatte → hätte}, \all{konnte → könnte}, \all{gab → gäbe}. Exception : \all{sollen} et \all{wollen} ne prennent pas d'Umlaut (\all{sollte}, \all{wollte}). Pour tous les autres verbes, on utilise \all{würde + Infinitiv} : \all{ich würde arbeiten, wir würden bleiben}.
\end{aretenir}

\begin{exoresolu}[Conjugaison — Perfekt, Präteritum, Konjunktiv II]
\textbf{Énoncé} : a) Mettez au Perfekt : \all{Er wandert nach Kanada aus. / Sie findet keine Arbeit.} b) Mettez au Präteritum : \all{Die Familie hat kein Geld. / Er muss im Land bleiben.} c) Mettez au Konjunktiv II : \all{Ich habe ein Visum. Ich reise nach Europa.}
\tcblower
\textbf{Raisonnement} :

a) \all{Er ist nach Kanada ausgewandert.} — \all{auswandern} exprime un déplacement → auxiliaire \all{sein} ; participe \all{ausgewandert} (verbe faible à particule : \all{aus-} + \all{gewandert}). / \all{Sie hat keine Arbeit gefunden.} — \all{finden} est fort : participe \all{gefunden}, auxiliaire \all{haben}.

b) \all{Die Familie hatte kein Geld.} — prétérit de \all{haben} : \all{hatte}. / \all{Er musste im Land bleiben.} — prétérit de \all{müssen} : \all{musste} (sans Umlaut au prétérit !) ; l'infinitif \all{bleiben} reste en fin.

c) \all{Wenn ich ein Visum hätte, würde ich nach Europa reisen.} — KII de \all{haben} : \all{hätte} ; dans la subordonnée introduite par \all{wenn}, \all{hätte} va en fin ; la principale commence par le verbe \all{würde} car la subordonnée occupe la première position.

\textbf{Conclusion} : retenir les trois réflexes : déplacement → \all{sein} au Perfekt ; \all{haben/sein/modaux} au Präteritum ; souhait → \all{hätte/wäre/würde}.
\end{exoresolu}

\courssub{Savoir-faire 6 : produire un court paragraphe argumenté}

\begin{methode}[Rédiger un paragraphe sur l'émigration (8 à 10 lignes)]
\begin{enumerate}
\item \textbf{Introduction} : une phrase générale avec le vocabulaire du thème : \all{Die Auswanderung ist heute ein wichtiges Thema in Kamerun.}
\item \textbf{Causes} : deux causes reliées par \all{weil/denn/da} : \all{Viele junge Leute wandern aus, weil sie keine Arbeit finden.}
\item \textbf{Buts et espoirs} : une phrase avec \all{um … zu} ou \all{damit} : \all{Sie gehen ins Ausland, um ein besseres Leben zu suchen.}
\item \textbf{Problèmes et dangers} : \all{Aber die Reise ist oft gefährlich, und das Visum ist schwer zu bekommen.}
\item \textbf{Avis personnel au Konjunktiv II} : \all{Ich würde gern in meinem Land bleiben, wenn es mehr Arbeit gäbe.}
\item \textbf{Relecture} : vérifier la place du verbe, les majuscules des noms, les articles, et compter au moins cinq mots du lexique de l'émigration.
\end{enumerate}
\end{methode}

\begin{exoresolu}[Production écrite guidée — type Bac]
\textbf{Énoncé} : \all{Schreiben Sie einen kurzen Text (8–10 Zeilen) über das Thema: „Warum wandern viele junge Kameruner aus? Was denken Sie darüber?“}
\tcblower
\textbf{Plan suivi} : phrase générale → causes (\all{weil}) → buts (\all{um … zu}) → dangers → avis personnel (\all{Konjunktiv II}).

\textbf{Texte rédigé} :

\all{Heute wandern viele junge Kameruner aus, weil sie in ihrer Heimat keine Arbeit finden. Die Arbeitslosigkeit ist sehr hoch, und viele Familien sind arm. Darum träumen die Jugendlichen von einem besseren Leben im Ausland. Sie lernen Deutsch oder Englisch, um in Europa oder in Amerika zu studieren und zu arbeiten. Aber die Auswanderung ist nicht einfach: Man braucht ein Visum, und die Reise ist oft teuer und gefährlich. Manche Flüchtlinge sterben sogar auf dem Weg nach Europa. Ich denke, dass die Regierung mehr Arbeitsplätze schaffen sollte. Wenn es mehr Arbeit gäbe, würden die jungen Leute gern in Kamerun bleiben, denn sie lieben ihr Land.}

« Aujourd'hui, beaucoup de jeunes Camerounais émigrent parce qu'ils ne trouvent pas de travail dans leur pays. Le chômage est très élevé et beaucoup de familles sont pauvres. C'est pourquoi les jeunes rêvent d'une vie meilleure à l'étranger. Ils apprennent l'allemand ou l'anglais pour étudier et travailler en Europe ou en Amérique. Mais l'émigration n'est pas facile : il faut un visa, et le voyage est souvent cher et dangereux. Certains réfugiés meurent même sur la route de l'Europe. Je pense que le gouvernement devrait créer plus d'emplois. S'il y avait plus de travail, les jeunes resteraient volontiers au Cameroun, car ils aiment leur pays. »

\textbf{Justifications grammaticales} : verbe en fin après \all{weil} (\all{finden}) et \all{dass} (\all{schaffen sollte}) ; \all{um … zu studieren und zu arbeiten} (même sujet) ; \all{würden … bleiben} et \all{gäbe} (KII dans l'hypothèse) ; \all{denn} suivi de l'ordre sujet–verbe.

\textbf{Conclusion} : le texte contient plus de cinq mots du champ lexical (\all{Auswanderung, Arbeitslosigkeit, Visum, Ausland, Flüchtlinge, Heimat}) et tous les points de grammaire de la leçon : c'est le niveau attendu pour la note maximale.
\end{exoresolu}

\begin{attention}[Les 5 erreurs qui coûtent des points au Bac]
\begin{enumerate}
\item \textbf{Verbe mal placé après \all{weil} et \all{damit}.} Faux : \all{… weil sie finden keine Arbeit.} Juste : \all{… weil sie keine Arbeit finden.} Le verbe conjugué va \textbf{en fin} de subordonnée. En revanche, après \all{denn}, l'ordre reste normal : \all{… denn sie finden keine Arbeit.}
\item \textbf{Confondre \all{um … zu} et \all{damit}.} On n'emploie \all{um … zu} que si le sujet est identique ; « pour que les jeunes restent » = \all{damit die Jugendlichen bleiben}, et non \all{um die Jugendlichen zu bleiben}.
\item \textbf{Oublier les majuscules et les articles des noms.} \all{ausland}, \all{visum}, \all{grenze} sans majuscule sont sanctionnés ; apprendre \all{das Visum, die Grenze, das Ausland, die Arbeitslosigkeit} avec leur article.
\item \textbf{Faux amis et calques du français.} « Demander un visa » = \all{ein Visum beantragen} (pas \all{fragen}) ; « à l'étranger » = \all{im Ausland} ; \all{die Emigration} existe mais on dit couramment \all{die Auswanderung}.
\item \textbf{Mauvais auxiliaire au Perfekt et KII déformé.} Verbes de déplacement → \all{sein} : \all{Er ist ausgewandert} (pas \all{hat ausgewandert}). KII : \all{hätte, wäre, könnte, gäbe} avec Umlaut ; ne pas confondre \all{wurde} (prétérit de \all{werden}, aussi auxiliaire du passif) avec \all{würde} (KII).
\end{enumerate}
\end{attention}

\newpage

\coursec{Exercices}

\courssub{Je vérifie mes connaissances}

\begin{exercice}[QCM : vocabulaire de l'émigration]
Entoure la bonne réponse.
\begin{enumerate}
\item \all{die Auswanderung} signifie :
\begin{enumerate}
\item l'immigration \item l'émigration \item le voyage d'agrément
\end{enumerate}
\item Quel est le pluriel de \all{die Grenze} ?
\begin{enumerate}
\item die Grenzen \item die Grenzer \item die Grenz
\end{enumerate}
\item \all{das Visum} est :
\begin{enumerate}
\item un document officiel autorisant l'entrée dans un pays \item un billet d'avion \item une carte d'identité
\end{enumerate}
\item Le contraire de \all{die Arbeitslosigkeit} est :
\begin{enumerate}
\item die Flucht \item die Beschäftigung \item das Ausland
\end{enumerate}
\item \all{jemandem die Flucht ermöglichen} veut dire :
\begin{enumerate}
\item empêcher quelqu'un de partir \item permettre à quelqu'un de fuir \item expulser quelqu'un
\end{enumerate}
\end{enumerate}
\end{exercice}

\begin{exercice}[Vrai ou faux ? Justifie ta réponse]
Réponds par \all{richtig} ou \all{falsch} et justifie chaque réponse par une phrase complète en allemand.
\begin{enumerate}
\item \all{weil} est une conjonction de coordination : le verbe conjugué reste à la deuxième place.
\item Après \all{denn}, le verbe conjugué se trouve à la deuxième place de la proposition.
\item \all{um ... zu} s'emploie quand le sujet des deux propositions est le même.
\item \all{damit} exprime la cause.
\item Le subjonctif II de \all{haben} à la 3\ieme{} personne du singulier est \all{hätte}.
\end{enumerate}
\end{exercice}

\begin{exercice}[Familles de mots]
Complète le tableau avec le mot de la même famille qui manque (nom, verbe ou adjectif).

\begin{center}
\begin{tabularx}{\linewidth}{|X|X|X|}
\hline
\rowcolor{popblueL} Nom & Verbe & Adjectif \\
\hline
die Auswanderung & \ldots & \ldots \\
\hline
\ldots & fliehen & \ldots \\
\hline
die Arbeitslosigkeit & \ldots & arbeitslos \\
\hline
\ldots & \ldots & fremd \\
\hline
die Hoffnung & \ldots & \ldots \\
\hline
\end{tabularx}
\end{center}
\end{exercice}

\begin{exercice}[Conjugaison : subjonctif II]
Mets les verbes entre parenthèses au subjonctif II.
\begin{enumerate}
\item Wenn ich mehr Geld (haben), (können) ich im Ausland studieren.
\item Wenn die Arbeitslosigkeit nicht so hoch (sein), (bleiben) viele junge Leute im Land.
\item Viele Jugendliche (wandern) nicht aus, wenn sie hier eine Arbeit (finden).
\item Wenn ich ein Visum (bekommen), (fliegen) ich sofort nach Deutschland.
\item (wissen) du, was ich an deiner Stelle (tun) ?
\end{enumerate}
\end{exercice}

\courssub{Je m'entraîne}

\begin{exercice}[Relie les phrases : weil, denn ou da]
Relie les deux phrases avec la conjonction indiquée entre parenthèses. Attention à la place du verbe !
\begin{enumerate}
\item Viele junge Kameruner verlassen ihr Land. Sie finden keine Arbeit. (\all{weil})
\item Er lernt Deutsch. Er möchte in Deutschland studieren. (\all{denn})
\item Sie beantragt ein Visum. Sie will ihre Familie in Europa besuchen. (\all{da})
\item Die Grenzen sind streng kontrolliert. Viele Flüchtlinge warten monatelang. (\all{weil})
\item Die Jugendlichen hoffen auf ein besseres Leben. Die Lage in ihrer Heimat ist schwierig. (\all{denn})
\end{enumerate}
\end{exercice}

\begin{exercice}[Transforme : damit ou um ... zu ?]
Remplace \all{damit} par \all{um ... zu} quand c'est possible. Si ce n'est pas possible, explique pourquoi.
\begin{enumerate}
\item Viele Eltern sparen Geld, damit ihre Kinder im Ausland studieren können.
\item Er arbeitet viel, damit er seiner Familie Geld schicken kann.
\item Sie lernt Deutsch, damit sie ein Visum bekommt.
\item Die Regierung schafft neue Jobs, damit die Jugendlichen im Land bleiben.
\item Ich schreibe einen Brief, damit ich Informationen über die Botschaft bekomme.
\end{enumerate}
\end{exercice}

\begin{exercice}[Texte à trous : le départ de Nadia]
Complète le texte avec les mots suivants : \all{Visum -- Grenze -- Auswanderung -- Arbeitslosigkeit -- Ausland -- Flucht}.

\begin{textebox}[Der Abschied]
Nadia träumt schon lange von der \ldots\ldots\ldots\ldots{} . Wegen der hohen \ldots\ldots\ldots\ldots{} in ihrer Stadt hat sie nach dem Abitur keine Arbeit gefunden. Sie möchte ins \ldots\ldots\ldots\ldots{} gehen und in Deutschland Krankenschwester werden. Zuerst muss sie bei der Botschaft ein \ldots\ldots\ldots\ldots{} beantragen. An der \ldots\ldots\ldots\ldots{} werden ihre Papiere kontrolliert. Für ihre Mutter ist der Abschied schwer, aber es ist keine \ldots\ldots\ldots\ldots{} : Nadia will eines Tages zurückkommen.
\end{textebox}
\end{exercice}

\begin{exercice}[Appariements : cause et but]
Relie chaque début de phrase (1--5) à la fin qui convient (a--e), puis traduis les phrases complètes en français.
\begin{enumerate}
\item Weil er keine Stelle findet, \ldots
\item Er lernt Deutsch, um \ldots
\item Da die Lage schwierig ist, \ldots
\item Er schickt Geld nach Hause, damit \ldots
\item Denn die Arbeitslosigkeit ist hoch, \ldots
\end{enumerate}
\begin{enumerate}
\item[a.] in Deutschland studieren zu können.
\item[b.] verlassen viele Jugendliche das Land.
\item[c.] seine Familie ein besseres Leben hat.
\item[d.] bewirbt er sich im Ausland.
\item[e.] und die Perspektiven fehlen.
\end{enumerate}
\end{exercice}

\courssub{Je me prépare au Bac}

\begin{exercice}[Compréhension de texte type Bac]
Lis le texte suivant, puis réponds aux questions en allemand par des phrases complètes.

\begin{textebox}[„Ich komme zurück, aber später“ -- Ein Gespräch mit einem jungen Auswanderer]
Jean-Claude, 24 Jahre alt, sitzt in einem Café in Yaoundé und erzählt von seinen Plänen. Nach dem Studium hat er zwei Jahre lang eine Stelle gesucht, aber er hat nichts gefunden. „Weil ich hier keine Perspektive habe, möchte ich nach Deutschland auswandern“, sagt er. „Dort gibt es Arbeit, und die Löhne sind höher. Ich habe schon Deutsch gelernt, damit ich mich integrieren kann. Mein Onkel, der seit zehn Jahren in München lebt, hilft mir bei den Papieren.“

Seine Mutter ist traurig, denn sie verliert ihren einzigen Sohn. Aber sie versteht ihn: „Wenn die jungen Leute hier Arbeit hätten, würden sie bleiben. Das ist das größte Problem unseres Landes.“ Jean-Claude verspricht: „Ich schicke Geld nach Hause, damit meine kleine Schwester studieren kann. Und eines Tages komme ich zurück, um hier ein Geschäft zu eröffnen.“

Experten warnen: Wenn zu viele gut ausgebildete junge Leute das Land verlassen, fehlen später Ärzte, Ingenieure und Lehrer. Die Auswanderung ist also nicht nur eine Chance, sondern auch ein Verlust für die Heimat.
\end{textebox}

\textbf{Questions de compréhension}
\begin{enumerate}
\item Warum möchte Jean-Claude auswandern ? (2 raisons)
\item Wer hilft ihm bei den Papieren ?
\item Was verspricht er seiner Familie ?
\item Wovor warnen die Experten ?
\end{enumerate}

\textbf{Vrai ou faux ? Justifie par une citation du texte.}
\begin{enumerate}
\item Jean-Claude hat nach dem Studium sofort eine Stelle gefunden.
\item Seine Mutter ist froh über seine Entscheidung.
\item Er will für immer im Ausland bleiben.
\end{enumerate}

\textbf{Questions de langue}
\begin{enumerate}
\item Trouve dans le texte deux connecteurs exprimant la cause et deux exprimant le but.
\item Mets la phrase suivante au subjonctif II : \all{Wenn die jungen Leute hier Arbeit haben, bleiben sie.}
\item Relève dans le texte deux verbes au Perfekt et donne leur infinitif.
\end{enumerate}
\end{exercice}

\begin{exercice}[Thème et version]
\textbf{A. Thème.} Traduis en allemand les phrases suivantes.
\begin{enumerate}
\item Beaucoup de jeunes Camerounais émigrent parce qu'ils ne trouvent pas de travail.
\item Il apprend l'allemand pour pouvoir étudier en Allemagne.
\item Elle envoie de l'argent à sa famille afin que ses enfants aillent à l'école.
\item Si j'avais un visa, je voyagerais en Europe.
\item Comme le chômage est très élevé, beaucoup de jeunes perdent espoir.
\end{enumerate}

\textbf{B. Version.} Traduis en français le texte suivant.

\begin{textebox}[Ein Brief aus Deutschland]
Liebe Mama, seit drei Monaten lebe ich jetzt in Hamburg. Ich habe eine Arbeit in einem Krankenhaus gefunden, und meine Kollegen sind sehr freundlich. Weil das Leben hier teuer ist, muss ich sparsam sein, aber ich schicke euch jeden Monat Geld, damit ihr die Schule für die Kleinen bezahlen könnt. Manchmal bin ich traurig, denn ich vermisse euch sehr. Wenn ich mehr verdiente, würde ich euch sofort besuchen kommen. Um mein Deutsch zu verbessern, besuche ich jeden Abend einen Sprachkurs. In zwei Jahren möchte ich zurückkommen, denn Kamerun bleibt meine Heimat. Dein Sohn Etienne.
\end{textebox}
\end{exercice}

\begin{exercice}[Production écrite type Bac]
Rédige en allemand un texte de 10 à 12 lignes sur le sujet suivant :

\begin{center}
\textbf{\all{Die Auswanderung der jungen Kameruner: Gründe und Folgen}}
\end{center}

\textbf{Consignes :}
\begin{enumerate}
\item Présente d'abord les raisons de l'émigration (au moins deux), puis ses conséquences pour le pays et pour les familles.
\item Termine par ton avis personnel : faut-il partir ou rester ?
\item Emploie obligatoirement les connecteurs suivants : \all{weil}, \all{denn}, \all{um ... zu}, \all{damit}, \all{deshalb}, et au moins une phrase au subjonctif II commençant par \all{Wenn ...}.
\end{enumerate}
\end{exercice}

\newpage

\coursec{Corrigés des exercices}

\coursec{Corrigés des exercices — Leçon AL28 : Les problèmes d'émigration}

\begin{corrige}[Exercice 1 — QCM : vocabulaire de l'émigration]
\begin{enumerate}
\item Réponse b : \all{die Auswanderung} signifie \emph{l'émigration}. Attention au faux sens : \all{die Einwanderung} = l'immigration (préfixe \all{ein-} = « entrer », \all{aus-} = « sortir »).
\item Réponse a : le pluriel de \all{die Grenze} est \all{die Grenzen} (pluriel en \all{-n}, comme la plupart des noms féminins en \all{-e}).
\item Réponse a : \all{das Visum} (pluriel : \all{die Visa}) est un document officiel autorisant l'entrée dans un pays. Un billet d'avion se dit \all{das Flugticket}, une carte d'identité \all{der Personalausweis}.
\item Réponse b : le contraire de \all{die Arbeitslosigkeit} (le chômage) est \all{die Beschäftigung} (l'emploi, l'occupation). \all{die Flucht} = la fuite ; \all{das Ausland} = l'étranger.
\item Réponse b : \all{jemandem die Flucht ermöglichen} = permettre à quelqu'un de fuir. « Empêcher » se dirait \all{verhindern}, « expulser » \all{abschieben}.
\end{enumerate}
\end{corrige}

\begin{corrige}[Exercice 2 — Vrai ou faux ?]
\begin{enumerate}
\item \all{falsch}. \all{Weil} est une conjonction de \emph{subordination} : le verbe conjugué va à la fin de la subordonnée. \all{Viele junge Leute wandern aus, weil sie keine Arbeit finden.} (« Beaucoup de jeunes émigrent parce qu'ils ne trouvent pas de travail. »)
\item \all{richtig}. \all{Denn} est une conjonction de \emph{coordination} : elle n'occupe pas de place dans la phrase, le verbe conjugué reste en deuxième position. \all{Er bleibt nicht, denn er findet keine Stelle.} (« Il ne reste pas, car il ne trouve pas de poste. »)
\item \all{richtig}. \all{Um ... zu} + infinitif ne s'emploie que si le sujet des deux propositions est identique. \all{Sie lernt Deutsch, um ein Visum zu bekommen.} (« Elle apprend l'allemand pour obtenir un visa. » — c'est elle qui apprend et elle qui obtient.)
\item \all{falsch}. \all{Damit} exprime le \emph{but}, non la cause. \all{Er spart Geld, damit seine Kinder studieren können.} (« Il économise de l'argent afin que ses enfants puissent étudier. »)
\item \all{richtig}. Le subjonctif II de \all{haben} se construit sur le prétérit \all{hatte} avec tréma : \all{ich hätte, du hättest, er/sie/es hätte}. \all{Wenn sie Arbeit hätte, bliebe sie hier.} (« Si elle avait du travail, elle resterait ici. »)
\end{enumerate}
\end{corrige}

\begin{corrige}[Exercice 3 — Familles de mots]
\begin{center}
\begin{tabularx}{\linewidth}{|X|X|X|}
\hline
\rowcolor{popblueL} Nom & Verbe & Adjectif \\
\hline
die Auswanderung & auswandern & ausgewandert \\
\hline
die Flucht & fliehen & flüchtig \\
\hline
die Arbeitslosigkeit & arbeitslos sein & arbeitslos \\
\hline
die Fremde & --- (fremd sein) & fremd \\
\hline
die Hoffnung & hoffen & hoffnungsvoll / hoffnungslos \\
\hline
\end{tabularx}
\end{center}
\textbf{Remarques :}
\begin{itemize}
\item \all{fliehen} est un verbe fort irrégulier : \all{fliehen -- floh -- ist geflohen} (auxiliaire \all{sein}).
\item Pour \all{die Fremde} (le pays étranger, l'étranger), il n'existe pas de verbe simple ; on accepte \all{fremd sein} ou l'expression \all{in der Fremde leben} (« vivre à l'étranger »).
\item \all{hoffnungsvoll} = plein d'espoir ; \all{hoffnungslos} = désespéré. Le suffixe \all{-los} exprime la privation : \all{arbeitslos}, \all{hoffnungslos}, \all{perspektivlos}.
\end{itemize}
\end{corrige}

\begin{corrige}[Exercice 4 — Conjugaison : subjonctif II]
\begin{enumerate}
\item \all{Wenn ich mehr Geld \textbf{hätte}, \textbf{könnte} ich im Ausland studieren.} (« Si j'avais plus d'argent, je pourrais étudier à l'étranger. ») — \all{haben} $\rightarrow$ \all{hätte} ; \all{können} $\rightarrow$ \all{könnte} (prétérit \all{konnte} + tréma).
\item \all{Wenn die Arbeitslosigkeit nicht so hoch \textbf{wäre}, \textbf{blieben} viele junge Leute im Land.} (« Si le chômage n'était pas si élevé, beaucoup de jeunes resteraient dans le pays. ») — \all{sein} $\rightarrow$ \all{wäre} (irrégulier) ; \all{bleiben} $\rightarrow$ \all{blieben} (prétérit \all{blieb} + tréma + \all{-en}). On accepte aussi \all{würden bleiben}.
\item \all{Viele Jugendliche \textbf{würden} nicht \textbf{auswandern}, wenn sie hier eine Arbeit \textbf{fänden}.} (« Beaucoup de jeunes n'émigreraient pas s'ils trouvaient un travail ici. ») — verbe faible : on préfère \all{würden auswandern} ; \all{finden} $\rightarrow$ \all{fänden} (prétérit \all{fand} + tréma).
\item \all{Wenn ich ein Visum \textbf{bekäme}, \textbf{würde} ich sofort nach Deutschland \textbf{fliegen}.} (« Si j'obtenais un visa, je prendrais immédiatement l'avion pour l'Allemagne. ») — \all{bekommen} $\rightarrow$ \all{bekäme} ; \all{flöge} est correct mais vieilli, on préfère \all{würde fliegen}.
\item \all{\textbf{Wüsstest} du, was ich an deiner Stelle \textbf{täte}?} (« Saurais-tu ce que je ferais à ta place ? ») — \all{wissen} $\rightarrow$ \all{wüsstest} ; \all{tun} $\rightarrow$ \all{täte}. On accepte aussi \all{würde ich tun}.
\end{enumerate}
\textbf{Point de vigilance :} le subjonctif II se forme sur le radical du \emph{prétérit} ; les voyelles radicales \all{a, o, u} prennent le tréma (\all{hatte $\rightarrow$ hätte}, \all{kam $\rightarrow$ käme}, \all{wusste $\rightarrow$ wüsste}). \all{Sein} est totalement irrégulier : \all{wäre}.
\end{corrige}

\begin{corrige}[Exercice 5 — Relie les phrases : weil, denn ou da]
\begin{enumerate}
\item \all{Viele junge Kameruner verlassen ihr Land, weil sie keine Arbeit finden.} (« Beaucoup de jeunes Camerounais quittent leur pays parce qu'ils ne trouvent pas de travail. ») — subordonnée introduite par \all{weil} : verbe conjugué \all{finden} à la fin.
\item \all{Er lernt Deutsch, denn er möchte in Deutschland studieren.} (« Il apprend l'allemand, car il voudrait étudier en Allemagne. ») — \all{denn} est une coordination : verbe \all{möchte} en deuxième place.
\item \all{Da sie ihre Familie in Europa besuchen will, beantragt sie ein Visum.} (« Comme elle veut rendre visite à sa famille en Europe, elle demande un visa. ») — \all{da} se place de préférence en tête de phrase ; la principale qui suit commence par le verbe conjugué (\all{beantragt}) : inversion sujet--verbe.
\item \all{Weil die Grenzen streng kontrolliert werden, warten viele Flüchtlinge monatelang.} (« Parce que les frontières sont strictement contrôlées, beaucoup de réfugiés attendent des mois. ») — subordonnée en tête : le verbe de la principale (\all{warten}) occupe la première place de la principale ; attention au passif \all{kontrolliert werden}.
\item \all{Die Jugendlichen hoffen auf ein besseres Leben, denn die Lage in ihrer Heimat ist schwierig.} (« Les jeunes espèrent une vie meilleure, car la situation dans leur patrie est difficile. ») — après \all{denn}, ordre normal : sujet \all{die Lage}, verbe \all{ist} en deuxième position.
\end{enumerate}
\textbf{Point de vigilance :} ne jamais écrire \all{weil sie finden keine Arbeit} : après \all{weil}, le verbe conjugué va obligatoirement en dernière position.
\end{corrige}

\begin{corrige}[Exercice 6 — Transforme : damit ou um ... zu ?]
\begin{enumerate}
\item \all{Viele Eltern sparen Geld, damit ihre Kinder im Ausland studieren können.} $\rightarrow$ \textbf{Transformation impossible} : les sujets sont différents (\all{die Eltern} / \all{ihre Kinder}). On conserve \all{damit}.
\item \all{Er arbeitet viel, um seiner Familie Geld schicken zu können.} (« Il travaille beaucoup pour pouvoir envoyer de l'argent à sa famille. ») — même sujet (\all{er}) : transformation possible. L'infinitif \all{schicken} précède \all{zu können}.
\item \all{Sie lernt Deutsch, um ein Visum zu bekommen.} (« Elle apprend l'allemand pour obtenir un visa. ») — même sujet (\all{sie}) : transformation possible.
\item \all{Die Regierung schafft neue Jobs, damit die Jugendlichen im Land bleiben.} $\rightarrow$ \textbf{Transformation impossible} : sujets différents (\all{die Regierung} / \all{die Jugendlichen}). On conserve \all{damit}.
\item \all{Ich schreibe einen Brief, um Informationen über die Botschaft zu bekommen.} (« J'écris une lettre pour obtenir des informations sur l'ambassade. ») — même sujet (\all{ich}) : transformation possible.
\end{enumerate}
\textbf{Règle à retenir :} \all{um ... zu} = même sujet dans les deux propositions ; \all{damit} = sujets différents (ou possibles dans tous les cas). En cas de doute, \all{damit} est toujours correct.
\end{corrige}

\begin{corrige}[Exercice 7 — Texte à trous : le départ de Nadia]
\begin{textebox}[Der Abschied — corrigé]
Nadia träumt schon lange von der \textbf{Auswanderung}. Wegen der hohen \textbf{Arbeitslosigkeit} in ihrer Stadt hat sie nach dem Abitur keine Arbeit gefunden. Sie möchte ins \textbf{Ausland} gehen und in Deutschland Krankenschwester werden. Zuerst muss sie bei der Botschaft ein \textbf{Visum} beantragen. An der \textbf{Grenze} werden ihre Papiere kontrolliert. Für ihre Mutter ist der Abschied schwer, aber es ist keine \textbf{Flucht}: Nadia will eines Tages zurückkommen.
\end{textebox}
\textbf{Traduction :} « Nadia rêve depuis longtemps de l'émigration. À cause du chômage élevé dans sa ville, elle n'a pas trouvé de travail après le baccalauréat. Elle voudrait aller à l'étranger et devenir infirmière en Allemagne. D'abord, elle doit demander un visa à l'ambassade. À la frontière, ses papiers seront contrôlés. Pour sa mère, l'adieu est dur, mais ce n'est pas une fuite : Nadia veut revenir un jour. »

\textbf{Justifications :}
\begin{itemize}
\item \all{von der Auswanderung} : datif féminin après \all{von} (\all{die Auswanderung} $\rightarrow$ \all{der}).
\item \all{wegen der hohen Arbeitslosigkeit} : \all{wegen} + génitif (datif admis à l'oral).
\item \all{ins Ausland} : contraction de \all{in das Ausland} ; \all{ins Ausland gehen} = partir à l'étranger (expression figée).
\item \all{ein Visum beantragen} : accusatif neutre ; \all{beantragen} = demander officiellement (verbe à particule \emph{inséparable} : \all{sie beantragt}).
\item \all{an der Grenze} : datif féminin (lieu, question « où ? »).
\item \all{keine Flucht} : la négation du nom avec article indéfini exige \all{kein-} (accusatif féminin \all{keine}).
\end{itemize}
\end{corrige}

\begin{corrige}[Exercice 8 — Appariements : cause et but]
\textbf{Appariements :} 1--d ; 2--a ; 3--b ; 4--c ; 5--e.

\textbf{Phrases complètes et traductions :}
\begin{enumerate}
\item \all{Weil er keine Stelle findet, bewirbt er sich im Ausland.} « Parce qu'il ne trouve pas de poste, il pose sa candidature à l'étranger. » — \all{sich bewerben} : verbe fort réfléchi (\all{bewirbt sich -- bewarb sich -- hat sich beworben}) ; subordonnée en tête, donc inversion dans la principale.
\item \all{Er lernt Deutsch, um in Deutschland studieren zu können.} « Il apprend l'allemand pour pouvoir étudier en Allemagne. » — même sujet : \all{um ... zu} + infinitif en fin de phrase.
\item \all{Da die Lage schwierig ist, verlassen viele Jugendliche das Land.} « Comme la situation est difficile, beaucoup de jeunes quittent le pays. » — \all{da} en tête ; verbe de la principale (\all{verlassen}) en première position.
\item \all{Er schickt Geld nach Hause, damit seine Familie ein besseres Leben hat.} « Il envoie de l'argent à la maison afin que sa famille ait une vie meilleure. » — sujets différents (\all{er} / \all{seine Familie}) : \all{damit} obligatoire.
\item \all{Denn die Arbeitslosigkeit ist hoch, und die Perspektiven fehlen.} « Car le chômage est élevé et les perspectives manquent. » — coordination : ordre des mots normal dans les deux propositions.
\end{enumerate}
\end{corrige}

\begin{corrige}[Exercice 9 — Compréhension de texte type Bac]
\textbf{Barème indicatif (sur 20) :} questions de compréhension : 8 pts (2 pts par question) ; vrai/faux justifié : 6 pts (2 pts par item : 1 pt pour la réponse, 1 pt pour la citation) ; questions de langue : 6 pts (2 pts par question).

\textbf{Questions de compréhension — réponses attendues :}
\begin{enumerate}
\item \all{Jean-Claude möchte auswandern, weil er in Kamerun keine Perspektive hat und weil es in Deutschland Arbeit und höhere Löhne gibt.} (« Jean-Claude veut émigrer parce qu'il n'a pas de perspective au Cameroun et parce qu'il y a du travail et des salaires plus élevés en Allemagne. »)
\item \all{Sein Onkel, der seit zehn Jahren in München lebt, hilft ihm bei den Papieren.} (« Son oncle, qui vit depuis dix ans à Munich, l'aide pour les papiers. »)
\item \all{Er verspricht, Geld nach Hause zu schicken, damit seine kleine Schwester studieren kann, und eines Tages zurückzukommen, um ein Geschäft zu eröffnen.} (« Il promet d'envoyer de l'argent à la maison afin que sa petite sœur puisse étudier, et de revenir un jour pour ouvrir un commerce. »)
\item \all{Die Experten warnen davor, dass später Ärzte, Ingenieure und Lehrer fehlen, wenn zu viele gut ausgebildete junge Leute das Land verlassen.} (« Les experts mettent en garde : si trop de jeunes bien formés quittent le pays, il manquera plus tard des médecins, des ingénieurs et des enseignants. »)
\end{enumerate}

\textbf{Vrai ou faux :}
\begin{enumerate}
\item \all{falsch} — citation : \all{Nach dem Studium hat er zwei Jahre lang eine Stelle gesucht, aber er hat nichts gefunden.} (Il a cherché un poste pendant deux ans, sans rien trouver.)
\item \all{falsch} — citation : \all{Seine Mutter ist traurig, denn sie verliert ihren einzigen Sohn.} (Sa mère est triste, non pas joyeuse.)
\item \all{falsch} — citation : \all{Und eines Tages komme ich zurück, um hier ein Geschäft zu eröffnen.} (Il compte revenir un jour.)
\end{enumerate}

\textbf{Questions de langue :}
\begin{enumerate}
\item Cause : \all{weil} (\all{Weil ich hier keine Perspektive habe, ...}) et \all{denn} (\all{..., denn sie verliert ihren einzigen Sohn}). But : \all{damit} (\all{..., damit meine kleine Schwester studieren kann}) et \all{um ... zu} (\all{..., um hier ein Geschäft zu eröffnen}).
\item \all{Wenn die jungen Leute hier Arbeit \textbf{hätten}, \textbf{würden} sie \textbf{bleiben}.} (ou \all{blieben sie}). C'est d'ailleurs la phrase du texte, déjà au subjonctif II : \all{hätten} (prétérit \all{hatten} + tréma) et \all{würden bleiben}.
\item Exemples de verbes au Perfekt : \all{hat ... gesucht} (infinitif \all{suchen}), \all{hat ... gefunden} (\all{finden}), \all{hat ... gelernt} (\all{lernen}). Deux d'entre eux suffisent.
\end{enumerate}

\textbf{Points de vigilance :} répondre par des phrases complètes avec verbe conjugué en deuxième position ; dans les citations, recopier exactement le texte avec les majuscules ; ne pas confondre \all{Perspektive} (perspective d'avenir) et le faux ami français « perspective » au sens pictural.
\end{corrige}

\begin{corrige}[Exercice 10 — Thème et version]
\textbf{A. Thème — propositions de traduction :}
\begin{enumerate}
\item \all{Viele junge Kameruner wandern aus, weil sie keine Arbeit finden.} — « parce que » = \all{weil} : verbe \all{finden} rejeté à la fin ; « émigrer » = \all{auswandern} (particule séparable : \all{wandern ... aus}).
\item \all{Er lernt Deutsch, um in Deutschland studieren zu können.} — même sujet : \all{um ... zu} ; « pouvoir étudier » = \all{studieren zu können}.
\item \all{Sie schickt ihrer Familie Geld, damit ihre Kinder zur Schule gehen.} — sujets différents : \all{damit} ; « à sa famille » = datif \all{ihrer Familie} ; « aller à l'école » = \all{zur Schule gehen}.
\item \all{Wenn ich ein Visum hätte, würde ich nach Europa reisen.} — condition irréelle : subjonctif II \all{hätte} / \all{würde reisen} ; « voyager » = \all{reisen} (attention : \all{fahren} = aller en véhicule).
\item \all{Da die Arbeitslosigkeit sehr hoch ist, verlieren viele junge Leute die Hoffnung.} — « comme » en tête = \all{da} ; la principale commence par le verbe \all{verlieren} ; « perdre espoir » = \all{die Hoffnung verlieren}.
\end{enumerate}

\textbf{B. Version — traduction proposée :}

« Chère maman, voilà trois mois que je vis à Hambourg. J'ai trouvé un travail dans un hôpital, et mes collègues sont très aimables. Comme la vie est chère ici, je dois être économe, mais je vous envoie de l'argent chaque mois afin que vous puissiez payer l'école des petits. Parfois je suis triste, car vous me manquez beaucoup. Si je gagnais davantage, je viendrais vous rendre visite immédiatement. Pour améliorer mon allemand, je suis un cours de langue tous les soirs. Dans deux ans, je voudrais revenir, car le Cameroun reste ma patrie. Ton fils Etienne. »

\textbf{Points de vigilance :}
\begin{itemize}
\item \all{seit drei Monaten lebe ich} : présent allemand + \all{seit} = français « voilà trois mois que » (ne pas traduire par un passé composé).
\item \all{ich vermisse euch} : « vous me manquez » — le sujet change entre les deux langues ; ne jamais traduire \all{vermissen} par « manquer » avec le même sujet.
\item \all{sparsam sein} = être économe ; \all{das Krankenhaus} = l'hôpital.
\item \all{Wenn ich mehr verdiente, würde ich euch sofort besuchen kommen} : subjonctif II = conditionnel français (« si je gagnais..., je viendrais... »).
\end{itemize}
\textbf{Barème indicatif :} thème 10 pts (2 pts par phrase : lexique 1 pt, grammaire et ordre des mots 1 pt) ; version 10 pts (fidélité 5 pts, qualité du français 3 pts, lexique 2 pts).
\end{corrige}

\begin{corrige}[Exercice 11 — Production écrite type Bac]
\textbf{Barème indicatif (sur 20) :} respect du sujet et de la structure (raisons, conséquences, avis personnel) : 5 pts ; emploi correct des cinq connecteurs imposés : 5 pts ; subjonctif II correct : 2 pts ; richesse du vocabulaire thématique : 4 pts ; correction grammaticale (ordre des mots, cas, conjugaisons) : 4 pts.

\textbf{Proposition de rédaction modèle :}

\begin{textebox}[Die Auswanderung der jungen Kameruner: Gründe und Folgen]
In Kamerun verlassen jedes Jahr viele junge Leute ihr Land. Warum wandern sie aus? Zuerst ist die Arbeitslosigkeit sehr hoch: Viele Jugendliche finden nach dem Studium keine Stelle, weil es zu wenige Jobs gibt. Außerdem sind die Löhne niedrig, denn die Wirtschaft ist schwach. Deshalb träumen viele von einem besseren Leben im Ausland. Um ein Visum zu bekommen, lernen sie Deutsch oder Englisch und sparen lange Geld.

Die Auswanderung hat aber auch Folgen. Die Familien sind traurig, denn sie verlieren ihre Kinder. Viele Eltern bekommen zwar Geld aus Europa, damit ihre Kinder die Schule besuchen können, aber das Geld ersetzt nicht die Familie. Für das Land ist die Auswanderung ein großer Verlust: Wenn die gut ausgebildeten jungen Leute blieben, hätte Kamerun mehr Ärzte, Ingenieure und Lehrer.

Meiner Meinung nach sollten die Jugendlichen nicht alle weggehen. Wenn ich die Wahl hätte, würde ich zuerst hier eine Chance suchen. Kamerun braucht seine Jugend, um sich zu entwickeln.
\end{textebox}

\textbf{Traduction du modèle :} « Au Cameroun, beaucoup de jeunes quittent leur pays chaque année. Pourquoi émigrent-ils ? D'abord, le chômage est très élevé : beaucoup de jeunes ne trouvent pas de poste après leurs études parce qu'il y a trop peu d'emplois. De plus, les salaires sont bas, car l'économie est faible. C'est pourquoi beaucoup rêvent d'une vie meilleure à l'étranger. Pour obtenir un visa, ils apprennent l'allemand ou l'anglais et économisent longtemps de l'argent. Mais l'émigration a aussi des conséquences. Les familles sont tristes, car elles perdent leurs enfants. Beaucoup de parents reçoivent certes de l'argent d'Europe afin que leurs enfants puissent aller à l'école, mais l'argent ne remplace pas la famille. Pour le pays, l'émigration est une grande perte : si les jeunes bien formés restaient, le Cameroun aurait plus de médecins, d'ingénieurs et d'enseignants. À mon avis, les jeunes ne devraient pas tous partir. Si j'avais le choix, je chercherais d'abord une chance ici. Le Cameroun a besoin de sa jeunesse pour se développer. »

\textbf{Points de vigilance :}
\begin{itemize}
\item Vérifier que les cinq connecteurs imposés apparaissent : \all{weil} (verbe final), \all{denn} (verbe en 2\up{e} position), \all{um ... zu} (même sujet), \all{damit} (sujets différents), \all{deshalb} (suivi de l'inversion : \all{deshalb träumen viele}).
\item La phrase au subjonctif II : \all{Wenn die gut ausgebildeten jungen Leute blieben, hätte Kamerun mehr Ärzte} — verbe de la subordonnée à la fin, verbe de la principale en tête.
\item Majuscules à tous les noms : \all{die Auswanderung, die Arbeitslosigkeit, die Löhne, das Visum, die Folgen}.
\item Éviter les calques du français : « chercher du travail » = \all{eine Stelle suchen} (et non \all{Arbeit suchen} est toléré mais \all{eine Arbeit finden} est plus idiomatique) ; « à mon avis » = \all{meiner Meinung nach}.
\item Compter les lignes : 10 à 12 lignes exigées ; une introduction, deux ou trois paragraphes, une conclusion personnelle.
\end{itemize}
\end{corrige}

\newpage

\coursec{Fiche bilan}

\begin{popfigure}
\begin{center}\tcbox[colback=popblue,colframe=popblue,coltext=white,arc=9pt,boxsep=4pt,left=6pt,right=6pt]{\bfseries\small Les probl\`emes d'\'emigration (AL28)}\end{center}
\vspace{-2pt}
\begin{tcbraster}[raster columns=3,raster equal height=rows,raster column skip=6pt,raster row skip=6pt,enhanced,blankest]
\begin{tcolorbox}[enhanced,colback=poporange,colframe=poporange,coltext=white,boxrule=0.9pt,arc=3pt,left=4pt,right=4pt,top=3pt,bottom=3pt,boxsep=1pt,halign=left]\bfseries\scriptsize Texte : \all{Warum so viele junge Leute auswandern}\end{tcolorbox}
\begin{tcolorbox}[enhanced,colback=popgreen,colframe=popgreen,coltext=white,boxrule=0.9pt,arc=3pt,left=4pt,right=4pt,top=3pt,bottom=3pt,boxsep=1pt,halign=left]\bfseries\scriptsize Lexique : \all{die Auswanderung}, \all{die Flucht}, \all{das Visum}, \all{die Grenze}\end{tcolorbox}
\begin{tcolorbox}[enhanced,colback=poppurple,colframe=poppurple,coltext=white,boxrule=0.9pt,arc=3pt,left=4pt,right=4pt,top=3pt,bottom=3pt,boxsep=1pt,halign=left]\bfseries\scriptsize Cause : \all{weil}, \all{denn}, \all{da}\end{tcolorbox}
\begin{tcolorbox}[enhanced,colback=poppink,colframe=poppink,coltext=white,boxrule=0.9pt,arc=3pt,left=4pt,right=4pt,top=3pt,bottom=3pt,boxsep=1pt,halign=left]\bfseries\scriptsize But : \all{damit}, \all{um \dots{} zu}\end{tcolorbox}
\begin{tcolorbox}[enhanced,colback=popgold,colframe=popgold,coltext=white,boxrule=0.9pt,arc=3pt,left=4pt,right=4pt,top=3pt,bottom=3pt,boxsep=1pt,halign=left]\bfseries\scriptsize Temps : Perfekt / Pr\"ateritum\end{tcolorbox}
\begin{tcolorbox}[enhanced,colback=popblue!60,colframe=popblue!60,coltext=white,boxrule=0.9pt,arc=3pt,left=4pt,right=4pt,top=3pt,bottom=3pt,boxsep=1pt,halign=left]\bfseries\scriptsize Konjunktiv II : \all{w\"urde}, \all{h\"atte}, \all{w\"are}\end{tcolorbox}
\end{tcbraster}

\legende{Carte mentale de la le\c{c}on AL28 : les probl\`emes d'\'emigration}
\end{popfigure}

\begin{aretenir}[L'essentiel de la le\c{c}on]
\textbf{1. Lexique cl\'e (article + pluriel \`a conna\^{\i}tre)}

\begin{center}
\begin{tabularx}{\linewidth}{|l|X|}
\hline
\rowcolor{popblueL} \textbf{Deutsch} & \textbf{Fran\c{c}ais} \\
\hline
\all{die Auswanderung, -en} & l'\'emigration (f.) \\
\hline
\all{der Auswanderer, - / die Auswanderin, -nen} & l'\'emigrant / l'\'emigrante \\
\hline
\all{die Flucht, -en} & la fuite \\
\hline
\all{der Fl\"uchtling, -e} & le r\'efugi\'e \\
\hline
\all{das Visum, -s / die Visa} & le visa \\
\hline
\all{die Grenze, -n} & la fronti\`ere \\
\hline
\all{die Arbeitslosigkeit (Sg.)} & le ch\^omage \\
\hline
\all{das Ausland (Sg.)} & l'\'etranger \\
\hline
\all{die Hoffnung, -en} & l'espoir, l'esp\'erance (f.) \\
\hline
\all{die Heimat (Sg.)} & la patrie, le pays natal \\
\hline
\all{auswandern / wanderte aus / ist ausgewandert} & \'emigrer \\
\hline
\all{fliehen / floh / ist geflohen} & fuir \\
\hline
\all{verlassen / verlie\ss{} / hat verlassen} & quitter \\
\hline
\all{sich integrieren / integrierte sich / hat sich integriert} & s'int\'egrer \\
\hline
\end{tabularx}
\end{center}

\textbf{2. Grammaire \`a conna\^{\i}tre par c\oe ur}

\begin{center}
\begin{tabularx}{\linewidth}{|l|X|X|}
\hline
\rowcolor{popgreenL} \textbf{Point} & \textbf{R\`egle} & \textbf{Exemple} \\
\hline
\all{weil} / \all{da} (cause) & subordonn\'ee : verbe conjugu\'e \`a la fin & \all{Viele gehen, weil sie keine Arbeit finden.} \\
\hline
\all{denn} (cause) & coordination : verbe en 2\up{e} position, pas d'inversion & \all{Sie gehen, denn sie haben keine Hoffnung.} \\
\hline
\all{damit} (but) & subordonn\'ee, sujets diff\'erents, verbe \`a la fin & \all{Er lernt Deutsch, damit seine Kinder es leichter haben.} \\
\hline
\all{um \dots{} zu} (but) & m\^eme sujet, infinitif + \all{zu} \`a la fin & \all{Sie lernt Deutsch, um in Deutschland zu studieren.} \\
\hline
Perfekt & oral : \all{haben/sein} + Partizip II ; \all{sein} avec d\'eplacement/changement & \all{Er ist nach Europa ausgewandert.} \\
\hline
Pr\"ateritum & \'ecrit, r\'ecit ; \all{sein/haben} : \all{war}, \all{hatte} & \all{Sie verlie\ss{} ihre Heimat.} \\
\hline
Konjunktiv II & souhait, hypoth\`ese : \all{w\"urde} + infinitif ; \all{h\"atte}, \all{w\"are}, \all{k\"onnte} & \all{Wenn ich ein Visum h\"atte, w\"urde ich bleiben.} \\
\hline
\end{tabularx}
\end{center}

\textbf{3. M\'ethodes express}
\begin{itemize}
  \item \textbf{Compr\'ehension} : lire le titre et la premi\`ere phrase de chaque paragraphe (sens global), puis rep\'erer les mots du lexique (sens d\'etaill\'e).
  \item \textbf{Cause ou but ?} La cause r\'epond \`a \all{warum?} (pourquoi) ; le but r\'epond \`a \all{wozu?} (pour quoi faire).
  \item \textbf{\all{um \dots{} zu} ou \all{damit} ?} M\^eme sujet dans les deux propositions $\rightarrow$ \all{um \dots{} zu} ; sujets diff\'erents $\rightarrow$ \all{damit}.
  \item \textbf{Perfekt ou Pr\"ateritum ?} Conversation, lettre $\rightarrow$ Perfekt ; article, r\'ecit au pass\'e $\rightarrow$ Pr\"ateritum.
  \item \textbf{Production guid\'ee} : 5--6 phrases, une id\'ee par phrase, verbe en 2\up{e} position, connecteurs vari\'es (\all{weil}, \all{damit}, \all{au\ss{}erdem}, \all{deshalb}).
\end{itemize}
\end{aretenir}

\begin{attention}[Pi\`eges fr\'equents des francophones]
\begin{itemize}
  \item Apr\`es \all{weil}, \all{da} et \all{damit}, le verbe conjugu\'e va \textbf{\`a la fin} : \all{Viele junge Leute wandern aus, weil sie zu Hause keine Arbeit finden.} « Beaucoup de jeunes \'emigrent parce qu'ils ne trouvent pas de travail chez eux. » Ne jamais calquer l'ordre fran\c{c}ais.
  \item Apr\`es \all{denn}, \textbf{pas d'inversion} : \all{Er bleibt nicht, denn er hat keine Hoffnung.} « Il ne reste pas, car il n'a pas d'espoir. »
  \item \all{auswandern} et \all{fliehen} sont des verbes de d\'eplacement : auxiliaire \all{sein} au Perfekt : \all{Sie ist nach Kanada ausgewandert.} « Elle a \'emigr\'e au Canada. »
  \item Attention au genre : \all{das Visum}, \all{das Ausland}, mais \all{die Grenze}, \all{die Flucht}. La majuscule est obligatoire pour tous les noms : \all{der Fl\"uchtling}, \all{die Arbeitslosigkeit}.
  \item Faux ami : \all{die Heimat} ne signifie pas « la maison » (\all{das Haus}) mais « la patrie, le pays natal ».
  \item Au Konjunktiv II, ne pas oublier les Umlaute : \all{h\"atte}, \all{w\"are}, \all{k\"onnte}, \all{w\"urde}.
\end{itemize}
\end{attention}

\begin{exemplebox}[Mod\`eles de phrases pour la production guid\'ee]
\all{Viele junge Kameruner verlassen ihr Land, weil die Arbeitslosigkeit sehr hoch ist.} « Beaucoup de jeunes Camerounais quittent leur pays parce que le ch\^omage est tr\`es \'elev\'e. »

\all{Sie wollen ins Ausland gehen, um dort ein besseres Leben zu finden.} « Ils veulent partir \`a l'\'etranger pour y trouver une vie meilleure. »

\all{Manche sparen lange Geld, damit ihre Familie das Visum bezahlen kann.} « Certains \'economisent longtemps de l'argent pour que leur famille puisse payer le visa. »

\all{Wenn ich die M\"oglichkeit h\"atte, w\"urde ich zuerst in meinem Land arbeiten.} « Si j'avais la possibilit\'e, je travaillerais d'abord dans mon pays. »
\end{exemplebox}

\begin{aretenir}[Checklist avant le Bac]
\begin{itemize}
  \item[$\square$] Je connais l'article et le pluriel des 10 noms cl\'es du th\`eme (\all{die Auswanderung}, \all{das Visum}, \all{die Grenze}\dots).
  \item[$\square$] Je sais conjuguer \all{fliehen}, \all{verlassen}, \all{auswandern} au Perfekt et au Pr\"ateritum.
  \item[$\square$] Je mets le verbe conjugu\'e \`a la fin apr\`es \all{weil}, \all{da} et \all{damit}.
  \item[$\square$] Je n'inverse jamais sujet et verbe apr\`es \all{denn}.
  \item[$\square$] Je choisis correctement entre \all{um \dots{} zu} (m\^eme sujet) et \all{damit} (sujets diff\'erents).
  \item[$\square$] Je forme le Perfekt avec \all{sein} pour les verbes de d\'eplacement (\all{ist ausgewandert}, \all{ist geflohen}).
  \item[$\square$] Je connais le Pr\"ateritum de \all{sein} (\all{war}) et de \all{haben} (\all{hatte}).
  \item[$\square$] Je forme le Konjunktiv II : \all{w\"urde} + infinitif, \all{h\"atte}, \all{w\"are}, \all{k\"onnte}.
  \item[$\square$] Je n'oublie pas la majuscule \`a tous les noms et les Umlaute (\all{Fl\"uchtling}, \all{m\"ochte}).
  \item[$\square$] Je sais r\'epondre aux questions de compr\'ehension par des phrases compl\`etes, verbe en 2\up{e} position.
  \item[$\square$] Je sais r\'ediger un court paragraphe (5--6 phrases) avec au moins une cause et un but.
  \item[$\square$] Je relis ma production : place du verbe, cas apr\`es pr\'epositions, orthographe des noms.
\end{itemize}
\end{aretenir}

\findecours

\end{document}
